% Lesson 7 — Reading: Texts about settling community disputes — Anglais 1ère A — Cours signé OnBuch+
% Programme officiel MINESEC. Compiler avec : tectonic ENA07-reading-texts-about-settling-community-disputes.tex
\newcommand{\DOCMATIERE}{Anglais}
\newcommand{\DOCNIVEAU}{1ère A}
\newcommand{\DOCMODULE}{Chapter 1 : Family and Social Life}
\newcommand{\DOCLECON}{ENA07}
\newcommand{\DOCTITRE}{Lesson 7 — Reading: Texts about settling community disputes}
% OnBuch+ — préambule « cours signés » ENRICHI (compilé avec tectonic / XeLaTeX)
% Système visuel « Pop » d'OnBuch+ : fond crème (jamais blanc), bordures encre
% épaisses, ombres « dures » décalées, titres Archivo Black en capitales.
% Version MATHÉMATIQUES : graphes (pgfplots/tikz), boîtes pédagogiques
% (définition, propriété, méthode, attention, le savais-tu, exercice résolu,
% exercice, TP, bilan), page de garde et signature OnBuch+ ancrée en bas.
%
% Macros attendues dans main.tex AVANT \input{preamble} :
%   \DOCMATIERE  \DOCNIVEAU  \DOCMODULE  \DOCLECON (numéro)  \DOCTITRE
\documentclass[11pt]{article}

\usepackage{fontspec}
\usepackage[english,french]{babel} % français = langue principale ; l'anglais via \eng{..} / textebox
\usepackage[a4paper,margin=2cm,top=3.4cm,bottom=2.8cm,headheight=1.4cm,footskip=1.3cm]{geometry}
\usepackage{amsmath,amssymb,amsfonts}
\usepackage{mathtools} % \xmapsto, \coloneqq... souvent utilisés par les modèles
\usepackage{cancel} % simplifications barrées (bilans de forces)
% \abs{z} (module, valeur absolue) et \norm{u} (norme), utilisés par les modèles
\providecommand{\abs}{}\let\abs\relax\DeclarePairedDelimiter{\abs}{\lvert}{\rvert}
\providecommand{\norm}{}\let\norm\relax\DeclarePairedDelimiter{\norm}{\lVert}{\rVert}
\usepackage[table]{xcolor} % [table] : fournit aussi \rowcolors (alternance de lignes)
\usepackage{pagecolor}
\usepackage{tikz}
\usetikzlibrary{arrows.meta,positioning,calc,shapes.geometric,shapes.misc,shapes.arrows,decorations.pathmorphing,decorations.pathreplacing,decorations.markings,patterns,shadows,mindmap,trees,fit,backgrounds,matrix,intersections,angles,quotes}
\usepackage{pgfplots}
\usepackage[version=4]{mhchem} % équations nucléaires \ce{^{235}_{92}U}
\usepackage[european,siunitx]{circuitikz} % schémas électriques
\pgfplotsset{compat=1.18}
\usepgfplotslibrary{fillbetween,groupplots} % aires entre courbes (calcul intégral)
\usepackage{enumitem}
\usepackage{fancyhdr}
% [most] charge toutes les bibliothèques tcolorbox (listings, minted, raster,
% documentation...) et gonfle inutilement la mémoire TeX sur un document long
% (~300 boîtes) : on ne charge que ce qui est réellement utilisé.
\usepackage[skins,breakable]{tcolorbox}
\usepackage{graphicx}
\usepackage{colortbl}
\usepackage{booktabs}
\usepackage{tabularx}
\usepackage{diagbox} % case d'en-tête diagonale des tableaux à double entrée
% Colonnes extensibles centrée / gauche / droite (C, L, R), souvent utilisées par les modèles
\newcolumntype{C}{>{\centering\arraybackslash}X}
\newcolumntype{L}{>{\raggedright\arraybackslash}X}
\newcolumntype{R}{>{\raggedleft\arraybackslash}X}
\usepackage{array}
\usepackage{multicol}
\usepackage{multirow}
\usepackage{siunitx}
% parse-numbers=false : accepte aussi \SI{2,5 \times 10^{-3}}{...}
\sisetup{locale=FR,output-decimal-marker={,},group-minimum-digits=4,parse-numbers=false}
\DeclareSIUnit\ohm{\ensuremath{\symup{\Omega}}} % U+2126 absent de la police
\DeclareSIUnit\division{div} % réglages d'oscilloscope (ms/div, V/div)
\DeclareSIUnit\tr{tr} % tours (vitesse de rotation, ex. \SI{1500}{\tr\per\minute})
\DeclareSIUnit\molar{M} % concentration molaire (ex. \SI{3}{\molar}, \SI{1}{\micro\molar})
\DeclareSIUnit\an{an} % année (le modèle confond parfois avec \year, un compteur TeX) — ex. \SI{5}{\kilo\meter\per\an}
\DeclareSIUnit\FCFA{FCFA} % franc CFA (ex. \SI{500}{\FCFA\per\kilo\gram})
\DeclareSIUnit\degreeBrix{\textdegree Brix} % teneur en sucre d'un sirop/jus (réfractométrie)
\DeclareSIUnit\month{mois} % le modèle utilise parfois \month, un compteur TeX, comme unité de durée
\usepackage{qrcode}
% \degree : commande fréquemment utilisée par les modèles pour un angle ou une
% température en dehors de \SI{}{} (non fournie par les paquets chargés).
\providecommand{\degree}{^{\circ}}
% \qed (fin de démonstration), utilisé par les modèles sans amsthm
\providecommand{\qed}{\hfill$\square$}
% \barre : double barre des valeurs interdites dans un tableau de variations (tkz-tab)
\providecommand{\barre}{\ensuremath{\|}}
% environnement proof (amsthm non chargé) utilisé par les modèles
\ifdefined\proof\else\newenvironment{proof}[1][Démonstration]{\par\noindent\textit{#1.}\ }{\qed\par}\fi
\usepackage{lastpage}
\usepackage{hyperref}
\hypersetup{hidelinks}

% ============================================================
% Palette « Pop » OnBuch+ — mêmes hex que la classe P (plus_ui.dart)
% ============================================================
\definecolor{poporange}{HTML}{F59321}
\definecolor{popdark}{HTML}{E07A0C}
\definecolor{popcream}{HTML}{FFF6E8}
\definecolor{popink}{HTML}{1C1714}
\definecolor{poppurple}{HTML}{7A5AE0}
\definecolor{popgreen}{HTML}{2E9E6A}
\definecolor{poppink}{HTML}{E0457B}
\definecolor{popblue}{HTML}{2F7DE1}
\definecolor{popgold}{HTML}{C98A12}
\definecolor{popmuted}{HTML}{8A7F76}
\definecolor{popline}{HTML}{EADFD2}
\definecolor{popgray}{HTML}{8A7F76}
\colorlet{popgrey}{popgray}
% Alias tolérants (le modèle nomme parfois les couleurs différemment)
\colorlet{popwhite}{white}
\colorlet{popblack}{popink}
\colorlet{popred}{poppink}
\colorlet{popyellow}{popgold}
% Teintes claires (fonds de tableaux/encadrés) — le modèle nomme parfois
% les couleurs avec un suffixe L (ex. popblueL) sans les avoir définies.
\colorlet{poporangeL}{poporange!20}
\colorlet{popdarkL}{popdark!20}
\colorlet{poppurpleL}{poppurple!20}
\colorlet{popgreenL}{popgreen!20}
\colorlet{poppinkL}{poppink!20}
\colorlet{popblueL}{popblue!20}
\colorlet{popgoldL}{popgold!20}
\colorlet{popredL}{popred!20}
\colorlet{popgrayL}{popgray!20}
% Teintes claires pour fonds de tableaux / zones
\colorlet{popblueL}{popblue!10!white}
\colorlet{poporangeL}{poporange!14!white}
\colorlet{popgreenL}{popgreen!12!white}
\colorlet{poppurpleL}{poppurple!12!white}
\colorlet{poppinkL}{poppink!12!white}
\colorlet{popgoldL}{popgold!14!white}

\pagecolor{popcream}

% ============================================================
% Typographie
% ============================================================
\defaultfontfeatures{Ligatures=TeX}
\setmainfont{PlusJakartaSans-Medium.ttf}[Path=../../../fonts/,BoldFont=PlusJakartaSans-Bold.ttf]
% Maths en sans-serif (Fira Math) avec chiffres et lettres droites de Plus
% Jakarta, pour que les formules se fondent dans le texte.
\usepackage[math-style=ISO,bold-style=ISO]{unicode-math}
\setmathfont{FiraMath-Regular.otf}[Path=../../../fonts/]
\setmathfont{PlusJakartaSans-Medium.ttf}[Path=../../../fonts/,range={up/num,up/{Latin,latin},\mathup}]
\usepackage{icomma} % virgule décimale sans espace en mode math
% Caractères mathématiques Unicode absents des polices de texte (Plus Jakarta,
% Archivo Black) : rendus par leur équivalent mathématique, en texte comme en
% formule — sinon ils s'affichent en carré vide (ex. « Arithmétique dans ℤ »).
\usepackage{newunicodechar}
\newunicodechar{ℝ}{\ensuremath{\mathbb{R}}}
\newunicodechar{ℤ}{\ensuremath{\mathbb{Z}}}
\newunicodechar{ℂ}{\ensuremath{\mathbb{C}}}
\newunicodechar{ℕ}{\ensuremath{\mathbb{N}}}
\newunicodechar{ℚ}{\ensuremath{\mathbb{Q}}}
\newunicodechar{𝔻}{\ensuremath{\mathbb{D}}}
\newunicodechar{α}{\ensuremath{\alpha}}
\newunicodechar{β}{\ensuremath{\beta}}
\newunicodechar{γ}{\ensuremath{\gamma}}
\newunicodechar{δ}{\ensuremath{\delta}}
\newunicodechar{ε}{\ensuremath{\varepsilon}}
\newunicodechar{θ}{\ensuremath{\theta}}
\newunicodechar{λ}{\ensuremath{\lambda}}
\newunicodechar{μ}{\ensuremath{\mu}}
\newunicodechar{σ}{\ensuremath{\sigma}}
\newunicodechar{τ}{\ensuremath{\tau}}
\newunicodechar{φ}{\ensuremath{\varphi}}
\newunicodechar{ψ}{\ensuremath{\psi}}
\newunicodechar{ω}{\ensuremath{\omega}}
\newunicodechar{Σ}{\ensuremath{\Sigma}}
% majuscules produites par \MakeUppercase dans les titres (α-aminés -> Α)
\newunicodechar{Α}{\ensuremath{\alpha}}
\newunicodechar{Β}{\ensuremath{\beta}}
\newunicodechar{Γ}{\ensuremath{\Gamma}}
\newunicodechar{Δ}{\ensuremath{\Delta}}
\newunicodechar{Ε}{\ensuremath{\varepsilon}}
\newunicodechar{Θ}{\ensuremath{\Theta}}
\newunicodechar{Λ}{\ensuremath{\Lambda}}
\newunicodechar{Μ}{\ensuremath{\mu}}
\newunicodechar{Τ}{\ensuremath{\tau}}
\newunicodechar{Φ}{\ensuremath{\Phi}}
\newunicodechar{Ψ}{\ensuremath{\Psi}}
\newunicodechar{Ω}{\ensuremath{\Omega}}
\newunicodechar{↦}{\ensuremath{\mapsto}}
\newunicodechar{∈}{\ensuremath{\in}}
\newunicodechar{∉}{\ensuremath{\notin}}
\newunicodechar{∧}{\ensuremath{\wedge}}
\newunicodechar{⇔}{\ensuremath{\Leftrightarrow}}
\newunicodechar{⟺}{\ensuremath{\Longleftrightarrow}}
\newunicodechar{⇒}{\ensuremath{\Rightarrow}}
\newunicodechar{⟹}{\ensuremath{\Longrightarrow}}
\newunicodechar{∘}{\ensuremath{\circ}}
\newunicodechar{∪}{\ensuremath{\cup}}
\newunicodechar{∩}{\ensuremath{\cap}}
\newunicodechar{≡}{\ensuremath{\equiv}}
\newunicodechar{⊂}{\ensuremath{\subset}}
\newunicodechar{⊥}{\ensuremath{\perp}}
\newunicodechar{∃}{\ensuremath{\exists}}
\newunicodechar{∀}{\ensuremath{\forall}}
\newunicodechar{∛}{\ensuremath{\sqrt[3]{\,}}}
\newunicodechar{∜}{\ensuremath{\sqrt[4]{\,}}}
\newunicodechar{⃗}{\ensuremath{{}^{\to}}}
\newunicodechar{ⁿ}{\ifmmode^{n}\else\textsuperscript{\ensuremath{n}}\fi}
\newunicodechar{ˣ}{\ifmmode^{x}\else\textsuperscript{\ensuremath{x}}\fi}
\newunicodechar{ᵇ}{\ifmmode^{b}\else\textsuperscript{\ensuremath{b}}\fi}
\newunicodechar{ʳ}{\ifmmode^{r}\else\textsuperscript{\ensuremath{r}}\fi}
\newunicodechar{ᵘ}{\ifmmode^{u}\else\textsuperscript{\ensuremath{u}}\fi}
\newunicodechar{ʸ}{\ifmmode^{y}\else\textsuperscript{\ensuremath{y}}\fi}
\newunicodechar{ᵃ}{\ifmmode^{a}\else\textsuperscript{\ensuremath{a}}\fi}
\newunicodechar{ᵉ}{\ifmmode^{e}\else\textsuperscript{\ensuremath{e}}\fi}
\newunicodechar{ᵏ}{\ifmmode^{k}\else\textsuperscript{\ensuremath{k}}\fi}
\newunicodechar{ᵖ}{\ifmmode^{p}\else\textsuperscript{\ensuremath{p}}\fi}
\newunicodechar{ᴺ}{\ifmmode^{N}\else\textsuperscript{\ensuremath{N}}\fi}
\newunicodechar{ᶜ}{\ifmmode^{c}\else\textsuperscript{\ensuremath{c}}\fi}
\newunicodechar{ʰ}{\ifmmode^{h}\else\textsuperscript{\ensuremath{h}}\fi}
\newunicodechar{ᵠ}{\ifmmode^{q}\else\textsuperscript{\ensuremath{q}}\fi}
\newunicodechar{ₙ}{\ifmmode_{n}\else\textsubscript{\ensuremath{n}}\fi}
\newunicodechar{ₐ}{\ifmmode_{a}\else\textsubscript{\ensuremath{a}}\fi}
\newunicodechar{ᵢ}{\ifmmode_{i}\else\textsubscript{\ensuremath{i}}\fi}
\newunicodechar{ᵣ}{\ifmmode_{r}\else\textsubscript{\ensuremath{r}}\fi}
\newunicodechar{ₖ}{\ifmmode_{k}\else\textsubscript{\ensuremath{k}}\fi}
\newunicodechar{ᵦ}{\ifmmode_{\beta}\else\textsubscript{\ensuremath{\beta}}\fi}
\newunicodechar{ₓ}{\ifmmode_{x}\else\textsubscript{\ensuremath{x}}\fi}
\newfontfamily\popdisplay{ArchivoBlack-Regular.ttf}[Path=../../../fonts/]
\newfontfamily\poplogo{SpaceGrotesk-Bold.ttf}[Path=../../../fonts/]
\newfontfamily\popsemi{PlusJakartaSans-SemiBold.ttf}[Path=../../../fonts/]
\newfontfamily\popxbold{PlusJakartaSans-ExtraBold.ttf}[Path=../../../fonts/]
\newfontfamily\popxboldls{PlusJakartaSans-ExtraBold.ttf}[Path=../../../fonts/,LetterSpace=3.0]

\setlist[enumerate]{leftmargin=1.7em,itemsep=5pt,topsep=4pt}
\setlist[itemize]{leftmargin=1.7em,itemsep=4pt,topsep=4pt,label={\color{poporange}\textbullet}}
\setlength{\parindent}{0pt}
\setlength{\parskip}{5pt}
\renewcommand{\arraystretch}{1.25}

% pgfplots : style Pop par défaut
\pgfplotsset{
  every axis/.append style={axis line style={popink,line width=1pt},tick style={popink},
    grid style={popline,line width=0.5pt},label style={font=\small},tick label style={font=\footnotesize}},
  popcurve/.style={poporange,line width=1.8pt,no markers,smooth},
}
% Même style disponible dans un \draw TikZ simple (hors pgfplots)
\tikzset{popcurve/.style={poporange,line width=1.8pt}}
\tikzset{popgrid/.style={popline,very thin}}
\colorlet{popcurve}{poporange} % le nom du style est parfois utilisé comme couleur

% ============================================================
% Pill / squiggle
% ============================================================
\newcommand{\poppill}[3]{%
  \tikz[baseline=-0.55ex]\node[draw=popink,line width=1.2pt,rounded corners=2.6pt,%
    fill=#2,inner xsep=6.5pt,inner ysep=3pt]{%
    \popxboldls\fontsize{7}{7}\selectfont\color{#3}\MakeUppercase{#1}};%
}
\newcommand{\popsquiggle}[1]{%
  \tikz[baseline=-0.4ex]{
    \draw[#1,line width=2.6pt,line cap=round]
      plot [smooth] coordinates {(0,0.09) (0.34,0.01) (0.68,0.21) (1.02,0.01) (1.36,0.10)};
  }%
}
% Mot-clé surligné (marqueur orange sous le texte)
\newcommand{\cle}[1]{{\popxbold\color{popdark}#1}}

% ============================================================
% En-tête / pied
% ============================================================
\pagestyle{fancy}
\fancyhf{}
\renewcommand{\headrulewidth}{0pt}
\renewcommand{\footrulewidth}{0pt}
\lhead{\raisebox{-0.15ex}{\poplogo\fontsize{13}{13}\selectfont\color{popink}OnBuch\color{poporange}+}}
\rhead{\poppill{\DOCMATIERE\ \textbullet\ \DOCNIVEAU\ \textbullet\ Leçon \DOCLECON}{white}{popink}}
\lfoot{\popxbold\fontsize{7.6}{7.6}\selectfont\color{popmuted}\MakeUppercase{Cours signé OnBuch+}\ \textbullet\ {\popsemi\DOCTITRE}}
\rfoot{\tikz[baseline=-0.6ex]\node[rounded corners=8.5pt,fill=popink,minimum height=17pt,minimum width=17pt,inner xsep=5pt,inner sep=0]{\popxbold\fontsize{8}{8}\selectfont\color{popcream}\thepage\,/\,\pageref*{LastPage}};}

% ============================================================
% Titres
% ============================================================
\newcommand{\chaptitre}[2]{%
  \begin{tcolorbox}[enhanced,colback=poporange,colframe=popink,boxrule=2.6pt,arc=11pt,
      drop shadow=popink,left=15pt,right=15pt,top=12pt,bottom=11pt,before skip=3pt,after skip=18pt]
    \poppill{#2}{popink}{popcream}\par\vspace{9pt}
    {\raggedright\hyphenpenalty=10000\exhyphenpenalty=10000\popdisplay\fontsize{22}{24}\selectfont\color{popcream}\MakeUppercase{#1}\par}
  \end{tcolorbox}
}
% Section numérotée : pastille orange avec le numéro + titre Archivo
\newcounter{popsec}
\newcommand{\coursec}[1]{%
  \stepcounter{popsec}\setcounter{popsub}{0}%
  \par\vspace{16pt}\noindent
  \begin{minipage}{\linewidth}
  \tikz[baseline=-0.7ex]\node[draw=popink,line width=1.6pt,fill=poporange,rounded corners=4pt,minimum size=22pt,inner sep=0]{\popdisplay\fontsize{12}{12}\selectfont\color{popcream}\thepopsec};%
  \hspace{8pt}{\popdisplay\fontsize{15}{17}\selectfont\color{popink}\MakeUppercase{#1}}\par
  \vspace{4pt}\noindent{\color{poporange}\rule{36pt}{3.6pt}}
  \end{minipage}\par\nopagebreak\vspace{8pt}
}
\newcounter{popsub}
\newcommand{\courssub}[1]{%
  \stepcounter{popsub}%
  \par\vspace{9pt}\noindent{\popxbold\fontsize{11.5}{13}\selectfont\color{popink}{\color{poporange}\thepopsec.\thepopsub}\hspace{6pt}#1}\par\nopagebreak\vspace{4pt}
}

% ============================================================
% Boîtes pédagogiques — même recette : fond blanc, bordure épaisse colorée,
% ombre dure encre, badge plein. Titre optionnel [#1].
% ============================================================
\tcbset{popbase/.style={breakable,enhanced,colback=white,boxrule=2.3pt,arc=9pt,
  shadow={3pt}{-3pt}{0pt}{popink},left=14pt,right=14pt,top=11pt,bottom=11pt,before skip=11pt,after skip=15pt,
  fontupper=\popsemi\fontsize{10.6}{14.5}\selectfont\color{popink},
  fontlower=\popsemi\fontsize{10.6}{14.5}\selectfont\color{popink}}}
% Coche dessinée (le glyphe ✓ n'existe pas dans les polices du document)
\newcommand{\popcheck}[1]{\tikz[baseline=-0.5ex]\draw[#1,line width=1.6pt,line cap=round,line join=round](-0.13,0.0)--(-0.04,-0.09)--(0.13,0.1);}
\providecommand{\coche}{\popcheck{popgreen}}
\providecommand{\resultat}[1]{\par\smallskip\noindent\fcolorbox{popgreen}{popgreenL}{\parbox{\dimexpr\linewidth-2\fboxsep-2\fboxrule\relax}{\textbf{Résultat.} #1}}\par\smallskip}
\newcommand{\popboxhead}[3]{\poppill{#1}{#2}{white}\ifx\relax#3\relax\else\hspace{6pt}{\popxbold\fontsize{10.6}{12}\selectfont\color{popink}#3}\fi\par\vspace{7pt}}

\newtcolorbox{aretenir}[1][]{popbase,colframe=popgreen,before upper={\popboxhead{À retenir}{popgreen}{#1}}}
% Texte en anglais (document de lecture, dialogue, extrait) : cadre
% d'étude. Un titre optionnel (source, auteur) s'affiche dans l'en-tête.
\newtcolorbox{textebox}[1][]{popbase,colframe=popblue,before upper={\popboxhead{Text}{popblue}{#1}\begin{otherlanguage}{english}},after upper={\end{otherlanguage}}}
% Marques ✓ / ✗ (forme correcte / fautive) souvent utilisées par les modèles
\providecommand{\cross}{\textcolor{poppink}{\ensuremath{\times}}}
\providecommand{\xmark}{\cross}\providecommand{\crossmark}{\cross}\providecommand{\cmark}{\ensuremath{\checkmark}}
% Ligne de réponse / justification à compléter (inventée par les modèles)
\providecommand{\justification}[1]{\par\noindent\textit{Justification~:} \dotfill\par}
% \textipa (package tipa non chargé) : la police ne couvre pas l'API -> simple passage
\providecommand{\textipa}[1]{#1}
% Soulignement (ulem non chargé) : \uline, \ul
\providecommand{\uline}[1]{\underline{#1}}\providecommand{\ul}[1]{\underline{#1}}
% \glossaire{\item ...} inventée par les modèles : liste de glossaire
\providecommand{\glossaire}[1]{\begin{itemize}#1\end{itemize}}
% \alert (beamer) inventée par les modèles : texte mis en évidence
\providecommand{\alert}[1]{\textbf{\textcolor{poppink}{#1}}}
% variantes ASCII d'ordinaux écrites en macro
\providecommand{\iere}{\textsuperscript{re}}\providecommand{\ieme}{\textsuperscript{e}}\providecommand{\ere}{\textsuperscript{re}}\providecommand{\eme}{\textsuperscript{e}}\providecommand{\er}{\textsuperscript{er}}\providecommand{\ie}{\textsuperscript{e}}
% Ordinaux français écrits en macro par les modèles : 1\ière, 3\ième
\newcommand{\ière}{\textsuperscript{re}}\newcommand{\ième}{\textsuperscript{e}}
% \exemple (inventée par les modèles) : étiquette « Exemple : »
\providecommand{\exemple}{\textbf{Exemple~:}\ }
% \square utilisable aussi hors mode math (cases à cocher)
\AtBeginDocument{\let\squareMath\square\renewcommand{\square}{\ensuremath{\squareMath}}}
% Mot ou phrase en anglais à l'intérieur d'un paragraphe français
\newcommand{\eng}[1]{\ifmmode\text{\textit{#1}}\else\begingroup\selectlanguage{english}\itshape #1\endgroup\fi}
\let\esp\eng % alias toléré
% flèches d'intonation inventées par les modèles
\providecommand{\textsearrow}{}\renewcommand{\textsearrow}{\ensuremath{\searrow}}\providecommand{\textnearrow}{}\renewcommand{\textnearrow}{\ensuremath{\nearrow}}
% \fr{..} : mot français (inventé par les modèles)
\providecommand{\fr}[1]{\textit{#1}}
\newtcolorbox{exemplebox}[1][]{popbase,colframe=poporange,before upper={\popboxhead{Exemple}{poporange}{#1}}}
\newtcolorbox{definition}[1][]{popbase,colframe=popblue,before upper={\popboxhead{Définition}{popblue}{#1}}}
\newtcolorbox{propriete}[1][]{popbase,colframe=poppurple,before upper={\popboxhead{Propriété}{poppurple}{#1}}}
\newtcolorbox{methode}[1][]{popbase,colframe=popgold,before upper={\popboxhead{Méthode}{popgold}{#1}}}
\newtcolorbox{attention}[1][]{popbase,colframe=poppink,before upper={\popboxhead{Attention}{poppink}{#1}}}
\newtcolorbox{experience}[1][]{popbase,colframe=popblue,colback=popblueL,before upper={\popboxhead{Expérience}{popblue}{#1}}}
% « Le savais-tu ? » : carte violette pleine (contexte, culture, Cameroun)
\newtcolorbox{savaistu}[1][]{popbase,colback=poppurple,colframe=popink,
  fontupper=\popsemi\fontsize{10.4}{14.2}\selectfont\color{white},
  before upper={\poppill{Le savais-tu\,?}{white}{poppurple}\ifx\relax#1\relax\else\hspace{6pt}{\popxbold\color{white}#1}\fi\par\vspace{7pt}}}
% Exercice résolu : énoncé en haut, \tcblower puis la solution sur fond crème
\newcounter{popexo}\newcounter{popexr}
\newtcolorbox{exoresolu}[1][]{popbase,colframe=poporange,colbacklower=popcream,
  segmentation style={popink,line width=1.2pt,dashed},
  before upper={\stepcounter{popexr}\popboxhead{Exercice résolu \thepopexr}{poporange}{#1}},
  before lower={\poppill{Solution}{popink}{popcream}\par\vspace{6pt}}}
% Exercice (non résolu en place) — la correction est dans la partie Corrigés
\newtcolorbox{exercice}[1][]{popbase,colframe=popink,
  before upper={\stepcounter{popexo}\popboxhead{Exercice \thepopexo}{popink}{#1}}}
% Corrigé
\newtcolorbox{corrige}[1][]{popbase,colframe=popgreen,colback=popgreenL,
  before upper={\popboxhead{Corrigé}{popgreen}{#1}}}
% Objectifs / prérequis / bilan
\newtcolorbox{objectifs}{popbase,colframe=popink,colback=poporangeL,
  before upper={\popboxhead{Objectifs de la leçon}{popink}{}}}
\newtcolorbox{prerequis}{popbase,colframe=popink,colback=popblueL,
  before upper={\popboxhead{Prérequis}{popblue}{}}}
\newtcolorbox{bilan}{popbase,colframe=popink,colback=white,boxrule=2.8pt,
  before upper={\popboxhead{Fiche bilan}{poporange}{L'essentiel à connaître pour l'examen}}}
% Figure encadrée avec légende
\newtcolorbox{popfigure}[1][]{enhanced,breakable=false,colback=white,colframe=popline,boxrule=1.4pt,arc=8pt,
  left=10pt,right=10pt,top=10pt,bottom=8pt,before skip=10pt,after skip=12pt,halign=center,
  lower separated=false,halign lower=center,
  fontlower=\popsemi\fontsize{9}{11.5}\selectfont\color{popmuted},
  #1}
\newcommand{\legende}[1]{\par\vspace{6pt}{\popsemi\fontsize{9}{11.5}\selectfont\color{popmuted}#1}}

% ============================================================
% Signature OnBuch+
% ============================================================
\newcommand{\poplogoinline}[1]{{\poplogo\fontsize{#1}{#1}\selectfont\color{popink}OnBuch\color{poporange}+}}
\newcommand{\signatureonbuch}{%
  \begin{tcolorbox}[enhanced,colback=popink,colframe=popink,boxrule=2.2pt,arc=10pt,
      drop shadow=poporange,left=14pt,right=14pt,top=10pt,bottom=10pt,before skip=0pt,after skip=0pt]
    \tikz[baseline=-0.7ex]\node[circle,fill=popgreen,draw=popcream,line width=1.3pt,minimum size=22pt,inner sep=0]{\popcheck{white}};%
    \hspace{9pt}\begin{minipage}[c]{0.62\linewidth}
      {\popxbold\fontsize{12}{13}\selectfont\color{popcream}Signé\ \textbullet\ {\poplogo OnBuch\color{poporange}+}}\par\vspace{2pt}
      {\raggedright\popsemi\fontsize{8}{10.5}\selectfont\color{popline}\MakeUppercase{Équipe pédagogique OnBuch+}\\\MakeUppercase{Conforme au programme officiel MINESEC}\par}
    \end{minipage}\hfill
    {\poplogo\fontsize{22}{22}\selectfont\color{popcream}OnBuch\color{poporange}+}
  \end{tcolorbox}%
}

% ============================================================
% Page de garde
% #1 titre · #2 matière · #3 niveau · #4 module · #5 n° leçon · #6 durée ·
% #7 description / objectifs (texte ou itemize)
% ============================================================
\newcommand{\pagegarde}[7]{%
  \thispagestyle{empty}
  \newgeometry{a4paper,margin=1.7cm,top=1.6cm,bottom=1.4cm}%
  \begin{tikzpicture}[remember picture,overlay]
    \fill[poporange!18] ([xshift=-3.2cm,yshift=-3.6cm]current page.north east) circle (4.6cm);
    \fill[poppurple!14] ([xshift=2.4cm,yshift=7.6cm]current page.south west) circle (3.8cm);
  \end{tikzpicture}%
  {\poplogo\fontsize{30}{30}\selectfont\color{popink}OnBuch\color{poporange}+}\hfill
  \raisebox{0.6ex}{\poppill{Cours signé \textbullet\ Premium}{popink}{popcream}}\par
  \vspace{1pt}\popsquiggle{poppurple}\par
  \vspace{8pt}
  \begin{tcolorbox}[enhanced,colback=poporange,colframe=popink,boxrule=2.8pt,arc=13pt,
      drop shadow=popink,left=18pt,right=18pt,top=12pt,bottom=13pt,after skip=10pt]
    \poppill{#2 \textbullet\ #3}{popink}{popcream}\hspace{5pt}\poppill{#4}{white}{popink}\par\vspace{8pt}
    {\popdisplay\fontsize{14}{14}\selectfont\color{popink}LEÇON #5}\par\vspace{4pt}
    {\raggedright\hyphenpenalty=10000\exhyphenpenalty=10000\popdisplay\fontsize{25}{27}\selectfont\color{popcream}\MakeUppercase{#1}\par}
    \vspace{8pt}
    {\popxbold\fontsize{9.5}{9.5}\selectfont\color{popink}\MakeUppercase{Durée conseillée : #6}}
  \end{tcolorbox}
  \begin{tcolorbox}[enhanced,colback=white,colframe=popink,boxrule=2.2pt,arc=10pt,
      drop shadow=popink,left=16pt,right=16pt,top=10pt,bottom=10pt,after skip=8pt]
    \poppill{Dans cette leçon}{poppurple}{white}\par\vspace{5pt}
    \popsemi\fontsize{9.3}{12.6}\selectfont\color{popink}#7
  \end{tcolorbox}
  \noindent\begin{minipage}[t]{0.487\linewidth}
    \begin{tcolorbox}[enhanced,colback=popgreen,colframe=popink,boxrule=2.2pt,arc=10pt,
        drop shadow=popink,left=11pt,right=11pt,top=10pt,bottom=10pt,height=3.1cm,valign=center]
      \tcbox[colback=white,colframe=popink,boxrule=1.3pt,arc=5pt,boxsep=0pt,nobeforeafter,
        left=3pt,right=3pt,top=3pt,bottom=3pt,valign=center]{\qrcode[height=1.9cm]{https://play.google.com/store/apps/details?id=com.luvvix.onbuch&hl=fr}}%
      \hspace{8pt}\begin{minipage}[c]{0.5\linewidth}
        {\popdisplay\fontsize{11}{12.5}\selectfont\color{white}\MakeUppercase{Télécharge OnBuch}}\par\vspace{4pt}
        {\raggedright\popsemi\fontsize{8.4}{11}\selectfont\color{white}Gratuit \textbullet\ Google Play\\Tuteur IA, annales, concours, résultats\par}
      \end{minipage}
    \end{tcolorbox}
  \end{minipage}\hfill
  \begin{minipage}[t]{0.487\linewidth}
    \begin{tcolorbox}[enhanced,colback=poppurple,colframe=popink,boxrule=2.2pt,arc=10pt,
        drop shadow=popink,left=11pt,right=11pt,top=10pt,bottom=10pt,height=3.1cm,valign=center]
      \tikz[baseline=-0.7ex]\node[circle,fill=white,draw=popink,line width=1.3pt,minimum size=1.6cm,inner sep=0]{\popdisplay\fontsize{24}{24}\selectfont\color{poppurple}+};%
      \hspace{8pt}\begin{minipage}[c]{0.6\linewidth}
        {\popdisplay\fontsize{11}{12.5}\selectfont\color{white}\MakeUppercase{Passe à OnBuch+}}\par\vspace{4pt}
        {\raggedright\popsemi\fontsize{8.4}{11}\selectfont\color{white}Tous les cours signés, Tuteur IA illimité, fiches PDF — onglet OnBuch+ de l'app\par}
      \end{minipage}
    \end{tcolorbox}
  \end{minipage}
  \vspace{8pt}
  \popcontenu
  \vfill
  \signatureonbuch
  \restoregeometry
  \newpage
}

% Rangée « Ce cours contient » (page de garde)
\newcommand{\poptile}[3]{% #1 couleur #2 gros texte #3 légende
  \begin{tcolorbox}[enhanced,colback=#1,colframe=popink,boxrule=2pt,arc=9pt,shadow={2.5pt}{-2.5pt}{0pt}{popink},
      left=6pt,right=6pt,top=9pt,bottom=9pt,halign=center,valign=center,height=1.75cm,width=\linewidth,nobeforeafter]
    {\popdisplay\fontsize{12.5}{13}\selectfont\color{white}\MakeUppercase{#2}}\par\vspace{3pt}
    {\centering\popsemi\fontsize{7.8}{9.5}\selectfont\color{white}#3\par}
  \end{tcolorbox}}
\newcommand{\popcontenu}{%
  \par\noindent{\popxboldls\fontsize{7.5}{7.5}\selectfont\color{popmuted}CE COURS SIGNÉ CONTIENT}\par\vspace{7pt}
  \noindent\begin{minipage}[t]{0.235\linewidth}\poptile{popblue}{Cours}{complet, illustré, pas à pas}\end{minipage}\hfill
  \begin{minipage}[t]{0.235\linewidth}\poptile{poporange}{Méthodes}{et exercices résolus}\end{minipage}\hfill
  \begin{minipage}[t]{0.235\linewidth}\poptile{poppink}{Exercices}{type examen + corrigés}\end{minipage}\hfill
  \begin{minipage}[t]{0.235\linewidth}\poptile{popgreen}{Fiche bilan}{l'essentiel à retenir}\end{minipage}\par}

% Sommaire
\newtcolorbox{sommaire}{enhanced,colback=white,colframe=popink,boxrule=2.2pt,arc=10pt,
  drop shadow=popink,left=18pt,right=18pt,top=13pt,bottom=13pt,before skip=6pt,after skip=20pt,
  before upper={\poppill{Sommaire}{poppurple}{white}\par\vspace{9pt}\popsemi\fontsize{10.6}{14.5}\selectfont\color{popink}}}

% Signature de fin de cours — poussée tout en bas de la dernière page
\newcommand{\findecours}{%
  \par\vfill\nopagebreak
  \begin{center}\popsquiggle{poporange}\end{center}\vspace{4pt}
  \signatureonbuch
}
\AtBeginDocument{\def\llbracket{\mathopen{[\mkern-3.5mu[}}\def\rrbracket{\mathclose{]\mkern-3.5mu]}}} % intervalles d'entiers (glyphe ⟦ absent de la police)
% Glyphes absents de Fira Math : redéfinis à partir de glyphes disponibles
\AtBeginDocument{%
  \def\setminus{\mathbin{\backslash}}\let\smallsetminus\setminus
  \def\vdots{\mathord{\vbox{\baselineskip3.2pt\lineskiplimit0pt\kern4pt\hbox{.}\hbox{.}\hbox{.}}}}%
  \def\ddots{\mathinner{\mkern1mu\raise6.5pt\hbox{.}\mkern2mu\raise3.6pt\hbox{.}\mkern2mu\raise0.7pt\hbox{.}\mkern1mu}}%
  \def\triangle{\mathord{\scalebox{0.9}{$\symup{\Delta}$}}}%
  \def\nabla{\mathord{\raisebox{\depth}{\scalebox{1}[-1]{$\symup{\Delta}$}}}}%
  \def\checkmark{\popcheck{popdark}}%
  \def\star{\mathbin{\ast}}\def\bigstar{\mathord{\ast}}%
  \def\diamond{\mathbin{\scalebox{0.6}{$\Diamond$}}}%
  \def\bigcup{\mathop{\mathchoice{\raisebox{-0.3em}{\scalebox{1.7}{$\cup$}}}{\scalebox{1.25}{$\cup$}}{\cup}{\cup}}\displaylimits}%
  \def\bigcap{\mathop{\mathchoice{\raisebox{-0.3em}{\scalebox{1.7}{$\cap$}}}{\scalebox{1.25}{$\cap$}}{\cap}{\cap}}\displaylimits}%
  \def\lesseqqgtr{\mathrel{\lessgtr}}\def\bigoplus{\mathop{\oplus}\displaylimits}%
}
\providecommand{\ssi}{si et seulement si} % abréviation fréquente
\AtBeginDocument{\providecommand{\wideparen}{\overparen}} % arc de cercle

%% FIGFIX-BEGIN v3 — correctifs globaux des figures (docs/onbuch-generation/tools/figfix)
\usepackage{adjustbox}\usepackage{etoolbox}
% toute figure TikZ/pgfplots plus large que la ligne est réduite à la largeur disponible
\BeforeBeginEnvironment{tikzpicture}{\begin{adjustbox}{max width=\linewidth}}
\AfterEndEnvironment{tikzpicture}{\end{adjustbox}}
% légendes pgfplots lisibles par-dessus les courbes
\pgfplotsset{every axis/.append style={legend style={fill=white,fill opacity=0.92,text opacity=1,draw=black!15,inner sep=3pt}}}
% étiquettes d'arêtes (syntaxe edge["..."]) avec fond blanc
\tikzset{every edge quotes/.append style={fill=white,inner sep=1.2pt}}
% bulles de cartes mentales : texte à la ligne dans la bulle, bulle un peu plus grande
\tikzset{every concept/.append style={text width=2.0cm,align=center,minimum size=2.3cm,inner sep=2pt}}
%% FIGFIX-END
\begin{document}

\pagegarde{Lesson 7 — Reading: Texts about settling community disputes}{Anglais}{1ère A}{Chapter 1 : Family and Social Life}{ENA07}{2 h}{Dans cette leçon de compréhension écrite, les élèves lisent un texte authentique sur le règlement des conflits communautaires (médiation traditionnelle, chefs, conseils de quartier). Ils s'entraînent aux stratégies de lecture (skimming, scanning, déduction du sens) et enrichissent leur lexique sur la paix et la résolution des conflits.
\vspace{4pt}
\begin{multicols}{2}\small
\begin{itemize}
\item À la fin de cette leçon, je sais identifier l'idée générale d'un texte grâce au skimming.
\item À la fin de cette leçon, je sais repérer des informations précises grâce au scanning.
\item À la fin de cette leçon, je sais déduire le sens d'un mot inconnu à partir du contexte.
\item À la fin de cette leçon, je sais utiliser le lexique du règlement des conflits (dispute, mediator, settle, agreement...).
\item À la fin de cette leçon, je sais répondre à des questions de compréhension de types variés (vrai/faux, QCM, réponses courtes).
\item À la fin de cette leçon, je sais justifier mes réponses en citant le texte.
\end{itemize}\end{multicols}}

\begin{sommaire}
\begin{multicols}{2}
\begin{enumerate}
\item Warm-up and pre-reading
\item Reading text with glossary
\item Reading strategies: skimming, scanning, guessing
\item Comprehension questions
\item Thematic vocabulary: peace and conflict
\item Guided production: summary and opinion
\item Activité d'intégration
\item Méthodes et exercices résolus
\item Exercices
\item Corrigés des exercices
\item Fiche bilan
\end{enumerate}
\end{multicols}
\end{sommaire}

\begin{objectifs}
\begin{itemize}
\item À la fin de cette leçon, je sais identifier l'idée générale d'un texte grâce au skimming.
\item À la fin de cette leçon, je sais repérer des informations précises grâce au scanning.
\item À la fin de cette leçon, je sais déduire le sens d'un mot inconnu à partir du contexte.
\item À la fin de cette leçon, je sais utiliser le lexique du règlement des conflits (dispute, mediator, settle, agreement...).
\item À la fin de cette leçon, je sais répondre à des questions de compréhension de types variés (vrai/faux, QCM, réponses courtes).
\item À la fin de cette leçon, je sais justifier mes réponses en citant le texte.
\item À la fin de cette leçon, je sais résumer un texte en quelques phrases avec mes propres mots.
\item À la fin de cette leçon, je sais donner mon opinion sur un conflit et proposer une solution pacifique.
\end{itemize}
\end{objectifs}

\begin{prerequis}
\begin{itemize}
\item Le lexique de base de la famille et de la vie sociale (Seconde, Chapter Family Life).
\item Les temps simples : present simple et past simple (Seconde).
\item Les mots interrogatifs : who, what, where, when, why, how (Seconde).
\item Les connecteurs simples : and, but, because, so (début d'année de Première).
\item La lecture de textes courts et la réponse à des questions WH- (Seconde).
\end{itemize}
\end{prerequis}

\newpage

\coursec{Warm-up and Pre-reading}

Imagine two families in Bamenda quarrelling over a piece of farmland for months. Instead of going to court, the village chief invites both families under the big mango tree, listens to each side, and helps them reach an agreement. In Cameroon, as in Ghana, Nigeria or the United Kingdom, communities have always found ways to settle disputes peacefully. Today, we are going to read a text about how communities solve their conflicts — and learn to read it like real detectives.

\courssub{Brainstorming: Types of Community Conflicts}

\begin{textebox}[Discussion]
In your village or quarter, who settles disputes? How? Discuss with your partner and note down three examples of conflicts you have seen or heard about.
\end{textebox}

\begin{exemplebox}[Brainstorming — What conflicts exist in a community?]
Work in small groups. Complete the mind map below with as many types of disputes as you can think of. Use the words in the box to help you.
\begin{center}
\begin{tabular}{|l|l|l|l|l|}\hline
land & noise & water & inheritance & money \\ \hline
football & animals & boundaries & parking & family \\ \hline
\end{tabular}
\end{center}
\end{exemplebox}

\begin{popfigure}
\begin{tikzpicture}[
  mindmap,
  concept color=popblue,
  level 1/.style={sibling angle=72, level distance=3.5cm},
  level 2/.style={sibling angle=45, level distance=2.5cm},
  every node/.style={concept, execute at begin node=\strut},
  text=white,
  font=\sffamily\bfseries
]
\node[align=center]{COMMUNITY\\ DISPUTES}
  child[concept color=poporange] { node[align=center]{LAND\\ DISPUTES} }
  child[concept color=popgreen] { node[align=center]{MONEY\\ ISSUES} }
  child[concept color=poppurple] { node[align=center]{NOISE\\ NUISANCE} }
  child[concept color=poppink] { node[align=center]{WATER\\ SHARING} }
  child[concept color=popgold] { node[align=center]{INHERITANCE\\ CONFLICTS} };
\end{tikzpicture}
\legende{Carte mentale : types de conflits communautaires}
\end{popfigure}

\begin{propriete}[Vocabulary: Types of Disputes]
Voici le lexique essentiel pour nommer les conflits courants. Observez la prononciation approximative entre guillemets.
\begin{center}
\begin{tabularx}{\linewidth}{|X|X|X|}\hline
\rowcolor{popblueL} \textbf{English} & \textbf{Français} & \textbf{Exemple / Contexte} \\ \hline
\eng{land dispute} « lend dis-pyou-t » & conflit foncier & \eng{A land dispute between two neighbours.} (Un conflit foncier entre deux voisins.) \\ \hline
\eng{inheritance conflict} « in-è-ri-tens kon-flikt » & conflit d'héritage & \eng{An inheritance conflict tore the family apart.} (Un conflit d'héritage a déchiré la famille.) \\ \hline
\eng{noise nuisance} « noiz nyu-sens » & nuisance sonore / tapage & \eng{Loud music at night is a noise nuisance.} (La musique forte la nuit est un tapage.) \\ \hline
\eng{water sharing} « wo-teu chè-ring » & partage de l'eau & \eng{Water sharing disputes are common in dry season.} (Les disputes sur le partage de l'eau sont fréquentes en saison sèche.) \\ \hline
\eng{boundary disagreement} « boun-da-ri dis-a-gri-ment » & désaccord sur les limites & \eng{A boundary disagreement over the fence.} (Un désaccord sur la clôture / limite.) \\ \hline
\eng{football quarrel} « foot-bol kwo-rel » & querelle de football & \eng{A football quarrel escalated into a fight.} (Une querelle de football a dégénéré en bagarre.) \\ \hline
\end{tabularx}
\end{center}
\end{propriete}

\begin{exoresolu}[Classifying Disputes]
\textbf{Énoncé :} Classify the following situations into the correct category from the mind map.
\begin{enumerate}
\item Two brothers argue over their father's cocoa farm after his death.
\item A neighbour plays loud music every night until 2 a.m.
\item Villagers disagree on the schedule for the communal standpipe.
\item A shopkeeper refuses to pay back a loan to the \eng{njangi} group.
\item A farmer plants plantains on land his neighbour claims belongs to him.
\end{enumerate}
\tcblower
\textbf{Solution détaillée :}
\begin{enumerate}
\item \textbf{Inheritance conflict} — \eng{father's cocoa farm after his death} indique un héritage.
\item \textbf{Noise nuisance} — \eng{loud music every night until 2 a.m.} = tapage nocturne.
\item \textbf{Water sharing} — \eng{communal standpipe} (borne fontaine commune) + \eng{schedule} (horaire).
\item \textbf{Money issues} — \eng{refuses to pay back a loan} + \eng{njangi} (tontine).
\item \textbf{Land dispute} — \eng{plants plantains on land his neighbour claims} = conflit de propriété foncière.
\end{enumerate}
\end{exoresolu}

\courssub{Who Settles Disputes in Cameroon?}

\begin{textebox}[Traditional and Modern Authorities]
Read the dialogue between two students, Ako and Mireille, discussing how disputes are handled in their quarter.
\eng{
Ako: "In my quarter in Buea, when there is a problem, we first go to the Quarter Head. He calls the elders and they sit under the 'palaver tree' to listen to both sides."
Mireille: "In Yaoundé, it's similar. My uncle had a land dispute with his neighbour. The Chief of the village invited them to his palace. They swore on the Bible and agreed on the boundary."
Ako: "And if they don't agree?"
Mireille: "Then it goes to the Court of First Instance or the police. But most times, the traditional way works better. It keeps the peace."
}
\end{textebox}

\begin{center}
\begin{tabularx}{\linewidth}{|X|X|X|}\hline
\rowcolor{popblueL} \textbf{Authority (English)} & \textbf{Autorité (Français)} & \textbf{Rôle typique} \\ \hline
\eng{the Quarter Head} & le chef de quartier & Premier niveau, médiation de proximité \\ \hline
\eng{the Village Chief / Traditional Ruler} & le chef de village / chef traditionnel & Autorité coutumière, serments, terres \\ \hline
\eng{the Council of Elders} & le conseil des sages / notables & Conseil, arbitrage, mémoire collective \\ \hline
\eng{the Court of First Instance} & le tribunal de première instance & Jugement formel, droit écrit \\ \hline
\eng{the Police / Gendarmerie} & la police / gendarmerie & Maintien de l'ordre, enquêtes, exécution \\ \hline
\end{tabularx}
\end{center}

\begin{savaistu}[Le saviez-vous ?]
In many Cameroonian communities, the \eng{quarter head} or the traditional council settles disputes before they ever reach a court of law. This is called \eng{alternative dispute resolution (ADR)}. The \eng{palaver tree} (or \eng{moot tree}) is a traditional meeting place where the community gathers to talk — \eng{palaver} comes from Portuguese \eng{palavra} (word/speech). In the UK, similar community mediation exists through \eng{Neighbourhood Mediation Services}.
\end{savaistu}

\courssub{Activating Prior Vocabulary}

\begin{definition}[Key Words Activation]
Avant de lire, réactivons le vocabulaire de base que vous connaissez déjà. Complétez le tableau avec la traduction et un synonyme simple en anglais.
\end{definition}

\begin{center}
\begin{tabularx}{\linewidth}{|X|X|X|}\hline
\rowcolor{popblueL} \textbf{Mot (English)} & \textbf{Traduction} & \textbf{Synonyme simple / Définition} \\ \hline
\eng{family} & famille & \eng{relatives, household} \\ \hline
\eng{problem} & problème & \eng{issue, difficulty, trouble} \\ \hline
\eng{talk} (v./n.) & parler / discussion & \eng{speak, discuss, conversation} \\ \hline
\eng{help} (v./n.) & aider / aide & \eng{assist, support, assistance} \\ \hline
\eng{peace} & paix & \eng{harmony, calm, truce} \\ \hline
\eng{settle} (v.) & régler, installer & \eng{resolve, solve, agree} \\ \hline
\eng{dispute} (n.) & dispute, litige & \eng{conflict, argument, quarrel} \\ \hline
\eng{agreement} (n.) & accord & \eng{deal, understanding, consensus} \\ \hline
\eng{listen} (v.) & écouter & \eng{hear attentively, pay attention} \\ \hline
\eng{chief} (n.) & chef & \eng{leader, head, ruler} \\ \hline
\end{tabularx}
\end{center}

\begin{exoresolu}[Word Families]
\textbf{Énoncé :} Complete the sentences with the correct form of the word in brackets.
\begin{enumerate}
\item The \_\_\_\_\_\_\_ (chief) invited both families to his palace.
\item We need a peaceful \_\_\_\_\_\_\_ (settle) of this matter.
\item The elders \_\_\_\_\_\_\_ (listen) carefully to each side.
\item They reached an \_\_\_\_\_\_\_ (agree) after three hours.
\item It is a serious land \_\_\_\_\_\_\_ (dispute).
\end{enumerate}
\tcblower
\textbf{Solution détaillée :}
\begin{enumerate}
\item \textbf{chief} (nom, invariable au singulier) — \eng{The chief invited...}
\item \textbf{settlement} (nom) — \eng{peaceful settlement}. Attention : \eng{settle} (verbe) $\to$ \eng{settlement} (nom).
\item \textbf{listened} (passé simple) — \eng{The elders listened...}
\item \textbf{agreement} (nom) — \eng{reached an agreement}. \eng{Agree} est le verbe.
\item \textbf{dispute} (nom, invariable) — \eng{land dispute}. Le verbe est \eng{dispute} aussi (\eng{they dispute the land}).
\end{enumerate}
\end{exoresolu}

\begin{attention}[Pièges fréquents — Faux amis et erreurs de forme]
\begin{itemize}
\item \textbf{\eng{Dispute} $\neq$ \eng{Disputer}} : \eng{Dispute} est un nom (litige) OU un verbe (contester). \eng{Disputer} n'existe pas en anglais. Pour « se disputer » (quereller), on utilise \eng{argue}, \eng{quarrel} ou \eng{fight}.
\item \textbf{\eng{Chief} vs \eng{Chef}} : \eng{Chief} = chef (autorité). \eng{Chef} (français) = \eng{cook} / \eng{chef} (cuisinier) en anglais.
\item \textbf{\eng{Settle} vs \eng{Install}} : \eng{Settle a dispute} = régler un conflit. \eng{Settle in a place} = s'installer. Ne confondez pas avec \eng{install} (installer un logiciel, un équipement).
\item \textbf{\eng{Agreement} (nom) vs \eng{Agree} (verbe)} : \eng{*They agreement} $\times$ $\to$ \eng{They agree} (verbe) ou \eng{They reached an agreement} (nom).
\item \textbf{\eng{Listen} + to} : \eng{listen to someone} (écouter qqn). Oublier \eng{to} est l'erreur n°1 : \eng{*listen someone} $\times$.
\end{itemize}
\end{attention}

\courssub{Predicting Text Content: The Detective Method}

\begin{methode}[How to predict what a text is about — 3 Steps]
Avant de lire un texte en détail, les bons lecteurs font des prédictions. Suivez ces étapes :
\begin{enumerate}
\item \textbf{Read the title and subtitle.} Identify key words. Ask: “What is the topic?” “What is the tone?” (serious, informative, narrative?).
\item \textbf{Look at the images, photos, or layout.} What do you see? People talking? A tree? A document? A courtroom? This gives the \eng{setting} (cadre) and the \eng{participants}.
\item \textbf{Read the first sentence of each paragraph.} This is often the \eng{topic sentence}. It tells you the main idea of each section. Write one-word summaries in the margin.
\end{enumerate}
\end{methode}

\begin{exemplebox}[Prediction Practice]
Imagine the text we will read has this title and picture:
\begin{center}
\textbf{“Under the Mango Tree: How Bafut Solves Land Disputes”}
\end{center}
\textbf{Picture:} An old chief sitting under a large tree, surrounded by men and women. One man points to a paper; another shakes hands with his neighbour.
\par\medskip
\textbf{Your predictions:}
\begin{enumerate}
\item Topic: \eng{Traditional land dispute resolution in Bafut.}
\item Setting: \eng{Outdoor, under a mango tree, village atmosphere.}
\item Key actors: \eng{The chief, elders, two disputing parties.}
\item Likely vocabulary: \eng{boundary, oath, witness, compromise, ancestor, tradition.}
\item Possible structure: Paragraph 1 = the conflict; Paragraph 2 = the meeting; Paragraph 3 = the resolution; Paragraph 4 = the significance today.
\end{enumerate}
\end{exemplebox}

\courssub{Reading Strategies: Skimming and Scanning}

\begin{propriete}[Reading Strategies: Skimming vs Scanning]
Deux stratégies complémentaires pour lire efficacement.
\begin{center}
\begin{tabularx}{\linewidth}{|X|X|X|}\hline
\rowcolor{popblueL} \textbf{Stratégie} & \textbf{Objectif} & \textbf{Comment faire} \\ \hline
\textbf{Skimming} (lecture rapide / survol) & Saisir l'\textbf{idée générale} (gist), le sujet principal, la structure du texte. & \begin{itemize}\item Lire le titre, l'intro, la conclusion.\item Lire la 1ère phrase de chaque paragraphe.\item Ignorer les mots inconnus.\item Ne pas lire chaque mot.\end{itemize} \\ \hline
\textbf{Scanning} (recherche ciblée) & Trouver une \textbf{information précise} : un nom, une date, un chiffre, un mot-clé. & \begin{itemize}\item Garder la question en tête.\item Faire courir les yeux sur le texte.\item Chercher des indices visuels (majuscules, chiffres, mots en gras).\item S'arrêter dès qu'on trouve.\end{itemize} \\ \hline
\end{tabularx}
\end{center}
\end{propriete}

\begin{exemplebox}[When to use which?]
\begin{itemize}
\item \textbf{Skimming :} “What is this text mainly about?” / “Choose the best title.” / “How many parts does the text have?”
\item \textbf{Scanning :} “When did the dispute start?” / “What is the name of the chief?” / “Which paragraph mentions the police?” / “Find a word that means ‘agreement’.”
\end{itemize}
\end{exemplebox}

\begin{exoresolu}[Skimming vs Scanning — Decision]
\textbf{Énoncé :} For each question, decide if you need to SKIM or SCAN. Write S (Skimming) or C (Scanning).
\begin{enumerate}
\item What is the main idea of the text?
\item In which year did the Bafut chief settle the big land case?
\item How many paragraphs talk about the role of women?
\item Find the word that means “official decision by a court”.
\item What is the author's opinion on traditional justice?
\end{enumerate}
\tcblower
\textbf{Solution détaillée :}
\begin{enumerate}
\item \textbf{S} (Skimming) — idée générale.
\item \textbf{C} (Scanning) — date précise (chiffre).
\item \textbf{C} (Scanning) — information localisée (compter / chercher “women”).
\item \textbf{C} (Scanning) — mot précis (synonyme).
\item \textbf{S} (Skimming) — opinion globale = lecture d'ensemble.
\end{enumerate}
\end{exoresolu}

\courssub{Reading Text: Under the Mango Tree}

\begin{textebox}[Reading Text: Under the Mango Tree — A Bafut Story]
\textbf{Title:} Under the Mango Tree — A Bafut Story
\par\medskip
In the heart of Bafut, a dispute had festered for two years. Two families, the Nfor and the Tuma, claimed the same strip of fertile land along the Menchum River. Each family planted maize and beans; each harvest season brought shouting, threats, and torn crops. The Quarter Head, Pa Nkwenti, had tried to mediate, but neither side would yield. “This land belonged to my grandfather,” insisted Old Man Nfor, waving a faded colonial document. “The river changed course,” countered Tuma’s son, “the boundary moved with it.”
\par\medskip
One harmattan morning, the Fon of Bafut sent his messenger: “The Fon requests your presence under the Great Mango Tree.” By noon, a semi-circle of elders, women, and curious children formed around the two families. The Fon, wrapped in his \eng{ndop} cloth, did not speak first. He let the wind stir the leaves. Then Old Man Nfor spoke, his voice cracking. Tuma’s son followed, calm but firm. The elders asked questions: “Where are the ancient markers?” “Who buried the placenta of the firstborn here?” “What did the white man’s map truly say?”
\par\medskip
After three hours, the Fon stood. “The river is the witness,” he declared. “Where the water flows today, the boundary lies. But the Nfor family keeps the mango trees they planted; the Tuma family keeps the yam barns they built. Peace is worth more than a strip of soil.” Hands shook. Kola nuts were shared. The dispute ended not in a courtroom, but under a tree, with words, respect, and the weight of tradition.
\end{textebox}

\begin{definition}[Glossary — Mots difficiles du texte]
\begin{center}
\begin{tabularx}{\linewidth}{|X|X|X|}\hline
\rowcolor{popblueL} \textbf{Mot (English)} & \textbf{Nature} & \textbf{Traduction / Explication} \\ \hline
\eng{fester} & verbe & pourrir, s'aggraver (pour un conflit, une blessure) \\ \hline
\eng{fertile} & adjectif & fertile, riche (sol) \\ \hline
\eng{yield} & verbe & céder, abandonner (sa position) \\ \hline
\eng{faded} & adjectif & décoloré, effacé (document, couleur) \\ \hline
\eng{counter} & verbe & répliquer, opposer un argument \\ \hline
\eng{harmattan} & nom & harmattan (vent sec et poussiéreux, nov.–mars) \\ \hline
\eng{mediator} & nom & médiateur (ici : Pa Nkwenti, Quarter Head) \\ \hline
\eng{ndop} & nom & tissu traditionnel teint à l'indigo (Bamileke/Bafut) \\ \hline
\eng{placenta} & nom & placenta (coutume : enterrer celui du premier-né lie la famille à la terre) \\ \hline
\eng{witness} & nom & témoin (ici : la rivière comme témoin naturel) \\ \hline
\eng{kola nut} & nom & noix de kola (symbole d'hospitalité, d'accord, de paix) \\ \hline
\end{tabularx}
\end{center}
\end{definition}

\courssub{Comprehension Questions}

\begin{exercice}[Reading Comprehension Check]
Answer these questions in full sentences based on the text “Under the Mango Tree”.
\begin{enumerate}
\item How long had the dispute between the Nfor and Tuma families lasted?
\item Why did the Quarter Head's mediation fail?
\item What object did Old Man Nfor use as proof of ownership?
\item Where did the Fon ask the families to meet?
\item What natural element did the Fon use as the final boundary marker?
\item What did each family keep in the final agreement?
\item What verb in paragraph 3 shows the Fon making a final, authoritative decision?
\end{enumerate}
\end{exercice}

\begin{corrige}[Exercice Reading Comprehension Check]
\begin{enumerate}
\item The dispute had lasted for two years.
\item The mediation failed because neither side would yield (give up their claim).
\item He used a faded colonial document.
\item Under the Great Mango Tree.
\item The river (the Menchum River) — “Where the water flows today, the boundary lies.”
\item The Nfor family kept the mango trees; the Tuma family kept the yam barns.
\item The verb is \eng{declared} (“he declared”). It performs the act of judgement.
\end{enumerate}
\end{corrige}

\courssub{Grammar Spotlight: Narrative Tenses}

\begin{propriete}[Narrative Tenses: Past Simple, Past Perfect, Past Continuous]
Les textes narratifs (récits) comme “Under the Mango Tree” utilisent un système de temps précis pour organiser la chronologie.
\begin{center}
\begin{tabularx}{\linewidth}{|X|X|X|}\hline
\rowcolor{popblueL} \textbf{Temps} & \textbf{Forme} & \textbf{Emploi narratif} \\ \hline
\textbf{Past Simple} & V-ed / 2ème colonne & Actions principales, successives, terminées dans le passé. \eng{The Fon stood. He declared.} \\ \hline
\textbf{Past Perfect} & had + V-en / 3ème colonne & Antériorité : action antérieure à une autre action passée. \eng{A dispute \textbf{had festered} for two years (before the meeting).} \\ \hline
\textbf{Past Continuous} & was/were + V-ing & Action en cours, cadre, description, action interrompue. \eng{The Fon \textbf{was wrapped} in his ndop cloth.} (Description). \eng{He \textbf{let} the wind \textbf{stir} the leaves.} \\ \hline
\end{tabularx}
\end{center}
\textbf{Règle d'or :} On utilise le \textbf{Past Perfect} pour le “flash-back” (ce qui s'est passé AVANT le moment principal du récit).
\end{propriete}

\begin{exoresolu}[Narrative Tenses Practice]
\textbf{Énoncé :} Put the verbs in brackets into the correct narrative tense (Past Simple, Past Perfect, or Past Continuous).
\begin{enumerate}
\item The families \_\_\_\_\_\_\_ (argue) for years before the Fon \_\_\_\_\_\_\_ (intervene).
\item While the elders \_\_\_\_\_\_\_ (listen), the children \_\_\_\_\_\_\_ (play) nearby.
\item Old Man Nfor \_\_\_\_\_\_\_ (wave) a document that he \_\_\_\_\_\_\_ (keep) for decades.
\item The river \_\_\_\_\_\_\_ (change) its course long before the dispute \_\_\_\_\_\_\_ (start).
\end{enumerate}
\tcblower
\textbf{Solution détaillée :}
\begin{enumerate}
\item \textbf{had been arguing} / \textbf{had argued} (Past Perfect — durée avant intervention) ; \textbf{intervened} (Past Simple — action principale).
\item \textbf{were listening} (Past Continuous — action en cours) ; \textbf{were playing} (Past Continuous — actions simultanées).
\item \textbf{waved} (Past Simple — action principale) ; \textbf{had kept} (Past Perfect — antériorité : il le gardait depuis longtemps avant de le montrer).
\item \textbf{had changed} (Past Perfect — antériorité par rapport au début du conflit) ; \textbf{started} (Past Simple — début du conflit, repère temporel).
\end{enumerate}
\end{exoresolu}

\courssub{Writing Workshop: Summary and Opinion}

\begin{methode}[How to write a summary — 4 Steps]
\begin{enumerate}
\item \textbf{Identify the main points:} Who? What? Where? When? How? Result? (Use your scanning notes).
\item \textbf{Select key information:} Delete details, examples, repetitions. Keep only the “skeleton”.
\item \textbf{Paraphrase:} Rewrite in your own words. Change vocabulary (\eng{dispute} $\to$ \eng{conflict}; \eng{settle} $\to$ \eng{resolve}). Change structure (active $\to$ passive, combine sentences).
\item \textbf{Connect:} Use linkers (\eng{First, Then, Finally, As a result, However}). Write in the \textbf{Present Simple} (standard for summaries) or \textbf{Past Simple} (if following the story timeline).
\end{enumerate}
\end{methode}

\begin{exemplebox}[Model Summary (Present Simple)]
\eng{
The text describes a land dispute in Bafut between the Nfor and Tuma families lasting two years. After the Quarter Head fails to mediate, the Fon summons both parties under the Great Mango Tree. Following hearings from both sides and the elders' questions, the Fon rules that the river serves as the boundary. Each family retains their improvements (mango trees, yam barns). The conflict ends peacefully through traditional mediation, highlighting the value of community harmony over legal battles.
}
\par\medskip
\textbf{Traduction :} Le texte décrit un conflit foncier à Bafut entre les familles Nfor et Tuma durant deux ans. Après l'échec de la médiation du chef de quartier, le Fon convoque les deux parties sous le Grand Manguier. Après les auditions et les questions des sages, le Fon décide que la rivière sert de frontière. Chaque famille conserve ses aménagements (manguiers, greniers à ignames). Le conflit se termine pacifiquement par la médiation traditionnelle, soulignant la valeur de l'harmonie communautaire par rapport aux batailles juridiques.
\end{exemplebox}

\begin{exercice}[Guided Writing: Summary and Opinion]
\begin{enumerate}
\item Write a summary of the text “Under the Mango Tree” in about 60–80 words. Use the Present Simple.
\item \textbf{Your opinion:} “Traditional justice is more effective than modern courts for community disputes.” Write a short paragraph (50–60 words) giving your opinion. Use \eng{I believe / In my opinion / For example / However}.
\end{enumerate}
\end{exercice}

\begin{corrige}[Exercice Guided Writing]
\begin{enumerate}
\item \textbf{Sample Summary:}
\eng{
In Bafut, a two-year land dispute between the Nfor and Tuma families escalates until the Fon intervenes. The Quarter Head's mediation having failed, the Fon hears both parties under the Great Mango Tree. He declares the river the natural boundary, allowing each family to keep their crops and buildings. The resolution restores peace through dialogue and tradition, avoiding the formal court system.
}
(73 words)
\item \textbf{Sample Opinion:}
\eng{
In my opinion, traditional justice is often more effective for community disputes. First, it is faster and cheaper than the Court of First Instance. Second, it preserves relationships because the goal is reconciliation, not punishment. For example, the sharing of kola nuts seals the peace. However, for serious crimes, modern courts are necessary to protect human rights.
}
(58 words)
\end{enumerate}
\end{corrige}

\courssub{Speaking Task: Mediation Role-Play}

\begin{experience}[Mediation Simulation — The Palaver Tree]
\textbf{Scenario:} A dispute between two neighbours in a quarter of Douala.
\begin{itemize}
\item \textbf{Neighbour A:} Accuses Neighbour B's goats of destroying his vegetable garden.
\item \textbf{Neighbour B:} Claims the garden encroaches on the path where goats always pass.
\item \textbf{Quarter Head:} Mediator. Must listen, ask questions, propose a compromise.
\item \textbf{Two Elders:} Advisors. Ask for evidence (photos, witnesses, boundary markers).
\end{itemize}
\par\medskip
\textbf{Instructions:}
\begin{enumerate}
\item Form groups of 5. Assign roles.
\item Prepare arguments for 5 minutes (use vocabulary: \eng{boundary, damage, compensation, fence, right of way, agreement}).
\item Role-play the meeting (5–7 minutes). The Quarter Head must announce a final decision.
\item Debrief: Did the mediation work? What language was useful?
\end{enumerate}
\textbf{Useful phrases:}
\eng{"I would like to explain my side..."} / \eng{"The issue is..."} / \eng{"Have you any proof?"} / \eng{"I propose that..."} / \eng{"We agree to..."} / \eng{"Let us shake hands."}
\end{experience}

\begin{aretenir}[Objectifs de cette leçon — Warm-up à Production]
\begin{itemize}
\item Identifier les types de conflits communautaires (\eng{land, inheritance, noise, water, boundaries}).
\item Connaître les acteurs de la résolution : \eng{Quarter Head, Chief, Elders, Court, Police}.
\item Réactiver le vocabulaire clé : \eng{dispute, settle, agreement, peace, chief, listen}.
\item Appliquer la \textbf{méthode de prédiction} : titre + images + premières phrases.
\item Choisir la bonne stratégie : \textbf{Skimming} (idée générale) vs \textbf{Scanning} (détail précis).
\item Comprendre un texte narratif authentique (\eng{Under the Mango Tree}) et son lexique culturel (\eng{ndop, placenta, kola nut, harmattan, Fon}).
\item Maîtriser les \textbf{temps du récit} : Past Simple, Past Perfect, Past Continuous.
\item Rédiger un \textbf{résumé} (paraphrase, connecteurs, temps présent) et exprimer un \textbf{avis personnel} argumenté.
\item Participer à un \textbf{jeu de rôle de médiation} en réinvestissant le lexique du thème.
\end{itemize}
\end{aretenir}


\coursec{Reading text with glossary}

Dans la section précédente (\emph{Warm-up and pre-reading}), nous avons activé nos connaissances sur les conflits de voisinage et prédit le contenu du texte à partir de son titre. Le moment est maintenant venu de lire le texte lui-même. Nous allons procéder en \textbf{deux lectures} : une première lecture globale et rapide (\emph{skimming}), sans dictionnaire, puis une deuxième lecture détaillée (\emph{scanning}) à l'aide du glossaire.

\courssub{First reading: skimming (lecture globale)}

\begin{methode}[Comment faire une lecture globale (skimming)]
\begin{enumerate}
\item Lis le texte \textbf{en silence}, du début à la fin, sans t'arrêter.
\item \textbf{N'ouvre pas le dictionnaire} et ne traduis rien mot à mot.
\item Repère visuellement les mots que tu connais déjà (noms propres, verbes de base, connecteurs) : ils suffisent souvent à comprendre l'idée générale.
\item À la fin, pose-toi trois questions : \emph{Who?} (qui ?), \emph{Where?} (où ?), \emph{What happens?} (que se passe-t-il ?).
\end{enumerate}
\end{methode}

\begin{attention}[Piège n° 1 : la traduction mot à mot]
L'erreur la plus fréquente des élèves francophones est de vouloir \textbf{traduire chaque mot} pendant la lecture. Résultat : on perd le fil, on oublie le début de la phrase et on abandonne. En anglais comme en français, on lit d'abord pour comprendre le \textbf{sens global}. Les mots inconnus seront traités ensuite, au cours de la deuxième lecture. Un bon lecteur accepte de ne pas tout comprendre du premier coup.
\end{attention}

\textbf{Consigne de première lecture.} Lis le texte ci-dessous en deux minutes maximum. Ne t'arrête sur aucun mot. Ensuite, réponds oralement : \emph{Who are the characters? Where does the story take place? What is the problem?}

\begin{textebox}[Peace Under the Mango Tree]
Last year, in the village of Bafut, two families quarrelled over a piece of land behind the market. The Mbuh family said their grandfather had planted coffee trees there fifty years ago. The Ndi family answered that the land belonged to their uncle. For months, nobody could farm the land, and the two families stopped greeting each other.

One Saturday morning, the chief invited both families to gather under the big mango tree near the palace. The elders listened to each side carefully. Then they asked questions and examined old documents. After three hours, the chief gave his decision: the land would be divided into two equal parts, and each family would pay a small fine for the months of trouble. Both families shook hands, shared kola nuts, forgave each other, and promised to live in peace.

Today, many towns in Cameroon also have mediation centres where trained mediators help people settle disputes without going to court. Going to court is slow and expensive; dialogue is faster and keeps friendship alive. As the elders say: “When two brothers fight, the stranger takes their land.”
\end{textebox}

\textbf{Vérification rapide après le skimming.} Même sans comprendre tous les mots, tu as dû repérer : deux familles (the Mbuh family, the Ndi family), un village (Bafut), un problème de terrain (\emph{land}), un chef (\emph{chief}), des anciens (\emph{elders}), une fin heureuse (\emph{shook hands}, \emph{forgave each other}, \emph{peace}) et une leçon sur la médiation. C'est suffisant pour une première lecture.

\courssub{Second reading: scanning with the glossary (lecture détaillée)}

La deuxième lecture est plus lente. Cette fois, tu peux t'arrêter sur les mots difficiles et utiliser le glossaire ci-dessous. Les mots du glossaire apparaissent dans l'ordre du texte.

\begin{definition}[Glossaire du texte]
\begin{center}
\begin{tabularx}{\linewidth}{|l|l|X|}
\hline
\rowcolor{popblueL} \textbf{English word} & \textbf{Français} & \textbf{Exemple tiré ou adapté du texte} \\
\hline
quarrel (v./n.) & se quereller / querelle & \eng{Two families quarrelled over land.} « Deux familles se sont querellées pour un terrain. » \\
\hline
elder (n.) & ancien, aîné (respecté) & \eng{The elders listened to each side.} « Les anciens ont écouté chaque partie. » \\
\hline
gather (v.) & se rassembler, réunir & \eng{Both families gathered under the mango tree.} « Les deux familles se sont rassemblées sous le manguier. » \\
\hline
examine (v.) & examiner, étudier & \eng{They examined old documents.} « Ils ont examiné de vieux documents. » \\
\hline
document (n.) & document, papier & \eng{Old documents proved the ownership.} « De vieux documents prouvaient la propriété. » \\
\hline
decision (n.) & décision & \eng{The chief gave his decision.} « Le chef a donné sa décision. » \\
\hline
fine (n.) & amende & \eng{Each family paid a small fine.} « Chaque famille a payé une petite amende. » \\
\hline
shake hands (expr.) & se serrer la main & \eng{Both families shook hands.} « Les deux familles se sont serré la main. » \\
\hline
forgive (v.) & pardonner & \eng{They forgave each other.} « Ils se sont pardonné. » \\
\hline
mediation (n.) & médiation & \eng{Many towns have mediation centres.} « De nombreuses villes ont des centres de médiation. » \\
\hline
settle (v.) & régler, résoudre & \eng{Mediators help people settle disputes.} « Les médiateurs aident à régler les conflits. » \\
\hline
dispute (n.) & conflit, différend & \eng{... settle disputes without going to court.} « ... régler les conflits sans aller au tribunal. » \\
\hline
stranger (n.) & étranger, inconnu & \eng{The stranger takes their land.} « L'étranger prend leur terre. » \\
\hline
\end{tabularx}
\end{center}
\end{definition}

\begin{attention}[Faux amis et pièges lexicaux]
\begin{itemize}
\item \eng{fine} ne signifie pas « fin » ni « bien » ici : c'est une \textbf{amende} (somme d'argent). \eng{She paid a fine.} « Elle a payé une amende. »
\item \eng{elder} n'est pas « électeur » ni simplement « vieux » : c'est un \textbf{ancien}, une personne âgée respectée pour sa sagesse.
\item \eng{to settle} ne veut pas dire « s'asseoir » (\eng{sit down}) : \eng{to settle a dispute} = « régler un conflit ». (Mais \eng{to settle in a village} = « s'installer dans un village ».)
\item \eng{court} n'est pas « court » (petit) : \eng{to go to court} = « aller au tribunal ».
\item \eng{stranger} ne signifie pas « danger » : c'un \textbf{étranger}, quelqu'un qui n'est pas de la communauté.
\item \eng{to shake hands} : ne dis pas \emph{to tighten the hand} (littéral incorrect). L'expression figée est \eng{shake hands} (prétérit \eng{shook hands}).
\end{itemize}
\end{attention}

\begin{definition}[Mediation]
\eng{Mediation is a way of settling a dispute in which a neutral person (the mediator) helps the two sides reach an agreement.} — La \textbf{médiation} est une manière de régler un conflit dans laquelle une personne neutre (le \emph{mediator}, médiateur) aide les deux parties à trouver un accord. Dans le texte, le chef et les anciens jouent le rôle de médiateurs traditionnels ; dans les villes modernes, ce sont les \emph{mediation centres} (centres de médiation) avec des \emph{trained mediators} (médiateurs formés).
\end{definition}

\courssub{Understanding the structure of the text}

Un bon lecteur repère aussi \textbf{comment le texte est organisé}. Notre récit suit une progression narrative très classique : le problème, la rencontre, la solution, la leçon.

\begin{popfigure}
\begin{tikzpicture}[
box/.style={draw=popblue, fill=popblueL, rounded corners, thick, align=center, text width=3.1cm, minimum height=1.5cm, font=\small},
arr/.style={-{Stealth[length=3mm]}, very thick, poporange}]
\node[box, align=center] (p1){\textbf{Paragraph 1}\\ The conflict\\ \emph{two families, one piece of land}};
\node[box, right=0.9cm of p1, align=center] (p2){\textbf{Paragraph 2}\\ The meeting\\ \emph{under the mango tree}};
\node[box, right=0.9cm of p2, align=center] (p3){\textbf{Paragraph 3}\\ The solution\\ \emph{land divided, fine paid, peace}};
\node[box, right=0.9cm of p3, align=center] (p4){\textbf{Paragraph 4}\\ Lesson learnt\\ \emph{dialogue better than court}};
\draw[arr] (p1) -- (p2);
\draw[arr] (p2) -- (p3);
\draw[arr] (p3) -- (p4);
\end{tikzpicture}
\legende{Structure du texte : du conflit à la leçon de sagesse.}
\end{popfigure}

\begin{aretenir}[La structure d'un récit de résolution de conflit]
La plupart des textes narratifs anglais sur ce thème suivent le schéma : \cle{conflict} (le conflit) $\rightarrow$ \cle{meeting} (la rencontre / la médiation) $\rightarrow$ \cle{solution} (la solution) $\rightarrow$ \cle{lesson} (la leçon ou morale). Repérer cette structure t'aide à répondre aux questions de compréhension et à résumer le texte.
\end{aretenir}

\courssub{Key sentences: observation and translation}

Relisons deux phrases importantes du texte et observons leur construction.

\begin{exemplebox}[Phrases clés traduites]
\begin{itemize}
\item \eng{They settled the dispute peacefully.} « Ils ont réglé le conflit pacifiquement. »
\item \eng{The elders listened to each side.} « Les anciens ont écouté chaque partie. »
\item \eng{Both families shook hands, forgave each other and promised to live in peace.} « Les deux familles se sont serré la main, se sont pardonné et ont promis de vivre en paix. »
\item \eng{Dialogue is faster and keeps friendship alive.} « Le dialogue est plus rapide et maintient l'amitié en vie. »
\end{itemize}
\end{exemplebox}

\begin{propriete}[Deux points de langue observés dans le texte]
\begin{enumerate}
\item \textbf{Le prétérit simple (Past Simple)} : le récit est au passé (\emph{quarrelled, said, answered, invited, gathered, listened, examined, gave, divided, paid, shook, shared, forgave, promised}). 
      \begin{itemize}
      \item Verbes réguliers : base + \textbf{-ed} (\emph{invite} $\rightarrow$ \emph{invited}, \emph{examine} $\rightarrow$ \emph{examined}). Attention à l'orthographe : \emph{quarrel} $\rightarrow$ \emph{quarrelled} (double \emph{l} en anglais britannique), \emph{stop} $\rightarrow$ \emph{stopped} (double \emph{p}).
      \item Verbes irréguliers fréquents : \emph{say} $\rightarrow$ \emph{said} (/sed/), \emph{give} $\rightarrow$ \emph{gave}, \emph{shake} $\rightarrow$ \emph{shook}, \emph{forgive} $\rightarrow$ \emph{forgave}, \emph{pay} $\rightarrow$ \emph{paid}.
      \end{itemize}
\item \textbf{Le verbe \eng{to listen} se construit avec la préposition \eng{to}} : \eng{to listen to somebody / something}. On ne dit jamais \emph{listen somebody}. \eng{The elders listened to each side.} « Les anciens ont écouté chaque partie. »
\end{enumerate}
\end{propriete}

\begin{attention}[Erreurs fréquentes des francophones]
\begin{itemize}
\item Ne dis pas \emph{the peoples} pour « les gens » : \eng{people} est déjà pluriel (singulier \eng{person}). \eng{Mediators help people.} « Les médiateurs aident les gens. »
\item « Pardonner à quelqu'un » se dit \eng{to forgive somebody}, \textbf{sans préposition} : \eng{She forgave her neighbour.} (Pas \emph{forgave to her neighbour}).
\item \eng{Agreement} (accord) vs \eng{Disagreement} (désaccord). Ne confonds pas avec \eng{Dispute} (conflit).
\end{itemize}
\end{attention}

\courssub{Quick check: from text to meaning}

\begin{exoresolu}[Retrouver les informations dans le texte (scanning)]
Énoncé. Réponds en anglais, par une phrase courte, en t'appuyant sur le texte.
\begin{enumerate}
\item Where did the two families quarrel?
\item Who listened to each side?
\item What was the chief's decision?
\item Why is going to court not the best solution, according to the text?
\end{enumerate}
\tcblower
Solution détaillée.
\begin{enumerate}
\item \eng{They quarrelled in the village of Bafut, over a piece of land behind the market.} (On repère \emph{Bafut} et \emph{land} au paragraphe 1.)
\item \eng{The elders listened to each side.} (Paragraphe 2 : le verbe \emph{listened} et le sujet \emph{the elders}.)
\item \eng{The chief decided that the land would be divided into two equal parts and that each family would pay a small fine.} (Paragraphe 2, après \emph{the chief gave his decision}.)
\item \eng{Going to court is slow and expensive, while dialogue is faster and keeps friendship alive.} (Paragraphe 3 : opposition \emph{slow and expensive} vs \emph{faster and keeps friendship alive}.)
\end{enumerate}
Méthode : pour chaque question, repère un \textbf{mot-clé} de la question (\emph{where}, \emph{who}, \emph{decision}, \emph{court}), localise-le dans le texte, puis lis la phrase entière autour de ce mot.
\end{exoresolu}

\begin{exercice}[À toi de jouer]
Relis le texte une dernière fois, puis complète ces phrases avec un mot ou une expression du glossaire (forme correcte) : \emph{quarrel, elder, gather, examine, document, decision, fine, shake hands, forgive, mediation, settle, dispute, stranger}.
\begin{enumerate}
\item The chief asked both families to \ldots\ under the mango tree.
\item The \ldots\ examined old documents before giving their opinion.
\item Each family had to pay a small \ldots\ for the months of trouble.
\item The two families decided to \ldots\ each other and live in peace.
\item A mediator helps people \ldots\ their \ldots\ without going to court.
\item The proverb says: when two brothers fight, the \ldots\ takes their land.
\end{enumerate}
\end{exercice}

\begin{corrige}[Exercice : compléter avec le glossaire]
\begin{enumerate}
\item \eng{gather} — « se rassembler » : \eng{The chief asked both families to gather under the mango tree.}
\item \eng{elders} — « les anciens » : \eng{The elders examined old documents.}
\item \eng{fine} — « amende » : \eng{Each family had to pay a small fine.}
\item \eng{forgive} — « pardonner » : \eng{The two families decided to forgive each other.}
\item \eng{settle} ; \eng{disputes} (pluriel dans le texte) — \eng{A mediator helps people settle their disputes without going to court.}
\item \eng{stranger} — « l'étranger » : \eng{The proverb says: when two brothers fight, the stranger takes their land.}
\end{enumerate}
\end{corrige}

\begin{savaistu}[La sagesse du manguier]
Dans de nombreuses communautés du Nord-Ouest du Cameroun, comme à Bafut, les conflits se règlent traditionnellement sous un grand arbre, souvent un manguier (\emph{mango tree}), devant le palais du chef (\emph{fondom}). Le proverbe final du texte — \eng{“When two brothers fight, the stranger takes their land.”} « Quand deux frères se battent, l'étranger prend leur terre. » — rappelle que la division affaiblit la communauté face aux menaces extérieures. C'est l'équivalent de notre « L'union fait la force », mais formulé par la menace de la perte. La médiation traditionnelle valorise le \emph{consensus} et la \emph{réconciliation} (\emph{forgiveness}) plutôt que la victoire de l'un sur l'autre.
\end{savaistu}

Tu as maintenant lu le texte deux fois, compris sa structure et son vocabulaire. Dans la prochaine section, nous approfondirons les stratégies de lecture (\emph{skimming}, \emph{scanning}, deviner le sens des mots) pour devenir un lecteur autonome.

\coursec{Reading strategies: skimming, scanning, guessing}

Après avoir découvert le texte support et son glossaire (section précédente), nous allons maintenant acquérir trois stratégies indispensables pour lire efficacement en anglais, surtout en situation d'examen où le temps est limité. Ces stratégies — \cle{skimming}, \cle{scanning} et \cle{guessing meaning from context} — vous permettent de comprendre l'essentiel sans connaître chaque mot.

Pour illustrer chaque stratégie, nous utiliserons le texte court ci-dessous, typique d'un sujet de compréhension sur le règlement des conflits communautaires.

\begin{textebox}[The Mango Tree Meeting]
Last Tuesday, a serious dispute arose between the Mbarga and the Nkolo families in Mvomeka'a village. The conflict started over a piece of fertile land near the river. Both families claimed ownership based on ancestral rights. For weeks, they quarrelled bitterly and stopped greeting each other at the market. The tension was so high that young men from both sides almost fought with machetes.

Fortunately, the village chief, His Majesty Fozo, decided to intervene. He summoned the heads of both families to a meeting under the big mango tree in the village square. The meeting lasted three hours. The chief listened patiently to both parties. He reminded them of the old proverb: “A family that eats together stays together.” He proposed a compromise: the land would be shared equally, and the border marked by a row of palm trees.

To everyone's relief, both families accepted the verdict. They shook hands publicly and agreed to organise a feast the following Saturday to seal the reconciliation. The chief fined each family a goat for the disturbance of peace. Peace returned to Mvomeka'a.
\end{textebox}

\begin{definition}[Glossaire clé du texte]
\begin{itemize}
    \item \textbf{arose} (v., past of arise) : a surgi, a éclaté
    \item \textbf{fertile} (adj.) : fertile
    \item \textbf{ownership} (n.) : propriété, droit de possession
    \item \textbf{ancestral rights} (n. pl.) : droits ancestraux
    \item \textbf{quarrelled} (v.) : se sont disputés
    \item \textbf{bitterly} (adv.) : amèrement, violemment
    \item \textbf{summoned} (v.) : a convoqué
    \item \textbf{verdict} (n.) : verdict, décision
    \item \textbf{seal} (v.) : sceller, officialiser
    \item \textbf{fined} (v., past of fine) : a infligé une amende à
\end{itemize}
\end{definition}

\courssub{Skimming: reading for the main idea}

Le \cle{skimming} (lecture rapide diagonale) consiste à parcourir un texte très vite pour en saisir l'\textbf{idée générale} (the gist), sans lire chaque mot. C'est la première chose à faire face à un texte inconnu.

\begin{methode}[Comment faire du skimming en 30 secondes]
\begin{enumerate}
    \item \textbf{Lisez le titre} (s'il y en a un) : il annonce le sujet.
    \item \textbf{Lisez la première phrase} de chaque paragraphe (souvent la \emph{topic sentence}).
    \item \textbf{Lisez la dernière phrase} du dernier paragraphe (souvent la conclusion).
    \item \textbf{Repérez les mots-clés} répétés (noms propres, noms communs forts, verbes d'action).
    \item \textbf{Formulez l'idée générale} en une seule phrase simple en anglais.
\end{enumerate}
\end{methode}

\begin{exemplebox}[Application au texte « The Mango Tree Meeting »]
\begin{itemize}
    \item \textbf{Titre} : \eng{The Mango Tree Meeting} → Une rencontre sous un manguier.
    \item \textbf{Premières phrases} : \eng{Last Tuesday, a serious dispute arose...} (conflit) ; \eng{Fortunately, the village chief... decided to intervene} (intervention du chef) ; \eng{To everyone's relief, both families accepted the verdict} (résolution).
    \item \textbf{Mots-clés} : \cle{dispute}, \cle{chief}, \cle{meeting}, \cle{accepted}, \cle{peace}.
    \item \textbf{Idée générale} : \eng{A land dispute between two families was settled by the village chief under a mango tree.}
\end{itemize}
\end{exemplebox}

\begin{attention}[Piège fréquent]
Ne confondez pas \textbf{skimming} et \textbf{lecture détaillée}. Au skimming, vous \textbf{ne lisez pas} les adjectifs, les exemples, les chiffres précis. Si vous mettez plus d'une minute sur un texte de 200 mots, vous lisez trop lentement.

\begin{center}
\begin{tabularx}{\linewidth}{|X|X|X|}
\hline
\rowcolor{popblueL} \textbf{Action} & \textbf{Skimming (Gist)} & \textbf{Reading for detail} \\
\hline
Vitesse & Très rapide (30--60 sec) & Lente (5--10 min) \\
\hline
Objectif & Idée principale & Compréhension totale \\
\hline
Focus & Titre, 1ère/dernière phrases, mots-clés & Chaque phrase, grammaire, nuances \\
\hline
\end{tabularx}
\end{center}
\end{attention}

\courssub{Scanning: locating specific information}

Le \cle{scanning} (lecture de repérage) consiste à chercher une \textbf{information précise} (un nom, une date, un chiffre, un lieu, une raison) sans relire tout le texte. Vos yeux « scannent » la page comme un scanner de supermarché.

\begin{methode}[Comment faire du scanning efficacement]
\begin{enumerate}
    \item \textbf{Analysez la question} : soulignez le \cle{mot-clé} ou la \cle{question word} (Who, Where, When, How much, Why...).
    \item \textbf{Cherchez le mot-clé} (ou un synonyme) dans le texte en faisant courir votre doigt ou votre regard verticalement.
    \item \textbf{Lisez seulement la phrase} où se trouve le mot-clé.
    \item \textbf{Extrayez la réponse} courte et précise.
\end{enumerate}
\end{methode}

\begin{exemplebox}[Questions de scanning sur le texte]
\begin{tabular}{|l|l|l|}
\hline
\textbf{Question} & \textbf{Mot-clé cherché} & \textbf{Réponse trouvée} \\
\hline
Where did the meeting take place? & \cle{Where} / \cle{meeting} / \cle{place} & \eng{Under the big mango tree in the village square.} \\
\hline
How long did the meeting last? & \cle{How long} / \cle{last} & \eng{Three hours.} \\
\hline
What did the chief propose? & \cle{propose} / \cle{chief} & \eng{The land would be shared equally...} \\
\hline
What was the fine? & \cle{fine} / \cle{fined} & \eng{A goat for each family.} \\
\hline
\end{tabular}
\end{exemplebox}

\begin{aretenir}[Astuce de pro]
Les \textbf{noms propres} (Majuscules : \eng{Mvomeka'a}, \eng{Fozo}, \eng{Tuesday}), les \textbf{chiffres} (\eng{three hours}, \eng{Last Tuesday}) et la \textbf{ponctuation} (deux-points, guillemets) sont des « phares » visuels pour le scanning. Repérez-les d'un coup d'œil.

\begin{center}
\begin{tabularx}{\linewidth}{|X|X|}
\hline
\rowcolor{poporangeL} \textbf{Question Word} & \textbf{Information cherchée (Scan target)} \\
\hline
Who / Whose & Personne (nom propre, pronom, rôle) \\
\hline
Where & Lieu (préposition + nom : \eng{in}, \eng{at}, \eng{under}) \\
\hline
When / What time & Moment (date, heure, durée, jour) \\
\hline
How much / How many & Quantité / Prix (chiffre + unité) \\
\hline
Why / What for & Raison / But (\eng{because}, \eng{to}, \eng{for}) \\
\hline
How & Manière / Moyen (\eng{by} + V-ing, adverbe en \eng{-ly}) \\
\hline
\end{tabularx}
\end{center}
\end{aretenir}

\courssub{Guessing meaning from context}

Au Bac et dans la vie réelle, \textbf{vous n'avez pas de dictionnaire}. La compétence vitale est de \cle{déduire le sens} d'un mot inconnu grâce aux indices textuels. C'est ce qu'on appelle \eng{inferring meaning} ou \eng{contextual guessing}.

\begin{propriete}[Règle de déduction lexicale : les 3 indices]
Le sens d'un mot inconnu se devine en croisant trois sources :
\begin{enumerate}
    \item \textbf{Le contexte situationnel} : le thème général, ce qui se passe avant/après (cause/effet, contraste).
    \item \textbf{La classe grammaticale} (part of speech) : est-ce un \cle{nom} (sujet/objet), un \cle{verbe} (action/état), un \cle{adjectif} (qualité), un \cle{adverbe} (manière) ? Les suffixes aident (\eng{-tion} = nom, \eng{-ly} = adv, \eng{-ed/-ing} = verbe/adj).
    \item \textbf{Les indices lexicaux immédiats} : synonymes (\eng{or}, \eng{that is to say}), antonymes (\eng{but}, \eng{however}, \eng{unlike}), définitions (\eng{is a kind of}, \eng{means}), exemples (\eng{such as}, \eng{for example}), explications entre parenthèses.
\end{enumerate}
\end{propriete}

\begin{popfigure}
\begin{tikzpicture}[
    node distance=1.2cm and 1.5cm,
    decision/.style={diamond, draw=popdark, thick, fill=poporangeL, text width=2.8cm, align=center, inner sep=4pt},
    process/.style={rectangle, rounded corners, draw=popblue, thick, fill=popblueL, text width=2.8cm, align=center, inner sep=4pt},
    start/.style={ellipse, draw=popgreen, thick, fill=popgreenL, text width=2.5cm, align=center, inner sep=4pt},
    arrow/.style={-Stealth, thick, popdark}
]
\node[start] (start) {Unknown word?};
\node[process, below of=start, align=center] (step1){Look at the words\\ before and after};
\node[decision, below of=step1, align=center] (step2){Is it a Noun,\\ Verb, Adj or Adv?};
\node[process, below left=1.5cm and 0.5cm of step2, align=center] (step3n){Noun: Who/What?\\ (person, thing, place)};
\node[process, below right=1.5cm and 0.5cm of step2, align=center] (step3v){Verb/Adj/Adv: Action,\\ State, Quality, Manner?};
\node[process, below of=step2, yshift=-3.2cm, align=center] (step4){Guess the meaning\\ from context clues};
\node[process, below of=step4, align=center] (step5){Check: Does it make\\ sense in the sentence?};
\node[start, below of=step5] (end) {Meaning found!};

\draw[arrow] (start) -- (step1);
\draw[arrow] (step1) -- (step2);
\draw[arrow] (step2) -| node[near start, above, font=\small, text=popdark] {Noun} (step3n);
\draw[arrow] (step2) -| node[near start, above, font=\small, text=popdark] {Verb/Adj/Adv} (step3v);
\draw[arrow] (step3n) |- (step4);
\draw[arrow] (step3v) |- (step4);
\draw[arrow] (step4) -- (step5);
\draw[arrow] (step5) -- (end);
\end{tikzpicture}
\legende{Organigramme : Stratégie de déduction du sens d'un mot inconnu}
\end{popfigure}

\begin{exemplebox}[Entraînement guidé sur le texte]
Appliquons la méthode à quatre mots du texte que vous ne connaissez peut-être pas.

\begin{enumerate}
    \item \textbf{quarrelled} (l. 4) : \eng{...they \cle{quarrelled} bitterly and stopped greeting each other...}
    \begin{itemize}
        \item \textit{Classe} : Verbe (passé en \eng{-ed}, sujet \eng{they}).
        \item \textit{Contexte} : Conflit pour une terre, tension, « stopped greeting » (arrêté de se saluer).
        \item \textit{Indice} : \eng{bitterly} (amèrement) + conséquence négative.
        \item \textbf{Déduction} : Se disputer violemment. $\rightarrow$ \eng{to quarrel} = se disputer.
    \end{itemize}

    \item \textbf{fined} (l. 15) : \eng{The chief \cle{fined} each family a goat...}
    \begin{itemize}
        \item \textit{Classe} : Verbe (passé, sujet \eng{chief}, objet \eng{each family}).
        \item \textit{Contexte} : Après la réconciliation, le chef sanctionne la perturbation.
        \item \textit{Indice} : Structure \eng{fine someone something} (verbe ditransitif). Un chèvre = paiement.
        \item \textbf{Déduction} : Infliger une amende. $\rightarrow$ \eng{to fine} = infliger une amende à qqn (attention : \eng{fine} adj. = bien/beau, \eng{fine} n. = amende).
    \end{itemize}

    \item \textbf{gather} (implicite : \eng{summoned} $\approx$ gather) : \eng{He \cle{summoned} the heads... to a meeting}. Si on ne connaît pas \eng{summoned} :
    \begin{itemize}
        \item \textit{Classe} : Verbe passé.
        \item \textit{Contexte} : Chef + chefs de famille + réunion.
        \item \textbf{Déduction} : Convoquer, rassembler. $\rightarrow$ \eng{to summon} = convoquer / \eng{to gather} = rassembler.
    \end{itemize}

    \item \textbf{seal} (l. 16) : \eng{...agreed to organise a feast to \cle{seal} the reconciliation.}
    \begin{itemize}
        \item \textit{Mot cible} : \eng{seal} (verbe).
        \item \textit{Contexte} : Fin du conflit, poignée de main, fête.
        \item \textit{Indice} : \eng{reconciliation} (réconciliation). Sceller une réconciliation = la rendre officielle, définitive.
        \item \textbf{Déduction} : Sceller, officialiser, conclure.
    \end{itemize}
\end{enumerate}
\end{exemplebox}

\begin{exoresolu}[Entraînement : Devinez le sens des mots soulignés]
\textbf{Texte} : \eng{After the verdict, the atmosphere was \textbf{tense} at first. But the feast \textbf{alleviated} the tension. Elders \textbf{urged} the youth to forget the past.}

\begin{enumerate}
    \item \textbf{tense} (l. 1) : \textit{Adjectif} (après \eng{was}). Contexte : « at first » (au début), mais fête après. Contraste. $\rightarrow$ \textbf{Tendu, crispé}. (Antonyme implicite de détendu après la fête).
    \item \textbf{alleviated} (l. 2) : \textit{Verbe passé} (sujet \eng{feast}). Objet \eng{tension}. La fête a fait quoi à la tension ? $\rightarrow$ \textbf{Allégé, apaisé, réduit}. (Synonyme contextuel : \eng{reduced}).
    \item \textbf{urged} (l. 3) : \textit{Verbe passé} (sujet \eng{Elders}, objet \eng{youth}, infinitif \eng{to forget}). Les anciens ont fait quoi aux jeunes ? $\rightarrow$ \textbf{Exhorté, pressé, fortement conseillé}. (Plus fort que \eng{advised}).
\end{enumerate}
\tcblower
\textbf{Solution détaillée} :
\begin{enumerate}
    \item \textbf{tense} = \emph{tendu / crispé}. Indice grammatical : adjectif après \eng{be}. Indice contextuel : opposition \eng{at first... But}.
    \item \textbf{alleviated} = \emph{allégé / apaisé}. Indice grammatical : verbe transitif. Indice logique : une fête après un conflit réduit la tension.
    \item \textbf{urged} = \emph{exhorté / pressé de}. Indice grammatical : \eng{urge someone to do something}. Indice sémantique : les anciens ont l'autorité morale pour pousser les jeunes à pardonner.
\end{enumerate}
\end{exoresolu}

\begin{attention}[Erreur fatale des francophones : le réflexe dictionnaire]
\begin{itemize}
    \item \textbf{Erreur} : S'arrêter à chaque mot inconnu, sortir le téléphone/dictionnaire, perdre 10 minutes, paniquer.
    \item \textbf{Réalité du Bac} : 0 dictionnaire autorisé. Le texte est conçu pour que les mots difficiles soient \textbf{devinables}.
    \item \textbf{Bonne habitude} : Entraînez-vous \textbf{maintenant} à cocher les mots inconnus, à deviner, à continuer la lecture. Vérifiez \textbf{après} seulement.
    \item \textbf{Faux ami} : \eng{fine} (adj.) = beau/bien ; \eng{fine} (n.) = amende ; \eng{to fine} = verbaliser. \eng{Actually} = en fait (pas « actuellement » $\rightarrow$ \eng{currently}).
\end{itemize}
\end{attention}

\begin{center}
\begin{tabularx}{\linewidth}{|X|X|X|}
\hline
\rowcolor{poppurpleL} \textbf{Stratégie} & \textbf{Question clé} & \textbf{Action concrète} \\
\hline
\textbf{Skimming} & What is the text \textbf{mainly} about? & Lire titre + 1ère phrase § + dernière phrase. Écrire 1 phrase résumé. \\
\hline
\textbf{Scanning} & \textbf{Where / When / Who / How much}...? & Souligner mot-clé question $\rightarrow$ Chasser le mot dans le texte $\rightarrow$ Lire la phrase cible. \\
\hline
\textbf{Guessing} & What does \textbf{X} mean? & 1. Classe du mot ? 2. Contexte (cause/effet/contraste) ? 3. Synonyme/indice dans la phrase ? 4. Vérifier si ça tient. \\
\hline
\end{tabularx}
\end{center}

\courssub{Practice: Integrated strategy workout}

Appliquez les trois stratégies séquentiellement sur le texte \eng{The Mango Tree Meeting} (relisez-le 30 secondes si besoin).

\begin{exercice}[Exercices d'application -- À faire sur votre cahier]
\begin{enumerate}
    \item \textbf{Skimming} : Écrivez l'idée générale du texte en \textbf{une seule phrase anglaise} (max 20 mots).
    \item \textbf{Scanning} : Répondez brièvement :
    \begin{itemize}
        \item[a)] What proverb did the chief quote?
        \item[b)] When is the feast planned?
        \item[c)] What marks the new border?
    \end{itemize}
    \item \textbf{Guessing} : Déduisez le sens de : \eng{ancestral} (l. 3), \eng{verdict} (l. 13), \eng{feast} (l. 15). Justifiez par l'indice textuel précis.
    \item \textbf{Traduction ciblée} : Traduisez en français : \eng{Both families shook hands publicly and agreed to organise a feast to seal the reconciliation.}
    \item \textbf{Expression d'opinion} (Speaking/Writing) : En 3--4 phrases anglaises, donnez votre avis : \eng{Do you think traditional mediation is more effective than modern courts for family disputes? Why?} Utilisez \eng{I believe / In my opinion / Because / For example}.
\end{enumerate}
\end{exercice}

\begin{corrige}[Exercices d'application]
\begin{enumerate}
    \item \eng{A village chief settled a land dispute between two families under a mango tree, restoring peace.}
    \item \textbf{a)} \eng{“A family that eats together stays together.”} \quad \textbf{b)} \eng{The following Saturday.} \quad \textbf{c)} \eng{A row of palm trees.}
    \item \textbf{ancestral} : adjectif (devant \eng{rights}). Contexte : « based on ancestral rights » = droits transmis par les ancêtres $\rightarrow$ \emph{ancestral / héréditaire}. \\
    \textbf{verdict} : nom (sujet \eng{accepted the verdict}). Contexte : décision du chef après écoute $\rightarrow$ \emph{verdict / décision / jugement}. \\
    \textbf{feast} : nom (objet de \eng{organise a feast}). Contexte : « to seal the reconciliation », après poignée de main $\rightarrow$ \emph{fête / banquet / festin}.
    \item \eng{Les deux familles se sont serré la main publiquement et ont accepté d'organiser une fête pour sceller la réconciliation.}
    \begin{itemize}
        \item \eng{shook hands} = se serrer la main (idiome, \textbf{pas} « secoué les mains »).
        \item \eng{publicly} = publiquement.
        \item \eng{seal} = sceller / officialiser.
    \end{itemize}
    \item \textbf{Exemple de réponse} : \eng{I believe traditional mediation is more effective for family disputes because it focuses on reconciliation, not punishment. For example, in the text, the feast and the handshake restored peace permanently. Modern courts often create winners and losers, which can deepen hatred.}
\end{enumerate}
\end{corrige}

\begin{experience}[Débat guidé : Justice traditionnelle vs justice moderne]
En groupes de 4, préparez un débat de 5 minutes. Deux élèves défendent la médiation traditionnelle (palabre, chef, réconciliation), deux défendent la justice formelle (tribunal, avocat, loi écrite). Utilisez le vocabulaire : \eng{mediation, reconciliation, verdict, fine, evidence, lawyer, judge, compromise, binding, enforceable}. Le professeur note les arguments au tableau en deux colonnes. Vote final de la classe.
\end{experience}

\begin{savaistu}[La palabre africaine et la justice restaurative]
Le texte illustre la \cle{palabre} (de l'ancien portugais \eng{palavra} = parole), institution majeure de résolution des conflits en Afrique centrale et occidentale. Contrairement à la justice occidentale (punitive, gagnant/perdant), la justice traditionnelle vise la \cle{réconciliation} et la \cle{réparation du lien social}. Le chef n'est pas un juge qui tranche sur des preuves, mais un médiateur qui cherche le consensus. La \eng{feast} (fête) et le \eng{shake hands} (poignée de main) sont des actes symboliques forts : ils marquent la fin de l'hostilité et le retour à la vie commune. Au Cameroun, le Code de procédure civile reconnaît d'ailleurs la valeur juridique des accords issus de la médiation traditionnelle (loi n°2005/007 du 27 juillet 2005).
\end{savaistu}

\coursec{Comprehension questions}

Maintenant que vous savez \emph{comment} lire le texte (skimming pour l'idée générale, scanning pour les détails, déduction pour les mots inconnus), il est temps d'utiliser ces stratégies pour \emph{répondre aux questions}. C'est l'étape qui rapporte des points : à l'examen comme en classe, une bonne compréhension ne vaut rien si la réponse est mal formulée ou mal justifiée. Dans cette section, nous allons apprendre à reconnaître les quatre grands types de questions de compréhension, à adopter la bonne stratégie pour chacun, et à rédiger des réponses complètes et correctes en anglais.

\courssub{The method: five golden steps}

Avant de voir les types de questions, retenons la démarche générale, valable pour toutes.

\begin{methode}[How to answer comprehension questions]
\begin{enumerate}
\item \cle{Read the questions before the text.} Lisez d'abord toutes les questions : elles vous indiquent quelles informations chercher pendant la lecture.
\item \cle{Underline the key words.} Soulignez les mots-clés de chaque question (noms propres, verbes, mots interrogatifs : \eng{who}, \eng{when}, \eng{why}, \eng{how long}).
\item \cle{Locate the passage.} Repérez dans le texte le passage qui contient la réponse. Astuce : les questions suivent presque toujours l'ordre du texte.
\item \cle{Answer with a complete sentence.} Rédigez une phrase complète en reprenant les mots de la question, pas un simple mot isolé.
\item \cle{Quote the text.} Justifiez votre réponse en citant la ligne du texte (entre guillemets anglais “ … ”, avec le numéro de ligne ou le paragraphe si le texte en comporte).
\end{enumerate}
\end{methode}

\begin{attention}[A “Yes” or “No” alone is not enough]
Au Bac, une réponse \eng{Yes} ou \eng{No} (ou \eng{True} / \eng{False}) \textbf{sans justification ne rapporte pas tous les points}. Il faut toujours ajouter la preuve tirée du texte. Mauvaise réponse : \eng{False.} Bonne réponse : \eng{False. The text says the dispute was about a piece of land, not money (paragraph 1: “a bitter quarrel over a piece of land behind the market”).} Retenez la formule magique : \cle{answer + quotation}.
\end{attention}

\courssub{The reading text (reminder)}

Pour travailler, reprenons le texte support de la leçon. Relisez-le rapidement (skimming) avant d'attaquer les questions.

\begin{textebox}[Settling a dispute the traditional way]
In the village of Mankon, near Bamenda, two families, the Ngongs and the Tantohs, stopped greeting each other because of a bitter quarrel over a piece of land behind the market. Both families claimed that their grandfather had planted the first mango tree on that land, and for three years nobody dared to farm it. The conflict became so serious that the children of the two families could not play together any more.

One Saturday morning, the village chief invited the two families to the palace square. The elders sat in a circle under the big fig tree, and each family presented its case calmly. The meeting lasted four hours. The elders listened carefully, asked questions, and then consulted the oldest man of the village, who remembered the history of the land. Finally, the chief gave his judgement: the land would be divided into two equal parts, and each family would pay a goat and a calabash of palm wine to seal the peace.

The two families shook hands, shared a meal, and the children played football together that same evening. “Going to court would have cost us money, time and friendship,” said Mr Ngong. “Here, dialogue cost us almost nothing and gave us back our neighbours.”

As the elders of Mankon like to say: “When two elephants fight, it is the grass that suffers; when they talk, the grass grows again.”
\end{textebox}

\begin{definition}[Glossary]
\begin{itemize}
\item \eng{a dispute} (n.) : un différend, un conflit
\item \eng{a quarrel} (n.) : une querelle, une dispute
\item \eng{to claim} (v.) : prétendre, affirmer
\item \eng{an elder} (n.) : un ancien, un notable
\item \eng{to seal the peace} (exp.) : sceller la paix
\item \eng{a calabash} (n.) : une calebasse
\item \eng{a judgement} (n.) : un jugement, une décision
\item \eng{to shake hands} (exp.) : se serrer la main
\end{itemize}
\end{definition}

\courssub{Type 1: True or False questions (with justification)}

Le principe : on vous donne une affirmation et vous devez dire si elle est \eng{True} ou \eng{False}, \textbf{puis citer la phrase du texte qui le prouve}. Ne vous contentez jamais de votre impression : cherchez la preuve écrite.

\begin{exemplebox}[Example]
\eng{The dispute was about money.} \par
Réponse : \eng{False. The dispute was about a piece of land: “a bitter quarrel over a piece of land behind the market” (paragraph 1).} \par
Traduction : « Faux. Le conflit portait sur un terrain : “une âpre querelle à propos d'un terrain derrière le marché”. »
\end{exemplebox}

Autres exemples sur le même texte :
\begin{itemize}
\item \eng{The two families stopped talking to each other.} → \eng{True: “two families … stopped greeting each other.”}
\item \eng{The meeting took place on a Sunday.} → \eng{False: “One Saturday morning, the village chief invited the two families.”}
\end{itemize}

\begin{attention}[Watch out for small changes]
Les affirmations fausses contiennent souvent \textbf{un seul détail modifié} : le jour (\eng{Sunday} au lieu de \eng{Saturday}), le lieu, la personne, la cause. Lisez l'affirmation mot par mot et comparez-la au texte. Ne répondez pas de mémoire !
\end{attention}

\courssub{Type 2: Multiple choice questions (MCQ)}

Le principe : une question, trois ou quatre options, une seule correcte. La stratégie gagnante est l'\cle{élimination} : barrez mentalement les options impossibles avant de choisir.

\begin{exemplebox}[Example]
\eng{Who settled the dispute? a) the police b) the chief and the elders c) a judge d) the families alone} \par
Élimination : la police n'apparaît nulle part dans le texte (a éliminé) ; aucun juge ni tribunal n'est mentionné — au contraire, \eng{“Going to court would have cost us money”} prouve qu'ils n'y sont \emph{pas} allés (c éliminé) ; les familles ne se sont pas réconciliées seules, puisque le chef les a invitées (d éliminé). \par
Réponse : \eng{b) the chief and the elders — “the village chief invited the two families … The elders sat in a circle.”}
\end{exemplebox}

\courssub{Type 3: Short answer questions}

Le principe : la question commence par un mot interrogatif (\eng{who}, \eng{what}, \eng{when}, \eng{how long}, \eng{how much}, \eng{where}, \eng{why}). Utilisez le \cle{scanning} : cherchez dans le texte le mot-clé de la question (un chiffre, une durée, un objet), puis répondez par une phrase complète.

\begin{exemplebox}[Examples]
\eng{How long did the meeting last?} → \eng{The meeting lasted four hours.} (« La réunion a duré quatre heures. ») \par
\eng{What did each family pay?} → \eng{Each family paid a goat and a calabash of palm wine to seal the peace.} (« Chaque famille a payé une chèvre et une calebasse de vin de palme pour sceller la paix. ») \par
\eng{Why did the two families stop greeting each other?} → \eng{Because they quarrelled over a piece of land.} (« Parce qu'elles se sont querellées à propos d'un terrain. »)
\end{exemplebox}

Remarquez la construction des réponses : on \textbf{reprend les mots de la question} pour former une phrase complète. Question : \eng{How long did the meeting last?} → Réponse : \eng{The meeting lasted four hours.} C'est la même structure, avec l'information ajoutée. Évitez la réponse télégraphique \eng{Four hours} : elle peut être acceptée, mais la phrase complète est plus sûre et montre votre maîtrise de l'anglais.

\courssub{Type 4: Interpretation and personal questions}

Ces questions demandent d'aller \emph{au-delà} du texte : expliquer une idée, dégager une leçon, donner son avis. La réponse n'est pas écrite mot pour mot dans le texte, mais elle doit \textbf{s'appuyer sur lui}.

\begin{exemplebox}[Example 1: interpretation]
\eng{Why is dialogue better than going to court, according to the text?} \par
Réponse : \eng{According to Mr Ngong, dialogue is better because going to court costs money, time and friendship, whereas dialogue costs almost nothing and gives people back their neighbours.} \par
Traduction : « Selon M. Ngong, le dialogue est meilleur parce que le tribunal coûte de l'argent, du temps et de l'amitié, alors que le dialogue ne coûte presque rien et rend aux gens leurs voisins. »
\end{exemplebox}

\begin{exemplebox}[Example 2: the final proverb]
\eng{Explain the proverb “When two elephants fight, it is the grass that suffers; when they talk, the grass grows again” in your own words.} \par
Réponse : \eng{In my own words, this proverb means that when powerful people quarrel, it is the weak and innocent people around them — like the children of the two families — who suffer. But when they choose dialogue, peace returns and everyone can live and prosper again.} \par
Traduction : « Ce proverbe signifie que quand les puissants se querellent, ce sont les faibles et les innocents autour d'eux — comme les enfants des deux familles — qui souffrent. Mais quand ils choisissent le dialogue, la paix revient et chacun peut vivre et prospérer à nouveau. »
\end{exemplebox}

Pour exprimer votre opinion ou paraphraser, utilisez ces amorces : \eng{In my opinion, …}, \eng{In my own words, …}, \eng{This means that …}, \eng{According to the text, …}, \eng{I think that … because …}.

\courssub{Summary table: one strategy per question type}

\begin{center}
\begin{tabularx}{\linewidth}{|X|X|X|}
\hline
\rowcolor{popblueL} \textbf{Type of question} & \textbf{What to do} & \textbf{Key phrase} \\
\hline
True / False & Justify with a line from the text & “The text says …” \\
\hline
MCQ & Eliminate the wrong options & “a and c are impossible because …” \\
\hline
Short answer & Find the key words, answer with a full sentence & “The meeting lasted …” \\
\hline
Interpretation / opinion & Use the text + your own words & “In my opinion, …” \\
\hline
Proverb / meaning & Paraphrase simply & “This means that …” \\
\hline
\end{tabularx}
\end{center}

\begin{popfigure}
\begin{tikzpicture}
\node[draw=popblue, fill=popblueL, rounded corners, align=center, font=\bfseries] (q) at (0,0) {Comprehension\\question};
\node[draw=popgreen, fill=popgreenL, rounded corners, align=center] (tf) at (-5.2,-2.2) {True / False\\justify with a line};
\node[draw=poporange, fill=poporangeL, rounded corners, align=center] (mcq) at (-1.75,-2.2) {MCQ\\eliminate wrong options};
\node[draw=poppurple, fill=poppurpleL, rounded corners, align=center] (sa) at (1.75,-2.2) {Short answer\\find key words};
\node[draw=poppink, fill=poppinkL, rounded corners, align=center] (op) at (5.2,-2.2) {Opinion / proverb\\use your own words};
\draw[-{Stealth}, thick, popdark] (q) -- (tf);
\draw[-{Stealth}, thick, popdark] (q) -- (mcq);
\draw[-{Stealth}, thick, popdark] (q) -- (sa);
\draw[-{Stealth}, thick, popdark] (q) -- (op);
\end{tikzpicture}
\legende{Chaque type de question appelle une stratégie différente.}
\end{popfigure}

\courssub{Worked example}

\begin{exoresolu}[A complete set of questions]
Répondez aux questions suivantes sur le texte “Settling a dispute the traditional way”. \par
1. True or False: \eng{The children of the two families kept playing together during the conflict.} \par
2. MCQ: \eng{Where did the meeting take place? a) in the market b) at the palace square c) in the court d) in the church} \par
3. Short answer: \eng{How was the land divided?} \par
4. Interpretation: \eng{Why could the children not play together any more?}
\tcblower
1. \eng{False. The text says “the children of the two families could not play together any more” (paragraph 1).} — On cite la ligne exacte qui contredit l'affirmation. \par
2. \eng{b) at the palace square — “the village chief invited the two families to the palace square.”} Le marché est le lieu du terrain disputé, pas de la réunion ; le tribunal et l'église ne sont pas mentionnés. \par
3. \eng{The land was divided into two equal parts.} — Phrase complète qui reprend les mots de la question. \par
4. \eng{They could not play together any more because the conflict between their parents had become very serious.} — La cause n'est pas écrite mot pour mot, mais elle se déduit de “The conflict became so serious that …”.
\end{exoresolu}

\begin{attention}[Common mistakes of francophone students]
\begin{itemize}
\item Traduire la question au lieu d'y répondre en anglais : toute la réponse doit être en anglais.
\item Recopier un long passage du texte sans répondre vraiment : citez \textbf{une courte phrase} pertinente, pas trois lignes.
\item Oublier la justification après \eng{True} ou \eng{False}.
\item Répondre \eng{Because of the land} au lieu d'une phrase complète : préférez \eng{Because they quarrelled over a piece of land.}
\item Inventer des informations absentes du texte : même pour les questions d'opinion, appuyez-vous sur ce que le texte dit.
\end{itemize}
\end{attention}

\begin{exercice}[Your turn]
Répondez en phrases complètes, avec justification quand c'est nécessaire. \par
1. True or False: \eng{The oldest man of the village remembered the history of the land.} \par
2. MCQ: \eng{What did the families do after the judgement? a) they went to court b) they shook hands and shared a meal c) they left the village d) they sold the land} \par
3. Short answer: \eng{Where was the disputed piece of land?} \par
4. Short answer: \eng{How long did the conflict last before the meeting?} \par
5. Interpretation: \eng{Do you think traditional chiefs are good judges? Why or why not? (two or three sentences)}
\end{exercice}

\begin{corrige}[Exercice « Your turn »]
1. \eng{True: “they consulted the oldest man of the village, who remembered the history of the land.”} \par
2. \eng{b) they shook hands and shared a meal — “The two families shook hands, shared a meal.”} \par
3. \eng{The disputed piece of land was behind the market.} \par
4. \eng{The conflict lasted three years before the meeting.} \par
5. Réponse possible : \eng{In my opinion, traditional chiefs are good judges because they know the families and the history of the village. Moreover, their solution costs almost nothing and restores friendship, unlike a court decision.}
\end{corrige}

\begin{savaistu}[Reading questions at the BAC]
At the BAC, reading comprehension questions are worth several points and almost always follow the order of the text: the answer to question 1 is near the beginning, and the answer to the last question is near the end. Use this simple trick to locate each answer quickly instead of reading the whole text again for every question!
\end{savaistu}

\begin{aretenir}[Key points]
\begin{itemize}
\item Lisez les questions \textbf{avant} le texte et soulignez les mots-clés.
\item \eng{True/False} : toujours ajouter la citation qui justifie.
\item QCM : éliminez les options impossibles avant de choisir.
\item Questions courtes : scanning + phrase complète qui reprend les mots de la question.
\item Questions d'interprétation et proverbes : paraphrasez avec \eng{In my own words, …} en restant fidèle au texte.
\item Formule d'or : \cle{answer + quotation}.
\end{itemize}
\end{aretenir}

\coursec{Thematic vocabulary: peace and conflict}

After checking your answers to the comprehension questions, you now have a good understanding of the text’s main ideas. To discuss the topic fluently — whether in speaking or writing — you need to master the specific vocabulary of conflict and resolution. This section builds your lexical toolkit through observation, classification, and practice.

\courssub{Conflict vocabulary: naming the problem}

\begin{textebox}{Lexical field: Conflict}
In many communities, a \cle{dispute} arises over land boundaries. It may start as a simple \cle{disagreement} but quickly turn into a heated \cle{quarrel}. If tempers flare, it becomes a \cle{fight}. The \cle{tension} can last for months if nobody intervenes.
\end{textebox}

\begin{definition}[Core conflict nouns]
\begin{itemize}
    \item \textbf{Dispute} « dis-pyou-t » : a serious argument or disagreement, often legal or formal (conflit, litige).
    \item \textbf{Quarrel} « kwo-rel » : an angry argument between people who know each other (querelle, dispute).
    \item \textbf{Argument} « ar-gyu-ment » : an exchange of opposing views, often heated (dispute, altercation). \textbf{Attention :} \eng{argument} $\neq$ \og argumentation \fg{} (raisonnement) mais \og dispute \fg.
    \item \textbf{Disagreement} « dis-a-gri-ment » : lack of consensus, milder than dispute (désaccord).
    \item \textbf{Fight} « fait » : physical or violent confrontation (bagarre, combat).
    \item \textbf{Tension} « ten-chone » : strained relations, latent conflict (tension).
\end{itemize}
\end{definition}

\begin{center}
\begin{tabularx}{\linewidth}{|X|X|X|}
\hline
\rowcolor{popblueL} \textbf{English} & \textbf{Français} & \textbf{Context / Collocation} \\
\hline
a land dispute & un litige foncier & \eng{They have a land dispute.} (Ils ont un litige foncier.) \\
a family quarrel & une querelle familiale & \eng{A petty quarrel over money.} (Une mesquine querelle pour de l'argent.) \\
a heated argument & une vive dispute & \eng{They had a heated argument.} (Ils ont eu une vive dispute.) \\
a minor disagreement & un léger désaccord & \eng{Despite minor disagreements, they agreed.} (Malgré de légers désaccords, ils se sont mis d'accord.) \\
a violent fight & une bagarre violente & \eng{The fight broke out at the market.} (La bagarre a éclaté au marché.) \\
rising tension & tension croissante & \eng{There is rising tension in the village.} (Il y a une tension croissante au village.) \\
\hline
\end{tabularx}
\end{center}

\courssub{Resolution vocabulary: fixing the problem}

\begin{definition}[Core resolution nouns and verbs]
\begin{itemize}
    \item \textbf{To settle} « se-tel » : to resolve a dispute definitively (régler, résoudre). \eng{They settled the dispute.} (Ils ont réglé le conflit.)
    \item \textbf{To solve} « solv » : to find an answer to a problem (résoudre). Used for \eng{problem}, \eng{mystery}, \eng{puzzle}, less for \eng{dispute}.
    \item \textbf{Mediation} « mi-di-eï-chone » : intervention by a neutral third party (médiation).
    \item \textbf{Mediator} « mi-di-eï-teur » : the neutral person (médiateur).
    \item \textbf{Negotiation} « ne-go-chi-eï-chone » : discussion to reach agreement (négociation).
    \item \textbf{Compromise} « kom-pro-maïz » : each side gives up something (compromis).
    \item \textbf{Agreement} « a-gri-ment » : the result of negotiation (accord).
    \item \textbf{Peace} « pis » : state of harmony (paix).
\end{itemize}
\end{definition}

\begin{exemplebox}{Verbs of resolution in context}
\eng{The chief helped \textbf{settle} the land dispute.} (Le chef a aidé à régler le litige foncier.) \\
\eng{We must \textbf{mediate} between the two families.} (Nous devons médier entre les deux familles.) \\
\eng{They \textbf{negotiated} for three hours.} (Ils ont négocié pendant trois heures.) \\
\eng{Both parties reached a \textbf{compromise}.} (Les deux parties sont parvenues à un compromis.) \\
\eng{They signed a peace \textbf{agreement}.} (Ils ont signé un accord de paix.)
\end{exemplebox}

\begin{propriete}[Reach an agreement vs Make an agreement]
\eng{Reach an agreement} (parvenir à un accord) insiste sur le processus et l'effort. \\
\eng{Make an agreement} (conclure un accord) est plus neutre. \\
\eng{Sign an agreement} (signer un accord) est l'acte formel final. \\
\textbf{Ne dites pas :} \eng{*do an agreement} (calque de \og faire un accord \fg).
\end{propriete}

\courssub{Actors and roles: who does what?}

\begin{center}
\begin{tabularx}{\linewidth}{|X|X|X|}
\hline
\rowcolor{popgreenL} \textbf{English} & \textbf{Français} & \textbf{Role in dispute settlement} \\
\hline
Chief / Traditional ruler & Chef traditionnel & Presides over the council, delivers judgement. \\
Elder & Ancien, doyen & Advises, shares wisdom, mediates informally. \\
Quarter head & Chef de quartier & First point of contact for local disputes. \\
Mediator & Médiateur & Neutral facilitator, helps parties talk. \\
Judge / Magistrate & Juge / Magistrat & Formal legal authority, binding decision. \\
Court & Tribunal & Formal institution (customary or statutory). \\
Witness & Témoin & Gives evidence, tells what they saw/heard. \\
Neighbour / Neighbor (US) & Voisin & Often a witness or party to the dispute. \\
\hline
\end{tabularx}
\end{center}

\begin{exemplebox}{Actors in action}
\eng{The \textbf{quarter head} summoned both neighbours.} (Le chef de quartier a convoqué les deux voisins.) \\
\eng{The \textbf{elder} acted as a \textbf{mediator}.} (L'ancien a agi en médiateur.) \\
\eng{The \textbf{witness} testified before the \textbf{court}.} (Le témoin a témoigné devant le tribunal.)
\end{exemplebox}

\courssub{Useful expressions and collocations}

\begin{methode}{Key phrases for talking about dispute settlement}
\begin{enumerate}
    \item \textbf{To reach an agreement} « ri-tch an a-gri-ment » = parvenir à un accord.
    \item \textbf{To shake hands} « cheik hendz » = se serrer la main (signe de réconciliation).
    \item \textbf{To listen to both sides} « li-sne tou bo-th saïdz » = écouter les deux parties.
    \item \textbf{To apologise (to sb)} « a-po-lo-dja-ïz » = présenter ses excuses (à qqn).
    \item \textbf{To forgive (sb)} « fer-giv » = pardonner (à qqn). \textit{Irregular: forgive – forgave – forgiven.}
    \item \textbf{To live in peace} « liv in pis » = vivre en paix.
    \item \textbf{To pay a fine} « peï eï faïn » = payer une amende.
    \item \textbf{To restore harmony} « ris-tor ar-mo-ni » = restaurer l'harmonie.
\end{enumerate}
\end{methode}

\begin{exemplebox}{Collocations in full sentences}
\eng{After the mediation, they \textbf{shook hands} and \textbf{apologised to each other}.} (Après la médiation, ils se sont serré la main et se sont présentés des excuses mutuelles.) \\
\eng{The mediator \textbf{listened to both sides} before proposing a solution.} (Le médiateur a écouté les deux parties avant de proposer une solution.) \\
\eng{He \textbf{forgave} his brother and they \textbf{live in peace} now.} (Il a pardonné à son frère et ils vivent en paix maintenant.)
\end{exemplebox}

\courssub{Word families: building your lexicon}

\begin{propriete}[Derivation: Nouns, Verbs, Adjectives, Adverbs]
English builds families using suffixes. Mastering them multiplies your vocabulary.
\begin{center}
\begin{tabular}{|l|l|l|l|l|}
\hline
\rowcolor{poppurpleL} \textbf{Verb} & \textbf{Noun (person/thing)} & \textbf{Noun (abstract)} & \textbf{Adjective} & \textbf{Adverb} \\
\hline
agree & agreement & — & agreeable & agreeably \\
disagree & disagreement & — & disagreeable & disagreeably \\
mediate & mediator & mediation & — & — \\
negotiate & negotiator & negotiation & negotiable & — \\
compromise & — & compromise & compromised & — \\
peace & peacemaker & peace & peaceful & peacefully \\
forgive & forgiveness & — & forgiving & forgivingly \\
settle & settler & settlement & settled & — \\
\hline
\end{tabular}
\end{center}
\end{propriete}

\begin{attention}[Faux amis : \eng{Argument}, \eng{Actually}, \eng{Assist}]
\begin{itemize}
    \item \textbf{\eng{Argument}} = \og dispute, querelle \fg{} (et non \og argumentation \fg).
        \begin{itemize}
            \item \eng{They had an argument.} = Ils ont eu une dispute.
            \item \eng{Argumentation} (raisonnement) = \eng{reasoning} ou \eng{line of argument}.
        \end{itemize}
    \item \textbf{\eng{Actually}} = \og en fait, en réalité \fg{} (et non \og actuellement \fg = \eng{currently} / \eng{at present}).
        \begin{itemize}
            \item \eng{Actually, the land belongs to my uncle.} = \og En fait, la terre appartient à mon oncle. \fg
            \item \eng{Currently, he lives in Douala.} = \og Actuellement, il habite à Douala. \fg
        \end{itemize}
    \item \textbf{\eng{To assist}} = \og aider, assister (qqn) \fg{} (et non \og assister à \fg = \eng{to attend}).
        \begin{itemize}
            \item \eng{The mediator assisted the parties.} = Le médiateur a aidé les parties.
            \item \eng{He attended the meeting.} = Il a assisté à la réunion.
        \end{itemize}
\end{itemize}
\end{attention}

\courssub{Lexical mind map: Conflict — Solution — People}

\begin{popfigure}
\begin{tikzpicture}[
    mindmap,
    grow cyclic,
    every node/.style={concept, execute at begin node=\hskip0pt},
    concept color=popblue!60,
    level 1/.style={level distance=4.5cm, sibling angle=120},
    level 2/.style={level distance=3cm, sibling angle=45},
]
\node[concept, text width=2.5cm, font=\bfseries\large] {DISPUTE SETTLEMENT}
    child[concept color=poporange!70] {
        node[concept, text width=2.2cm, font=\bfseries] {CONFLICT}
        child { node[concept, text width=2cm] {dispute} }
        child { node[concept, text width=2cm] {quarrel} }
        child { node[concept, text width=2cm] {argument} }
        child { node[concept, text width=2cm] {fight} }
        child { node[concept, text width=2cm] {tension} }
        child { node[concept, text width=2cm] {disagreement} }
    }
    child[concept color=popgreen!70] {
        node[concept, text width=2.2cm, font=\bfseries] {SOLUTION}
        child { node[concept, text width=2cm] {settle} }
        child { node[concept, text width=2cm] {mediate} }
        child { node[concept, text width=2cm] {negotiate} }
        child { node[concept, text width=2cm] {compromise} }
        child { node[concept, text width=2cm] {agreement} }
        child { node[concept, text width=2cm] {peace} }
        child { node[concept, text width=2cm] {forgive} }
    }
    child[concept color=poppurple!70] {
        node[concept, text width=2.2cm, font=\bfseries] {PEOPLE}
        child { node[concept, text width=2cm] {chief} }
        child { node[concept, text width=2cm] {elder} }
        child { node[concept, text width=2cm] {mediator} }
        child { node[concept, text width=2cm] {judge} }
        child { node[concept, text width=2cm] {witness} }
        child { node[concept, text width=2cm] {quarter head} }
        child { node[concept, text width=2cm] {neighbour} }
    };
\end{tikzpicture}
\legende{Carte mentale lexicale : trois pôles pour organiser le vocabulaire du règlement des conflits.}
\end{popfigure}

\courssub{Reading practice: a new text on mediation}

\begin{textebox}{Traditional Mediation in Buea}
In the quarter of Molyko, Buea, a serious \cle{dispute} broke out between two families over a piece of farmland. For weeks, \cle{tension} was high. The \cle{neighbours} avoided each other, and a \cle{violent fight} almost happened at the water point. The \cle{quarter head}, Mr. Ekema, intervened quickly. He summoned the family heads and \cle{witnesses} to the community hall. He acted as a \cle{mediator}: he \cle{listened to both sides} without taking sides. After two hours of \cle{negotiation}, he proposed a \cle{compromise}: the land would be shared, and the family who had cleared it would pay a small \cle{fine} to the other. Both parties \cle{accepted} the \cle{settlement}. They \cle{shook hands}, \cle{apologised}, and promised to \cle{live in peace}. The \cle{elder} closed the session with a prayer for \cle{harmony}.

\vspace{0.3cm}
\noindent\textbf{Glossaire :} \textit{farmland} (terre agricole), \textit{water point} (point d'eau), \textit{summoned} (convoqué), \textit{community hall} (salle communautaire), \textit{cleared} (défriché), \textit{session} (séance), \textit{harmony} (harmonie).
\end{textebox}

\begin{exercice}{Vocabulary check (from the text above)}
Find words in the text that match these definitions:
\begin{enumerate}
    \item A person who sees an event and can describe it. (\_ \_ \_ \_ \_ \_ \_ )
    \item An amount of money paid as punishment. (\_ \_ \_ \_ )
    \item To say sorry. (\_ \_ \_ \_ \_ \_ \_ \_ \_ \_ )
    \item A person who helps conflicting parties agree. (\_ \_ \_ \_ \_ \_ \_ \_ )
    \item The head of a neighbourhood. (\_ \_ \_ \_ \_ \_ \_ \_ \_ \_ \_ )
\end{enumerate}
\end{exercice}

\begin{corrige}{Exercice Vocabulary check}
\begin{enumerate}
    \item \textbf{Witness} (l. 5)
    \item \textbf{Fine} (l. 8)
    \item \textbf{Apologised} (l. 9) — \textit{UK spelling: apologise / US: apologize}
    \item \textbf{Mediator} (l. 6)
    \item \textbf{Quarter head} (l. 4)
\end{enumerate}
\end{corrige}

\courssub{Guided production: use the vocabulary}

\begin{exoresolu}{Writing a short report}
\textbf{Consigne :} Write 5 sentences reporting the Molyko dispute using at least 8 words from the mind map (conflict, solution, people).
\vspace{0.2cm}

\textbf{Solution détaillée :}
\begin{enumerate}
    \item A \textbf{dispute} over \textbf{farmland} arose between two families in Molyko. (conflict vocab)
    \item \textbf{Tension} rose until a \textbf{fight} nearly broke out. (conflict vocab)
    \item The \textbf{quarter head} acted as a \textbf{mediator} and \textbf{listened to both sides}. (people + solution)
    \item After \textbf{negotiation}, they reached a \textbf{compromise} and \textbf{shook hands}. (solution vocab)
    \item The \textbf{elder} prayed for \textbf{peace} and \textbf{harmony} in the community. (solution + people)
\end{enumerate}
\textit{Words used: dispute, farmland, tension, fight, quarter head, mediator, listened to both sides, negotiation, compromise, shook hands, elder, peace, harmony (13 words).}
\end{exoresolu}

\begin{savaistu}{African traditions of peace}
In \textbf{Ghana}, the \textbf{National Peace Council} (established 2011) mediates political and chieftaincy disputes nationwide. In \textbf{Rwanda}, after the 1994 genocide, traditional \textbf{Gacaca courts} (\eng{Inkiko Gacaca}) handled over a million cases — community elders judged perpetrators face-to-face, focusing on truth, \eng{apology}, and \eng{reintegration}. In \textbf{Cameroon}, the \textbf{Chefferies traditionnelles} (chiefdoms) in the Grassfields and the \eng{Lamidats} in the North still settle land and family disputes daily. Africa’s \eng{restorative justice} — repairing harm rather than just punishing — inspires modern mediation worldwide.
\end{savaistu}

\begin{aretenir}[Checklist : Vocabulary of Peace and Conflict]
\begin{itemize}
    \item \textbf{Conflict words :} dispute, quarrel, argument, fight, disagreement, tension.
    \item \textbf{Solution words :} settle, mediate, negotiate, compromise, agreement, peace, forgive, fine.
    \item \textbf{People :} chief, elder, quarter head, mediator, judge, witness, neighbour.
    \item \textbf{Key phrases :} reach an agreement, shake hands, listen to both sides, apologise, live in peace.
    \item \textbf{Faux amis :} argument = dispute ; actually = en fait ; assist = aider.
    \item \textbf{Familles :} agree $\to$ agreement / disagreement ; peace $\to$ peaceful / peacefully ; mediate $\to$ mediator / mediation.
\end{itemize}
\end{aretenir}

\coursec{Guided production: summary and opinion}

Dans la section précédente, nous avons construit ensemble le vocabulaire de la paix et du conflit : \eng{dispute}, \eng{mediator}, \eng{reconcile}, \eng{compromise}, \eng{peaceful settlement}. Ces mots ne doivent pas rester des listes apprises par cœur : ils doivent devenir des outils actifs. C'est exactement l'objectif de cette dernière étape de la leçon. Vous avez lu le texte, vous l'avez compris, vous avez enrichi votre lexique ; il s'agit maintenant de \cle{produire} : d'abord \cle{résumer} le texte avec vos propres mots, puis \cle{donner votre opinion} sur le règlement des conflits communautaires. C'est la compétence reine au baccalauréat et dans la vie : lire, digérer, reformuler, juger.

\courssub{What is a summary? Observation first}

Avant toute règle, observons. Voici deux versions d'un même résumé du texte étudié sur le règlement d'un conflit communautaire. Laquelle est la bonne ?

\begin{exemplebox}[Two summaries — which one is correct?]
\textbf{Version A :} “The text says that two families quarrelled over a piece of land behind the market. The text says that the chief called a meeting under the big tree. The text says that both families explained their side. The text says that they finally agreed to share the land.”

\textbf{Version B :} “Two neighbours disagreed about a plot of land. First, the village chief organised a meeting where each side could speak. Then, after a long discussion, they reached a compromise. Finally, peace returned to the community.”

\textbf{Réponse :} la version B. Pourquoi ? Elle est courte, elle ne copie aucune phrase du texte, elle utilise des mots personnels (\eng{neighbours} au lieu de \eng{families}, \eng{disagreed} au lieu de \eng{quarrelled}) et elle est structurée par des connecteurs (\eng{first}, \eng{then}, \eng{finally}). La version A répète \eng{the text says that} et colle au texte : c'est une paraphrase paresseuse, pas un résumé.
\end{exemplebox}

\begin{definition}[Summary — définition]
A \cle{summary} (un résumé) is a short text that gives only the \emph{main ideas} of a longer text, \emph{in your own words}, without details, examples or quotations. En français : un résumé est un texte bref qui ne reprend que les idées principales d'un texte plus long, avec vos propres mots, sans détails, sans exemples et sans citations.
\end{definition}

\courssub{Method: how to write a good summary}

\begin{methode}[Résumer un texte en 4 étapes]
\begin{enumerate}
\item \textbf{Repérer les idées principales.} Relisez le texte et soulignez une seule idée par paragraphe : le conflit, la réunion, la solution, le résultat. Ignorez les détails (noms propres, dates, descriptions).
\item \textbf{Changer les mots.} Pour chaque idée, cherchez un synonyme ou une reformulation : \eng{quarrel} $\rightarrow$ \eng{disagree} ; \eng{a piece of land} $\rightarrow$ \eng{a plot of land} ; \eng{they stopped fighting} $\rightarrow$ \eng{peace returned}.
\item \textbf{Assembler avec des connecteurs.} Écrivez une phrase par idée, reliées par \eng{first}, \eng{then}, \eng{after that}, \eng{finally}. Jamais plus de 3 ou 4 phrases.
\item \textbf{Vérifier.} Relisez votre résumé : aucune phrase ne doit être identique au texte ; quelqu'un qui n'a pas lu le texte doit comprendre l'essentiel.
\end{enumerate}
\end{methode}

\begin{propriete}[Les connecteurs de séquence]
Pour structurer un résumé chronologique (ou tout récit en étapes), l'anglais utilise des \cle{connecteurs de séquence} placés en début de phrase, suivis d'une virgule :

\begin{center}
\begin{tabularx}{\linewidth}{|l|X|X|}\hline
\rowcolor{popblueL} \textbf{Connecteur} & \textbf{Français} & \textbf{Exemple} \\ \hline
\eng{First,} & d'abord, tout d'abord & \eng{First, the two families argued about the land.} \\ \hline
\eng{Then,} & ensuite, puis & \eng{Then, the chief called a meeting.} \\ \hline
\eng{After that,} & après cela, ensuite & \eng{After that, each side explained its point of view.} \\ \hline
\eng{Finally,} & finalement, enfin & \eng{Finally, they signed an agreement.} \\ \hline
\end{tabularx}
\end{center}

\textbf{Emploi :} on les utilise dans l'ordre chronologique des événements ; on ne commence jamais par \eng{finally}. \textbf{Attention :} \eng{firstly} existe mais est soutenu ; préférez \eng{first}. Ne confondez pas \eng{finally} (enfin, à la fin) et \eng{at last} (enfin, après une longue attente) : \eng{Finally, they agreed.} = « Finalement, ils se sont mis d'accord. »
\end{propriete}

\begin{attention}[Piège n°1 : copier le texte au lieu de reformuler]
L'erreur la plus fréquente des élèves francophones est de recopier des phrases entières du texte en changeant un mot ou deux. Au baccalauréat, le correcteur compare votre production avec le texte : toute phrase copiée est sanctionnée, car le résumé doit être \cle{personnel}. Autre piège : traduire mot à mot depuis le français. « Les deux familles ont fait la paix » ne se traduit pas par \eng{the two families made the peace} mais par \eng{the two families made peace} (sans article) ou mieux : \eng{the two families were reconciled}. De même, « selon moi » ne se dit pas \eng{according to me} (faux ami !) mais \eng{in my opinion} ou \eng{to my mind}.
\end{attention}

\courssub{A model text: summary and opinion}

Lisons maintenant un modèle complet de production écrite, rédigé par une élève de Première après la lecture du texte sur le conflit de quartier. Observez la structure : quatre phrases de résumé, puis l'opinion, puis une conclusion.

\begin{textebox}[A student's production — “Settling a dispute in my community”]
The text is about a land dispute between two families in a neighbourhood of Bamenda. First, both families claimed the same plot of land and the quarrel became violent. Then, the traditional chief invited them to a palaver under the village tree, where everybody could speak freely. After that, the elders proposed a compromise: the land would be divided into two equal parts. Finally, the two families shook hands and peace returned to the community.

In my opinion, traditional chiefs are the best mediators because they know everybody and they are respected. I believe that dialogue is better than court because it is faster, cheaper and it keeps people friends. However, I think that serious crimes should go to court, because the law must protect everybody. To conclude, I agree with the villagers: talking under a tree can solve more problems than shouting in a courtroom. Peace is not the absence of conflict; it is the ability to handle conflict by peaceful means.
\end{textebox}

\textbf{Glossaire :}

\begin{center}
\begin{tabularx}{\linewidth}{|l|l|X|}\hline
\rowcolor{popblueL} \textbf{Mot anglais} & \textbf{Nature} & \textbf{Traduction} \\ \hline
\eng{to claim} & verbe & revendiquer, prétendre posséder \\ \hline
\eng{a palaver} & nom & palabre (réunion de discussion traditionnelle) \\ \hline
\eng{the elders} & nom pluriel & les anciens, les sages du village \\ \hline
\eng{to shake hands} & expression & se serrer la main (signe de réconciliation) \\ \hline
\eng{a courtroom} & nom & salle d'audience, tribunal \\ \hline
\eng{to handle} & verbe & gérer, traiter (un problème) \\ \hline
\end{tabularx}
\end{center}

\textbf{Questions sur le modèle :} (1) How many sentences does the summary contain? (2) Underline the four sequence connectors. (3) Which expressions show the writer's opinion? (4) Does the writer agree or disagree with traditional mediation? Justify with one sentence from the text.

\courssub{Expressing your opinion}

Après le résumé, l'examinateur attend votre \cle{opinion personnelle}, justifiée. L'anglais dispose d'expressions codifiées qu'il faut connaître par cœur.

\begin{propriete}[Les expressions d'opinion]
\begin{center}
\begin{tabularx}{\linewidth}{|X|X|X|}\hline
\rowcolor{popblueL} \textbf{Expression} & \textbf{Français} & \textbf{Exemple} \\ \hline
\eng{I think (that)...} & je pense que... & \eng{I think that dialogue is the best solution.} \\ \hline
\eng{In my opinion,...} & à mon avis,... & \eng{In my opinion, chiefs are good mediators.} \\ \hline
\eng{I believe (that)...} & je crois que... & \eng{I believe that peace starts at home.} \\ \hline
\eng{I agree (with...) because...} & je suis d'accord (avec...) parce que... & \eng{I agree with the chief because he was fair.} \\ \hline
\eng{I disagree because...} & je ne suis pas d'accord parce que... & \eng{I disagree because courts are too slow.} \\ \hline
\eng{To conclude,...} & pour conclure,... & \eng{To conclude, peace is everybody's business.} \\ \hline
\end{tabularx}
\end{center}

\textbf{Règle :} après \eng{I think}, \eng{I believe}, on peut garder ou supprimer \eng{that} : \eng{I think (that) dialogue is better.} Une opinion doit toujours être \cle{justifiée} par \eng{because} : une opinion sans raison ne vaut presque rien à l'examen.
\end{propriete}

\begin{exemplebox}[Exemples traduits]
\eng{In my opinion, dialogue is better than court because it is faster and cheaper.} « À mon avis, le dialogue vaut mieux que le tribunal parce qu'il est plus rapide et moins cher. »

\eng{I think that traditional chiefs are the best mediators because they know the customs of the village.} « Je pense que les chefs traditionnels sont les meilleurs médiateurs parce qu'ils connaissent les coutumes du village. »

\eng{I disagree with violent demonstrations because they destroy the community.} « Je ne suis pas d'accord avec les manifestations violentes parce qu'elles détruisent la communauté. »

\eng{I believe that every dispute can be settled peacefully if people listen to each other.} « Je crois que tout conflit peut être réglé pacifiquement si les gens s'écoutent les uns les autres. »
\end{exemplebox}

\begin{attention}[Piège n°2 : les faux amis de l'opinion]
\begin{itemize}
\item « Selon moi » ne se dit \textbf{jamais} \eng{according to me} : dites \eng{in my opinion} ou \eng{to me}.
\item « Je suis d'accord » = \eng{I agree} (sans \eng{am} !). Ne dites pas \eng{I am agree} : \eng{agree} est un verbe, pas un adjectif.
\item « À mon avis » au singulier : \eng{in my opinion}, jamais \eng{in my opinions}.
\item \eng{to think} au passé : \eng{thought} — \eng{I thought the meeting was useful.} « Je pensais que la réunion était utile. »
\end{itemize}
\end{attention}

\courssub{The staircase of production: from summary to conclusion}

Une bonne production écrite monte trois marches : d'abord le résumé en quatre étapes, puis l'opinion, puis la conclusion. Visualisons cet escalier.

\begin{popfigure}
\begin{tikzpicture}[font=\small]
\draw[fill=popblueL, draw=popblue] (0,0) rectangle (3.2,1);
\node[align=center, text=popdark] at (1.6,0.5) {1. Summary\\ \eng{First... Then...}};
\draw[fill=popgreenL, draw=popgreen] (3.2,1) rectangle (6.4,2);
\node[align=center, text=popdark] at (4.8,1.5) {2. Summary (suite)\\ \eng{After that... Finally...}};
\draw[fill=poporangeL, draw=poporange] (6.4,2) rectangle (9.6,3);
\node[align=center, text=popdark] at (8,2.5) {3. Opinion\\ \eng{I think... because...}};
\draw[fill=poppurpleL, draw=poppurple] (9.6,3) rectangle (12.8,4);
\node[align=center, text=popdark] at (11.2,3.5) {4. Conclusion\\ \eng{Peace is...}};
\draw[-{Stealth}, very thick, popmuted] (0,-0.4) -- (12.8,-0.4);
\node[text=popmuted] at (6.4,-0.8) {de la compréhension vers l'expression personnelle};
\end{tikzpicture}
\legende{L'escalier de la production écrite : quatre étapes de résumé, une opinion justifiée, une conclusion.}
\end{popfigure}

\begin{exoresolu}[Guided writing — corrigé type]
\textbf{Énoncé :} Read the text about the land dispute in Bamenda again. Write 5--6 lines: first summarise the text in your own words (use \eng{first}, \eng{then}, \eng{after that}, \eng{finally}), then give your opinion about traditional mediation.
\tcblower
\textbf{Solution détaillée (modèle de réponse) :}

\eng{The text is about a conflict between two families over a piece of land. First, the families quarrelled and the situation became dangerous. Then, the chief organised a palaver where both sides spoke. After that, the elders suggested sharing the land equally. Finally, the families made peace and the community was united again. In my opinion, traditional mediation is very useful because it is free and it keeps neighbours together. I believe that dialogue is better than court because a palaver ends with a handshake, not with anger.}

\textbf{Commentaires de méthode :} (1) le résumé contient exactement quatre phrases, une par étape du plan ; (2) aucune phrase n'est copiée du texte (\eng{quarrelled} remplace \eng{fought}, \eng{suggested sharing} remplace \eng{proposed a compromise}) ; (3) les quatre connecteurs de séquence structurent le résumé ; (4) l'opinion est introduite par \eng{in my opinion} et \eng{I believe}, et chaque opinion est justifiée par \eng{because} ; (5) la production fait bien 5--6 lignes : ni trop courte, ni hors sujet.
\end{exoresolu}

\begin{exercice}[Your turn — summary and opinion]
\begin{enumerate}
\item Réécrivez ce mauvais résumé en le raccourcissant et en le reformulant : \eng{“The text says that the two families fought for the land and the text says that the chief came and the text says that they agreed.”}
\item Complétez avec \eng{first}, \eng{then}, \eng{after that}, \eng{finally} : \eng{..., the neighbours stopped talking to each other. ..., the chief invited them to his palace. ..., they discussed for three hours. ..., they signed a peace agreement.}
\item Traduisez en anglais : « À mon avis, le dialogue est meilleur que la violence parce qu'il construit la paix. »
\item Rédigez votre propre production de 5--6 lignes : résumé du texte + votre opinion sur la question \eng{“Are traditional chiefs better mediators than judges?”}
\end{enumerate}
\end{exercice}

\begin{corrige}[Exercice : Your turn]
\begin{enumerate}
\item \eng{First, two families fought over a piece of land. Then, the chief intervened. Finally, they reached an agreement.} (Résumé court, reformulé, sans \eng{the text says that}.)
\item \eng{First, the neighbours stopped talking to each other. Then, the chief invited them to his palace. After that, they discussed for three hours. Finally, they signed a peace agreement.}
\item \eng{In my opinion, dialogue is better than violence because it builds peace.}
\item Production libre : vérifier la présence des quatre connecteurs, d'au moins une expression d'opinion (\eng{I think} / \eng{in my opinion} / \eng{I believe}) et d'une justification avec \eng{because}.
\end{enumerate}
\end{corrige}

\begin{aretenir}[L'essentiel de la section]
\begin{itemize}
\item Un \cle{résumé} = les idées principales uniquement, avec vos propres mots, en 3--4 phrases, sans rien copier du texte.
\item Le plan type : (1) the conflict, (2) the meeting, (3) the solution, (4) the result.
\item Les connecteurs de séquence : \eng{first}, \eng{then}, \eng{after that}, \eng{finally}.
\item Une opinion = une expression d'opinion (\eng{in my opinion}, \eng{I think}, \eng{I believe}, \eng{I agree/disagree}) + une justification avec \eng{because}.
\item Pièges : jamais \eng{according to me}, jamais \eng{I am agree}, jamais de phrase copiée du texte.
\end{itemize}
\end{aretenir}

\newpage

\coursec{Activité d'intégration}

\coursec{Lesson 7 — Reading: Texts about settling community disputes}

\courssub{Warm-up and Pre-reading}
\noindent Avant d'aborder les textes, activons les connaissances antérieures et le vocabulaire du thème.

\begin{experience}[Brainstorming : \eng{Peace and Conflict in our Communities}]
\noindent \textbf{Consigne :} En classe entière, l'enseignant écrit au tableau le mot \cle{Dispute}. Les élèves proposent des mots, expressions ou situations associés (en anglais si possible). L'enseignant regroupe en clusters au fur et à mesure.
\begin{center}
\begin{tikzpicture}[mindmap, grow cyclic, every node/.style=concept, concept color=popblue!60,
    level 1/.append style={level distance=4cm, sibling angle=90},
    level 2/.append style={level distance=3cm, sibling angle=45}]
\node {Dispute}
    child[concept color=poporange!70] {node[align=center]{Family\\ \eng{Inheritance, Land}}}
    child[concept color=popgreen!70] {node[align=center]{School\\ \eng{Prefects, Sports}}}
    child[concept color=poppurple!70] {node[align=center]{Neighbourhood\\ \eng{Boundary, Noise}}}
    child[concept color=poppink!70] {node[align=center]{Resolution\\ \eng{Mediation, Council, Compromise}}};
\end{tikzpicture}
\end{center}
\noindent \textbf{Question orale :} \eng{``In your village or quarter, who settles disputes before they go to the modern court?''} (Réponses attendues : \eng{Quarter Head, Fon, Elders, Family Head, Njangi}).
\end{experience}

\begin{propriete}[Stratégies de lecture : Rappel]
\begin{itemize}
    \item \textbf{Skimming (Lecture diagonale) :} Lire le titre, la 1\textsuperscript{re} phrase de chaque paragraphe, les noms propres. Objectif : saisir le \eng{gist} (l'idée générale).
    \item \textbf{Scanning (Repérage ciblé) :} Chercher une info précise (un nom, un chiffre, un lieu, un verbe d'action). Balayer le texte sans tout lire.
    \item \textbf{Guessing meaning from context (Déduction) :} 1. Identifier la nature du mot (N, V, Adj). 2. Analyser le co-texte (mots voisins, connecteurs \eng{but, because, such as}). 3. Formuler une hypothèse. 4. Vérifier dans le glossaire.
\end{itemize}
\end{propriete}

\courssub{Reading Texts: Case Studies for the Peace Council}
\noindent Voici trois cas authentiques inspirés du contexte camerounais (Régions du Nord-Ouest et du Sud-Ouest). Lisez-les en appliquant le \emph{skimming} (2 min par texte).

\begin{textebox}[Case Studies for the Peace Council of Bafut]
\textbf{Case A : The Kola Nut Tree Dispute}
\eng{Mr. Tanga and Mr. Ngwa are neighbours in the village of Agyati. A large kola nut tree stands exactly on the boundary line between their two compounds. Mr. Tanga claims the trunk is on his land, so the nuts belong to him. Mr. Ngwa argues that the branches hang over his compound, so he has a right to half the harvest. Last week, Mr. Tanga's son climbed the tree and picked all the ripe nuts. Mr. Ngwa's wife shouted at the boy, and Mr. Tanga pushed her away. The quarter head has referred the case to the Fon's Council.}

\textbf{Case B : The School Football Final}
\eng{At Government High School Buea, the final of the inter-class football tournament opposed Form 5 Arts and Form 5 Science. The match ended 1-1 after extra time. During the penalty shootout, the Form 5 Arts goalkeeper saved the last shot, but the Form 5 Science captain insists the goalkeeper moved off his line before the ball was kicked. The referee (a Form 3 student) did not see it. Form 5 Science refuses to accept the silver medals. The Principal has asked the Disciplinary Board (students and teachers) to mediate.}

\textbf{Case C : The Family Inheritance}
\eng{In a family in Bamenda, the father passed away six months ago. He left a house in Nkwen and a cocoa farm in Muyuka. He wrote no will. The eldest son (the heir according to tradition) wants to sell the cocoa farm and share the money. The three daughters and the younger son want to keep the farm and share the yearly profit. The eldest son has locked the gate of the farm. The family Njangi has failed. The case is now before the Family Council presided by the Quarter Head.}
\end{textebox}

\noindent \textbf{Glossaire des mots difficiles (pour l'élève) :}
\begin{center}
\begin{tabularx}{\linewidth}{|X|X|X|}
\hline
\rowcolor{popblueL} \textbf{Mot / Expression} & \textbf{Nature} & \textbf{Traduction / Sens contextuel} \\
\hline
\eng{boundary line} & n. & ligne de limite, borne séparative \\
\hline
\eng{compound} & n. & concession familiale close (maison + cour) \\
\hline
\eng{quarter head} & n. & chef de quartier (autorité administrative coutumière) \\
\hline
\eng{penalty shootout} & n. & séance de tirs au but \\
\hline
\eng{goalkeeper} & n. & gardien de but \\
\hline
\eng{off his line} & adv. & hors de sa ligne (de but), en position irrégulière \\
\hline
\eng{will} & n. & testament \\
\hline
\eng{heir} & n. & héritier (désigné par la coutume ou la loi) \\
\hline
\eng{yearly profit} & n. & bénéfice annuel, revenu de l'année \\
\hline
\eng{Njangi} & n. & réunion familiale / tontine (terme camerounais, anglais du Cameroun) \\
\hline
\eng{refer (a case)} & v. & déférer, soumettre (une affaire) à une autorité supérieure \\
\hline
\eng{mediate} & v. & médier, intervenir pour réconcilier \\
\hline
\end{tabularx}
\end{center}

\courssub{Reading Strategies in Action: Skimming, Scanning, Guessing}
\noindent Travail en groupes de 5. Chaque groupe reçoit \textbf{un} cas (A, B ou C).

\begin{methode}[Phase 1 — Analyse des faits (15 min)]
\begin{enumerate}
    \item \textbf{Skimming (3 min) :} Lisez votre cas rapidement. Répondez oralement dans le groupe :
    \begin{itemize}
        \item What is the general topic? (Land / Sport / Inheritance)
        \item Who are the two main parties?
        \item What is the core object of the dispute?
    \end{itemize}
    \item \textbf{Scanning (7 min) :} Relisez pour remplir le \textbf{Tableau des faits} ci-dessous dans votre cahier. Relevez les mots-clés exacts du texte.
    \item \textbf{Guessing (5 min) :} Sans dictionnaire, déduisez le sens de : \eng{compound, quarter head, heir, Njangi, referred, mediate}. Vérifiez avec le glossaire.
\end{enumerate}
\end{methode}

\begin{center}
\begin{tabularx}{\linewidth}{|X|X|}
\hline
\rowcolor{popblueL} \textbf{Information recherchée} & \textbf{Extrait du texte (Mots-clés)} \\
\hline
Parties en conflit (Party A / Party B) &  \\
\hline
Lieu précis du litige &  \\
\hline
Objet du conflit (arbre, match, ferme, maison) &  \\
\hline
Action déclenchante (dernier incident) &  \\
\hline
Autorité saisie (Quarter Head, Principal, Fon's Council) &  \\
\hline
Conséquence actuelle (violence, blocage, refus) &  \\
\hline
\end{tabularx}
\end{center}

\courssub{Guided Production: Role-play and Peace Agreement}

\begin{methode}[Phase 2 — Simulation de la médiation : \eng{The Peace Council} (20 min)]
\begin{enumerate}
    \item \textbf{Distribution des rôles (1 min) :}
    \begin{itemize}
        \item 1 élève = \eng{The Chairperson} (\eng{Fon} / \eng{Principal} / \eng{Quarter Head}) : préside, donne la parole, synthétise.
        \item 2 élèves = \eng{The Mediators / Elders} : posent des questions, rappellent la coutume ou le règlement, proposent des solutions.
        \item 1 élève = \eng{Party A} (\eng{Mr. Tanga} / \eng{Form 5 Science Captain} / \eng{Eldest Son}).
        \item 1 élève = \eng{Party B} (\eng{Mr. Ngwa} / \eng{Form 5 Arts Captain} / \eng{Daughters \& Younger Son}).
    \end{itemize}
    \item \textbf{Déroulement codifié (15 min) :}
    \begin{enumerate}
        \item \eng{Opening Statement} (Chairperson) : \eng{``We are here today to settle the dispute between [Party A] and [Party B] concerning [object]. Let us hear both sides.''}
        \item \eng{Hearing} (2 min par partie) : Chaque partie expose sa version. \textbf{Langue :} \eng{I claim that... / It is my right to... / According to custom...}
        \item \eng{Clarification} (Mediators) : \eng{``Can you prove the trunk is on your land?'' / ``What does the school regulation say about goalkeepers?'' / ``Where is the will?''}
        \item \eng{Deliberation} (Mediators, 2 min) : Concertation rapide. Proposition de \textbf{deux options} : une stricte (droit coutumier / règlement), une de compromis.
        \item \eng{Judgment / Ruling} (Chairperson) : Décision finale motivée. \eng{``Based on customary law / school rules, we rule that...''}
    \end{enumerate}
    \item \textbf{Outils linguistiques obligatoires :} Modaux (\eng{shall, must, should, have to}), connecteurs (\eng{however, therefore, consequently, whereas}), vocabulaire cible (\eng{settle, mediate, compromise, boundary, inheritance, unbiased, ruling, undertake, abide by}).
\end{enumerate}
\end{methode}

\begin{methode}[Phase 3 — Rédaction du \eng{Peace Agreement} (15 min)]
\noindent Le groupe rédige \textbf{ensemble} l'accord officiel sur une feuille A4 (format imposé).
\begin{enumerate}
    \item \textbf{Titre :} \eng{PEACE AGREEMENT — [Date] — [Place (ex. Palace of the Fon of Bafut / GHS Buea / Quarter Head's Office)]}.
    \item \textbf{Parties :} \eng{Between [Name/Role A] (Party A) and [Name/Role B] (Party B)}.
    \item \textbf{Attendus (\eng{Whereas}) :} \eng{Whereas a dispute arose concerning [précis objet]...}
    \item \textbf{Décision (\eng{Ruling}) :} \eng{The Council rules as follows:} + clauses numérotées avec \eng{shall} (obligation) ou \eng{agree to} (engagement mutuel).
    \item \textbf{Engagement :} \eng{Both parties undertake to respect this agreement and abide by the ruling.}
    \item \textbf{Signatures :} \eng{Signed: [Chairperson] \quad [Party A] \quad [Party B] \quad [Mediators]}.
\end{enumerate}
\noindent \textbf{Auto-correction :} Vérifiez \eng{shall} (pas de \eng{to}), \eng{agree to + V}, \eng{abide by + N}, passif (\eng{is decided}), orthographe du vocabulaire clé.
\end{methode}

\begin{methode}[Phase 4 — Restitution et évaluation par les pairs (10 min)]
\begin{enumerate}
    \item Le \eng{Chairperson} de chaque groupe lit l'accord (1 min max).
    \item La classe écoute (\eng{Listening for gist}) et remplit une grille d'évaluation :
    \begin{center}
    \begin{tabular}{|c|c|c|c|}
    \hline
    \rowcolor{popblueL} \textbf{Groupe / Cas} & \textbf{Solution juste ? (Oui/Non)} & \textbf{Meilleur vocabulaire ? (Note/5)} & \textbf{Clarté de l'oral (Note/5)} \\
    \hline
    A & & & \\
    \hline
    B & & & \\
    \hline
    C & & & \\
    \hline
    \end{tabular}
    \end{center}
    \item L'enseignant proclame le \eng{``Most Fair and Well-Worded Agreement''} et fait un retour linguistique global.
\end{enumerate}
\end{methode}

\courssub{Ressources linguistiques mobilisées}

\begin{propriete}[Lexique thématique : \eng{Peace and Conflict Resolution}]
\begin{center}
\begin{tabularx}{\linewidth}{|X|X|X|}
\hline
\rowcolor{popblueL} \textbf{Nom / Verbe} & \textbf{Collocations utiles (Contexte)} & \textbf{Traduction} \\
\hline
\eng{dispute} (n/v) & \eng{to settle a dispute, a land dispute, a boundary dispute} & différend / contester \\
\hline
\eng{conflict} (n) / \eng{resolve} (v) & \eng{to resolve a conflict, conflict resolution, peaceful resolution} & conflit / résoudre \\
\hline
\eng{mediation} (n) / \eng{mediate} (v) & \eng{traditional mediation, to mediate between two parties} & médiation / médier \\
\hline
\eng{compromise} (n/v) & \eng{to reach a compromise, a fair compromise, to compromise on sth} & compromis / transiger \\
\hline
\eng{boundary / boundary line} (n) & \eng{to mark the boundary, a boundary dispute, boundary marker} & limite, borne \\
\hline
\eng{inheritance} (n) / \eng{inherit} (v) & \eng{inheritance law, to inherit land, rightful heir} & héritage / hériter \\
\hline
\eng{verdict / ruling} (n) & \eng{to deliver a verdict, a final ruling, to accept the ruling} & verdict, jugement \\
\hline
\eng{undertake} (v) & \eng{to undertake to do sth (formel), an undertaking} & s'engager à, engagement \\
\hline
\eng{abide by} (phr. v) & \eng{to abide by the decision / the rules / the agreement} & respecter, se conformer à \\
\hline
\eng{impartial / unbiased} (adj) & \eng{an impartial mediator, an unbiased ruling} & impartial \\
\hline
\end{tabularx}
\end{center}
\end{propriete}

\begin{propriete}[Grammaire : Modaux et structures pour l'accord formel (\eng{Formal Agreement Register})]
\begin{itemize}
    \item \textbf{\eng{Shall}} : Valeur juridique, obligation formelle contraignante (toutes personnes dans le style légal). \eng{``Each party \textbf{shall} share the costs.''} (Chacune des parties \textbf{devra} partager...). \textit{Attention : pas de 'to' après shall.}
    \item \textbf{\eng{Must / Have to}} : Obligation forte, contrainte externe ou règle. \eng{``The gate \textbf{must} remain unlocked.''} / \eng{``They \textbf{have to} attend the next meeting.''}
    \item \textbf{\eng{Should / Ought to}} : Conseil, recommandation morale, avis. \eng{``They \textbf{should} plant a new tree.''}
    \item \textbf{\eng{Agree} : trois constructions distinctes} :
    \begin{center}
    \begin{tabular}{|l|l|l|}
    \hline
    \rowcolor{popblueL} \textbf{Structure} & \textbf{Exemple} & \textbf{Traduction} \\
    \hline
    \eng{agree \textbf{that} + clause} & \eng{We agree \textbf{that} the tree is shared.} & Nous convenons \textbf{que}... \\
    \hline
    \eng{agree \textbf{to} + infinitif} & \eng{They agreed \textbf{to} share the harvest.} & Ils ont accepté \textbf{de} partager... \\
    \hline
    \eng{agree \textbf{on / upon} + N / V-ing} & \eng{They agreed \textbf{on} a date / \textbf{on} sharing.} & Ils se sont mis d'accord \textbf{sur} une date / le partage. \\
    \hline
    \end{tabular}
    \end{center}
    \item \textbf{Passif formel (\eng{Present Simple Passive})} : Focus sur l'objet/décision. \eng{``The land \textbf{is divided} equally.''} / \eng{``The fine \textbf{will be paid} by the offender.''}
    \item \textbf{\eng{Whereas}} : Conjonction d'attendu (style juridique), introduit les faits. \eng{``\textbf{Whereas} a dispute arose...''} = « Attendu qu'un différend est survenu... / Considérant que... »
\end{itemize}
\end{propriete}

\begin{attention}[Erreurs fréquentes des francophones — \eng{False Friends \& Traps}]
\begin{itemize}
    \item \textbf{\eng{Compromise} $\neq$ \eng{Compromis}} (souvent péjoratif en français = arrangement boiteux). En anglais, \eng{compromise} est \textbf{positif/neutre} : solution mutuellement acceptable. \textit{Correct :} \eng{``They reached a \textbf{fair compromise}.''} \textit{Évitez :} \eng{``It's a bad compromise''} (préférez \eng{``It's an unfair settlement''}).
    \item \textbf{\eng{Undertake} $\neq$ \eng{Entreprendre}} (créer une entreprise). \eng{Undertake to do} = \textbf{s'engager formellement à faire}. \eng{``I \textbf{undertake to} pay the fine.''} = « Je m'engage à payer l'amende. »
    \item \textbf{\eng{Abide by}} : La préposition \textbf{by} est obligatoire. \textit{Faux :} \eng{abide the rules.} \textit{Juste :} \eng{abide \textbf{by} the rules.}
    \item \textbf{\eng{Inheritance}} (héritage successoral : argent, terre) vs \textbf{\eng{Heritage}} (héritage culturel, patrimonial, ex. \eng{World Heritage Site}). Pour la succession : \eng{inheritance}.
    \item \textbf{\eng{Party}} (juridique : partie au contrat/procès) vs \textbf{\eng{Part}} (morceau, rôle). \eng{``Both \textbf{parties} signed the agreement.''}
    \item \textbf{Ordre des mots :} \eng{``We agree \textbf{that}...''} (pas \eng{*agree to that...} pour introduire une proposition subordonnée).
    \item \textbf{\eng{Advice} (nom) / \eng{Advise} (verbe)} : \eng{``The Elders gave good \textbf{advice}.''} / \eng{``I \textbf{advise} you to accept.''}
\end{itemize}
\end{attention}

\courssub{Modèle de production et corrigé commenté}
\noindent Production attendue pour le \textbf{Cas A (The Kola Nut Tree Dispute)}.

\begin{exoresolu}[Modèle de \eng{Peace Agreement} — Cas A]
\textbf{Énoncé :} Rédiger l'accord de paix issu de la médiation du Conseil du Fon pour le litige de l'arbre à kola (Cas A).

\tcblower

\textbf{Solution détaillée :}

\begin{textebox}[PEACE AGREEMENT — 14th March 2025 — Palace of the Fon of Bafut]
\eng{Between Mr. Tanga (Party A) and Mr. Ngwa (Party B).}

\eng{Whereas a dispute arose concerning the ownership of the kola nut tree standing on the boundary line of their compounds.}

\eng{The Council, after hearing both parties and consulting customary land tenure rules, rules as follows:}

\eng{1. The tree belongs to both families as a shared resource.}

\eng{2. Each party \textbf{shall} harvest the nuts on the branches hanging over their own compound.}

\eng{3. Mr. Tanga \textbf{undertakes to} apologise to Mr. Ngwa's wife for the push.}

\eng{4. Both parties \textbf{agree to} maintain the tree and \textbf{abide by} this ruling.}

\eng{Signed: The Fon (Chairperson) \quad Mr. Tanga \quad Mr. Ngwa \quad Two Elders}
\end{textebox}

\textbf{Commentaire linguistique des choix opérés :}
\begin{enumerate}
    \item \textbf{Registre formel (\eng{Formal/Legal Register}) :} \eng{Whereas} (attendu que), \eng{rules as follows} (décide comme suit), \eng{undertakes to} (s'engage à — plus formel que \eng{promises to}), \eng{customary land tenure rules} (règles coutumières de tenure foncière).
    \item \textbf{\eng{Shall} (obligation partagée) :} \eng{``Each party shall harvest...''} impose une règle de droit coutumier contraignante pour les deux parties. Marque du \eng{Legal English}.
    \item \textbf{\eng{Agree to + V} :} \eng{``agree to maintain''} (action future convenue). Structure correcte (pas \eng{*agree maintaining}).
    \item \textbf{\eng{Abide by} + N :} Correctement construit. Objet : \eng{this ruling} (cette décision).
    \item \textbf{Participes adjectivaux :} \eng{branches hanging over} (branches qui surplombent / qui pendent au-dessus).
    \item \textbf{Mise en page juridique :} Deux-points après le titre, points fermes pour les articles numérotés, alignement des signatures.
\end{enumerate}
\end{exoresolu}

\courssub{Activité d'approfondissement (Devoir / Travail personnel)}

\begin{exercice}[Writing: The other cases]
\noindent Rédigez, \textbf{individuellement}, le \eng{Peace Agreement} pour le \textbf{Cas B} (School Football) \textbf{OU} le \textbf{Cas C} (Family Inheritance). Respectez scrupuleusement le format imposé (6 lignes minimum). Utilisez au moins 5 mots du tableau de vocabulaire (\eng{dispute, mediate, compromise, ruling, undertake, abide by, inheritance, boundary, unbiased, heir}).
\end{exercice}

\begin{exercice}[Grammar Transformation]
\noindent Reformulez les phrases suivantes en utilisant la structure indiquée entre parenthèses.
\begin{enumerate}
    \item The Council decided that the land would be shared. (\eng{It was decided that...} — Passif)
    \item "I promise I will pay the fine," said the offender. (\eng{The offender undertook...} — \eng{undertake to})
    \item You must respect the decision. (\eng{You have to...} / \eng{You are to...})
    \item They agreed on the date of the next meeting. (\eng{They agreed...} — \eng{that}-clause)
    \item The mediator was fair. He listened to both sides. (\eng{The mediator listened... unbiasedly} — Adverbe)
\end{enumerate}
\end{exercice}

\begin{corrige}[Exercice Grammar Transformation]
\begin{enumerate}
    \item \eng{It was decided that the land would be shared (by the Council).}
    \item \eng{The offender undertook to pay the fine.}
    \item \eng{You have to abide by the decision.} / \eng{You are to abide by the decision.}
    \item \eng{They agreed that the next meeting would be held on [date].}
    \item \eng{The mediator listened to both sides unbiasedly.} (ou \eng{in an unbiased manner})
\end{enumerate}
\end{corrige}

\courssub{Bilan et Auto-évaluation}

\begin{aretenir}[Bilan de la leçon — Lesson 7 : Settling Community Disputes]
\noindent Cette leçon a permis de travailler la compétence \textbf{Reading} (stratégies \emph{skimming, scanning, guessing}) sur des textes ancrés dans le réel camerounais (chefferie Bafut, école à Buea, famille à Bamenda). Le lexique de la paix et du conflit (\eng{dispute, mediation, compromise, inheritance, boundary, ruling, undertake, abide by}) est passé de la reconnaissance à la production active (oral \rightarrow écrit). La grammaire (\eng{shall, modals, agree structures, passive, whereas}) a été contextualisée dans la rédaction d'un acte juridique simulé (\eng{Peace Agreement}), préparant aux exigences du Baccalauréat (Composition, Guided Writing).
\end{aretenir}

\begin{aretenir}[Auto-évaluation de l'élève (à coller dans le cahier)]
\noindent À la fin de cette leçon, je suis capable de :
\begin{itemize}
    \item[\eng{$\square$}] Lire un texte narratif sur un conflit communautaire et en extraire l'idée générale (\emph{skimming}).
    \item[\eng{$\square$}] Repérer des informations précises (noms, lieux, actions, conséquences) par \emph{scanning}.
    \item[\eng{$\square$}] Deviner le sens de 4 mots inconnus grâce au contexte (\eng{compound, quarter head, heir, Njangi, mediate...}).
    \item[\eng{$\square$}] Participer à un jeu de rôle de médiation en utilisant \eng{should, must, shall, agree to/that/on}.
    \item[\eng{$\square$}] Rédiger un \eng{Peace Agreement} formel (5-6 lignes) avec le vocabulaire précis (\eng{Whereas, ruling, undertake to, abide by, shared resource}).
    \item[\eng{$\square$}] Expliquer en anglais pourquoi la solution proposée est \eng{fair} (juste) et \eng{unbiased} (impartiale).
\end{itemize}
\end{aretenir}

\begin{savaistu}[Culture : Justice coutumière au Cameroun (\eng{Customary Law in Cameroon})]
\noindent Au Cameroun, le système juridique est \textbf{bijuridique} : droit moderne (hérité du droit français/anglais) et droit coutumier (\eng{Customary Law}) coexistent. L'\eng{Administration of Justice} reconnaît les \eng{Traditional Courts} (Tribunaux coutumiers) présidés par les \eng{Fons}, \eng{Chiefs} ou \eng{Quarter Heads} pour les affaires de droit personnel (mariage, succession, terre) tant qu'elles ne sont pas \eng{repugnant to natural justice, equity and good conscience} (contraires à la justice naturelle, l'équité et la bonne conscience). Le \eng{Njangi} (ou \eng{Njangi house}) est à la fois une institution financière (tontine) et un cadre de régulation sociale et de médiation familiale très puissant dans les \eng{Grassfields} (Nord-Ouest, Ouest). L'arbre à kola (\eng{kola nut tree}) a une valeur symbolique et économique majeure : il marque souvent les limites foncières et ses fruits sont source de revenus.
\end{savaistu}

\newpage

\coursec{Méthodes et exercices résolus}

\coursec{Lesson 7 — Reading: Settling Community Disputes}

\courssub{Warm-up : Community Palaver}
\begin{experience}[Discussion à chaud — 10 min]
\begin{enumerate}
\item \textbf{Brainstorming.} Au tableau, écris le mot \eng{dispute}. Demande aux élèves de donner tous les mots qui leur viennent à l'esprit en anglais (ex. \eng{quarrel, fight, argument, land, family, chief, peace, agreement}). Classe-les en deux colonnes : \eng{Conflict} / \eng{Resolution}.
\item \textbf{Mise en situation orale.} En binômes, les élèves imaginent une scène : « Two neighbours argue over a mango tree. A third person intervenes. » Ils jouent le dialogue 1 minute, puis 2-3 binômes présentent devant la classe.
\item \textbf{Transition.} « Today we will read texts about how communities in Cameroon and elsewhere settle disputes peacefully. We will learn strategies to understand such texts and vocabulary to talk about peace and conflict. »
\end{enumerate}
\end{experience}

\courssub{Reading Text : Traditional Conflict Resolution in the Grassfields}
\begin{textebox}[Texte support — Lecture cursive]
“In the Grassfields region of Cameroon, traditional institutions play a vital role in maintaining social harmony. When a dispute arises between two quarters over land boundaries or water rights, the fon — the traditional ruler — summons the quarter heads and notables to a palaver at the palace. This gathering is not a modern court: there are no lawyers, no written codes, only the wisdom of the elders and the respect for custom.

First, each party states its grievances without interruption. The fon listens carefully, then asks the elders to mediate. They investigate the history of the land, consult the ancestors through libation, and propose a compromise. Often, the solution involves sharing the resource or alternating its use. The process ends with a ceremony of reconciliation: the parties shake hands, share a meal, and swear an oath to live in peace. This traditional mechanism restores relationships rather than merely punishing offenders, and the whole community witnesses the return to harmony.”
\end{textebox}

\begin{definition}[Glossaire / Vocabulaire clé]
\begin{itemize}
\item \eng{Grassfields} (n.) : région des Hauts-Plateaux (Nord-Ouest, Ouest Cameroun).
\item \eng{quarter} (n.) : quartier, subdivision d'un village ou d'une chefferie.
\item \eng{fon} (n.) : chef traditionnel, roi (dans les chefferies des Grassfields).
\item \eng{notable} (n.) : notable, personne respectée pour sa sagesse ou son statut.
\item \eng{palaver} (n.) : palabre, assemblée traditionnelle de discussion et de règlement de conflits.
\item \eng{grievance} (n.) : grief, plainte, motif de mécontentement.
\item \eng{libation} (n.) : libation, offrande liquide (vin de palme, eau) aux ancêtres.
\item \eng{compromise} (n.) : compromis, accord où chaque partie cède un peu.
\item \eng{reconciliation} (n.) : réconciliation (attention : un seul \emph{l}).
\item \eng{binding} (adj.) : contraignant, qui oblige (ex. \eng{a binding decision}).
\end{itemize}
\end{definition}

\courssub{Méthode 1 : Lire et comprendre un texte sur les conflits communautaires (skimming et scanning)}
\begin{methode}[Réussir la compréhension d'un texte sur \eng{settling community disputes}]
\begin{enumerate}
\item \textbf{Lis d'abord le titre et la première phrase de chaque paragraphe} (\cle{skimming}) : cela te donne le thème général (qui ? où ? quel conflit ? quelle solution ?) en moins d'une minute.
\item \textbf{Lis ensuite les questions avant le texte complet} : tu sauras quelles informations chercher (dates, noms, raisons, résultats).
\item \textbf{Relis le texte en cherchant les mots-clés des questions} (\cle{scanning}) : ne traduis pas tout, repère les noms propres, les chiffres, les verbes d'action.
\item \textbf{Repère les marqueurs logiques} : \eng{because} (cause), \eng{so / therefore} (conséquence), \eng{but / however} (opposition), \eng{finally} (étape finale d'un règlement de conflit).
\item \textbf{Rédige tes réponses en phrases complètes} en anglais, en reprenant les mots de la question, sans recopier de longs passages du texte.
\end{enumerate}
\end{methode}

\begin{exoresolu}[Compréhension de texte — type examen]
\textbf{Read the text and answer the questions.}

\emph{“In the village of Mankon, two families had quarrelled for years over a piece of farmland. The dispute became so serious that neighbours could no longer work peacefully. Finally, the fon called a palaver under the big tree. After listening to both sides, the elders decided that the land would be shared equally, and the two families shook hands in front of the whole community.”}

\textbf{Questions:} (a) What was the cause of the dispute? (b) Who helped to settle it? (c) What was the final decision?
\tcblower
\textbf{Étape 1 — Skimming.} Le titre implicite et la première phrase annoncent : un village (Mankon), deux familles, un conflit de terre. Thème : \eng{settling a land dispute}.

\textbf{Étape 2 — Scanning.} Je repère : \eng{quarrelled... over a piece of farmland} (cause), \eng{the fon called a palaver}, \eng{the elders decided} (médiateurs), \eng{the land would be shared equally} (décision).

\textbf{Étape 3 — Réponses rédigées en phrases complètes :}

(a) \eng{The dispute was caused by a disagreement between two families over a piece of farmland.} (« Le conflit venait d'un désaccord entre deux familles à propos d'un terrain. ») — Je reprends le mot \eng{caused} de la question et je ne recopie pas tout le texte.

(b) \eng{The fon and the village elders helped to settle the dispute by calling a palaver.} — Attention à la préposition : \eng{help to settle}, pas « help for settle ».

(c) \eng{The elders decided that the land would be shared equally between the two families.} — Discours rapporté au passé : \eng{would be shared} (passif, futur dans le passé).

\textbf{Conclusion :} chaque réponse est courte, exacte, et reprend les mots de la question. C'est ce que le correcteur attend.
\end{exoresolu}

\courssub{Méthode 2 : Deviner le sens des mots inconnus}
\begin{methode}[Déduire le sens d'un mot sans dictionnaire]
\begin{enumerate}
\item \textbf{Identifie la nature du mot} : nom, verbe, adjectif ? La terminaison aide : \eng{-tion}, \eng{-ment} (noms) ; \eng{-ful}, \eng{-ous} (adjectifs) ; \eng{-ise/-ize} (verbes).
\item \textbf{Cherche les indices dans la phrase} : synonyme (\eng{or}, \eng{that is}), contraire (\eng{but}, \eng{unlike}), exemple (\eng{such as}).
\item \textbf{Utilise le contexte du thème} : dans un texte sur les conflits, \eng{dispute}, \eng{quarrel}, \eng{clash} seront proches de « conflit » ; \eng{settle}, \eng{resolve}, \eng{mediate} proches de « régler ».
\item \textbf{Méfie-toi des faux amis} : \eng{actually} = « en fait » (pas « actuellement ») ; \eng{sensible} = « raisonnable » (pas « sensible »).
\item \textbf{Vérifie ton hypothèse} : remplace le mot par ta traduction et relis la phrase. Si elle a du sens, c'est bon.
\end{enumerate}
\end{methode}

\begin{exoresolu}[Vocabulaire en contexte — type examen]
\textbf{Find the meaning of the underlined words from the context.}

\eng{“The two neighbours were involved in a bitter \textbf{dispute} over a mango tree, but the village chief acted as a \textbf{mediator} and helped them reach an \textbf{amicable} agreement, so they did not need to go to court.”}
\tcblower
\textbf{dispute} : nom (précédé de l'adjectif \eng{bitter}). Contexte : deux voisins, un arbre, un désaccord. Indice : \eng{over a mango tree}. Sens : « conflit, différend ». Vérification : « un conflit au sujet d'un manguier » — la phrase a du sens.

\textbf{mediator} : nom de personne (suffixe \eng{-or}, comme \eng{actor}). Contexte : \eng{the village chief acted as...} et \eng{helped them reach an agreement}. Sens : « médiateur », personne qui aide deux parties à s'entendre.

\textbf{amicable} : adjectif (avant le nom \eng{agreement}). Indice : \eng{so they did not need to go to court} — l'accord est obtenu sans tribunal, donc à l'amiable. Sens : « amical, à l'amiable ».

\textbf{Réponses rédigées :} \eng{A dispute is a serious disagreement or quarrel. A mediator is a person who helps two sides reach an agreement. Amicable means friendly, settled without fighting or going to court.}

\textbf{Conclusion :} on ne traduit jamais un mot isolé : on s'appuie toujours sur la nature grammaticale et les indices de la phrase.
\end{exoresolu}

\courssub{Méthode 3 : Répondre aux questions de compréhension et d'interprétation}
\begin{methode}[Questions factuelles et questions d'interprétation]
\begin{enumerate}
\item \textbf{Distingue les deux types de questions} : les questions factuelles (\eng{What happened? When? Who?}) ont leur réponse dans le texte ; les questions d'interprétation (\eng{Why do you think...? What can we learn?}) demandent ton raisonnement appuyé sur le texte.
\item \textbf{Pour les questions factuelles} : cite ou paraphrase le passage précis, avec la ligne ou le paragraphe si demandé (\eng{According to paragraph 2...}).
\item \textbf{Pour les questions d'interprétation} : commence par \eng{I think that... / In my opinion... / This shows that...} puis justifie avec un élément du texte (\eng{because the text says that...}).
\item \textbf{Pour les questions vrai/faux} : justifie toujours par une courte citation du texte, sinon la réponse ne vaut rien.
\item \textbf{Respecte le temps verbal de la question} : question au past simple, réponse au past simple.
\end{enumerate}
\end{methode}

\begin{exoresolu}[Questions factuelles et d'interprétation — type examen]
\textbf{Texte :} \eng{“When the water pump in Nkambe broke down, the villagers first blamed one another. Then a youth leader suggested a community meeting. Everyone contributed money and labour, and within a week the pump was repaired. The villagers learnt that dialogue is more useful than accusations.”}

\textbf{Questions:} (a) True or False: the villagers repaired the pump alone, without meeting. Justify. (b) Why do you think the meeting was successful?
\tcblower
\textbf{(a)} Question vrai/faux : je cherche le passage correspondant. Le texte dit : \eng{a youth leader suggested a community meeting} et \eng{everyone contributed money and labour}. La phrase de l'énoncé est donc fausse.

Réponse rédigée : \eng{False. The text says that a youth leader suggested a community meeting and that everyone contributed money and labour to repair the pump.} (« Faux. Le texte dit qu'un jeune leader a proposé une réunion communautaire et que tout le monde a contribué en argent et en travail. ») — Justification obligatoire par citation.

\textbf{(b)} Question d'interprétation : je donne mon avis + justification textuelle.

Réponse rédigée : \eng{I think the meeting was successful because it allowed the villagers to stop blaming one another and to work together. The text shows that when everyone contributed, the pump was repaired within a week.} — Structure : \eng{I think... because...} + preuve tirée du texte.

\textbf{Conclusion :} une réponse d'interprétation sans appui sur le texte est hors sujet ; une réponse factuelle sans justification perd la moitié des points.
\end{exoresolu}

\courssub{Méthode 4 : Maîtriser le lexique de la paix et du conflit}
\begin{methode}[Apprendre et réutiliser le vocabulaire thématique]
\begin{enumerate}
\item \textbf{Classe le vocabulaire en deux colonnes} : mots du conflit (\eng{dispute, quarrel, clash, tension, grievance, blame}) et mots de la paix (\eng{settlement, dialogue, reconciliation, compromise, mediation, peace}).
\item \textbf{Apprends les verbes avec leurs prépositions} : \eng{to quarrel with somebody about something}, \eng{to settle a dispute}, \eng{to apologise for something}, \eng{to agree on a solution}, \eng{to listen to both sides}.
\item \textbf{Apprends les familles de mots} : \eng{to reconcile} (verbe) → \eng{reconciliation} (nom) ; \eng{to mediate} → \eng{mediation / mediator} ; \eng{peaceful} (adjectif) → \eng{peacefully} (adverbe).
\item \textbf{Réutilise au moins cinq mots du thème} dans chaque production écrite ou orale : c'est un critère de notation explicite.
\end{enumerate}
\end{methode}

\begin{exoresolu}[Vocabulaire — type examen]
\textbf{Complete the sentences with the correct word from the box:} \eng{dispute — mediator — apologise — reconciliation — compromise}

1. The two villages signed a peace agreement after years of \_\_\_\_\_\_\_\_ over the river.
2. A respected elder served as a \_\_\_\_\_\_\_\_ between the two families.
3. Neither side won everything: they finally reached a \_\_\_\_\_\_\_\_.
4. The young man decided to \_\_\_\_\_\_\_\_ for his rude behaviour.
5. The ceremony marked the \_\_\_\_\_\_\_\_ of the two communities.
\tcblower
1. \eng{dispute} — nom après la préposition \eng{after years of} ; sens : « un différend au sujet de la rivière ». (\eng{dispute over something} : bonne préposition.)

2. \eng{mediator} — article indéfini \eng{a} + nom de personne ; le suffixe \eng{-or} désigne celui qui fait l'action de \eng{mediate}. Sens : « médiateur ».

3. \eng{compromise} — \eng{reached a compromise} : collocation à connaître par cœur (« parvenir à un compromis »). Indice : \eng{neither side won everything}.

4. \eng{apologise} — après \eng{decided to}, il faut la base verbale (infinitif sans \eng{to}). \eng{Apologise for something} = « s'excuser de quelque chose ». Attention : pas « excuse for ».

5. \eng{reconciliation} — nom après l'article \eng{the} ; famille de \eng{to reconcile}. Sens : « réconciliation ».

\textbf{Conclusion :} pour ce type d'exercice, vérifie toujours deux choses : la nature grammaticale exigée par la phrase (nom, verbe, adjectif) et la préposition qui accompagne le mot.
\end{exoresolu}

\courssub{Méthode 5 : Résumer un texte et donner son opinion}
\begin{methode}[Rédiger un résumé suivi d'une opinion personnelle]
\begin{enumerate}
\item \textbf{Repère les quatre éléments essentiels} : qui (les personnages), quoi (le conflit), comment (la médiation), résultat (la solution).
\item \textbf{Rédige le résumé en trois ou quatre phrases} avec tes propres mots : commence par \eng{The text is about... / The passage deals with...}. N'utilise ni citations longues ni détails secondaires (dates, descriptions).
\item \textbf{Utilise les connecteurs} : \eng{first, then, after that, finally} pour la chronologie du règlement du conflit.
\item \textbf{Ajoute ton opinion en deux ou trois phrases} : \eng{In my opinion, this way of settling disputes is... because... / I believe that dialogue is better than violence because...}
\item \textbf{Relis-toi} : accords sujet-verbe, temps (past simple pour le résumé, present simple pour l'opinion générale), un mot du thème par phrase.
\end{enumerate}
\end{methode}

\begin{exoresolu}[Résumé et opinion — type composition au BEPC/Probatoire]
\textbf{Summarise the following text in about four sentences, then give your opinion on traditional ways of settling disputes (two or three sentences).}

\emph{“In Bafut, a serious conflict opposed two quarters over the use of a community field. Insults were exchanged and young people nearly fought. The fon invited the notables of both quarters to the palace. Each side explained its grievances. After a long discussion, the notables agreed on a calendar for sharing the field, and a ceremony of reconciliation closed the palaver.”}
\tcblower
\textbf{Étape 1 — Éléments essentiels.} Qui : two quarters of Bafut. Quoi : a conflict over a community field. Comment : the fon and the notables held a palaver. Résultat : a sharing calendar and a reconciliation ceremony.

\textbf{Étape 2 — Résumé rédigé (past simple, connecteurs) :}

\eng{The text is about a serious dispute between two quarters of Bafut over a community field. First, the tension became so high that young people nearly fought. Then, the fon called the notables of both sides to the palace, and each side explained its grievances. Finally, they agreed on a calendar for sharing the field and celebrated their reconciliation.}

(« Le texte parle d'un grave conflit entre deux quartiers de Bafut à propos d'un champ communautaire. D'abord, la tension est devenue telle que des jeunes ont failli se battre. Ensuite, le fon a convoqué les notables des deux camps au palais, et chaque camp a exposé ses griefs. Enfin, ils se sont mis d'accord sur un calendrier de partage du champ et ont célébré leur réconciliation. »)

\textbf{Étape 3 — Opinion (present simple) :}

\eng{In my opinion, traditional ways of settling disputes are very effective because the whole community takes part in the solution. I believe that dialogue under the authority of elders restores peace more durably than a court decision.}

\textbf{Conclusion :} le résumé est au passé et sans détails inutiles ; l'opinion est au présent, introduite par \eng{In my opinion} et justifiée par \eng{because}. Deux parties clairement séparées = barème respecté.
\end{exoresolu}

\begin{attention}[Les 5 erreurs qui coûtent des points au Bac]
\begin{itemize}
\item \textbf{Faux amis.} \eng{Actually} signifie « en fait », pas « actuellement » (qui se dit \eng{currently, nowadays}). \eng{To assist} signifie « aider » ; « assister à une réunion » se dit \eng{to attend a meeting}.
\item \textbf{Prépositions.} On dit \eng{to listen to both sides}, \eng{to apologise for his behaviour}, \eng{to agree on a solution}, \eng{a dispute over land}. Les calques du français (« listen both sides », « excuse for ») sont systématiquement sanctionnés.
\item \textbf{Temps verbaux.} Le résumé d'un récit se fait au \cle{past simple} (\eng{the fon called, they agreed}), l'opinion générale au \cle{present simple} (\eng{dialogue is better than violence}). Ne mélange pas les deux dans la même partie.
\item \textbf{Orthographe des mots du thème.} \eng{reconciliation} (un seul \emph{l}), \eng{mediator} (avec \emph{-or}), \eng{peaceful} (un seul \emph{l}), \eng{agreement} (deux \emph{e}). Vérifie chaque mot du lexique en te relisant.
\item \textbf{Ordre des mots et questions.} Dans une réponse, sujet + verbe : \eng{The elders decided...}, jamais « Decided the elders... ». Et dans les questions d'interprétation, écris \eng{Why do you think...?} et non « Why you think...? » — l'auxiliaire \eng{do} est obligatoire.
\end{itemize}
\end{attention}

\newpage

\coursec{Exercices}

\courssub{Je vérifie mes connaissances}

\begin{exercice}[QCM : comprendre le texte]
Read the following short text, then choose the correct answer (a, b or c).

\begin{textebox}[A quarrel at the market]
Last Saturday, two traders at Mokolo market quarrelled over a piece of land where they both wanted to display their goods. The quarrel became so loud that customers ran away. The market master called a mediator, an old woman respected by everybody. She listened to both sides patiently. After two hours of discussion, the two traders reached an agreement: each of them would use the land on alternate days. They apologised to each other and shook hands. Peace returned to the market.
\end{textebox}

\begin{enumerate}
\item The two traders quarrelled over: \quad a) money \quad b) a piece of land \quad c) a customer
\item The mediator was: \quad a) a policeman \quad b) a young man \quad c) a respected old woman
\item The discussion lasted: \quad a) two hours \quad b) two days \quad c) twenty minutes
\item At the end, the two traders: \quad a) went to court \quad b) reached an agreement \quad c) left the market
\end{enumerate}
\end{exercice}

\begin{exercice}[Vrai ou faux, avec justification]
Read the text of the previous exercise again. Say if the following statements are \eng{true} or \eng{false} and justify each answer with a short quotation from the text.

\begin{enumerate}
\item The quarrel frightened the customers.
\item The mediator listened to only one trader.
\item The two traders will share the land.
\item The conflict ended with a fight.
\end{enumerate}
\end{exercice}

\begin{exercice}[Vocabulaire : déduire le sens des mots]
The following words are taken from the lesson text \emph{Peace Under the Mango Tree}. For each word, choose the correct meaning (a, b or c) using the context.

\begin{enumerate}
\item \eng{quarrelled} : \quad a) argued angrily \quad b) sang together \quad c) worked silently
\item \eng{gather} : \quad a) to run away \quad b) to come together \quad c) to sell goods
\item \eng{fine} : \quad a) a sum of money paid as a punishment \quad b) a gift \quad c) a meal
\item \eng{forgive} : \quad a) to punish \quad b) to forget a debt of money \quad c) to stop being angry with somebody
\item \eng{peacefully} : \quad a) with violence \quad b) without violence \quad c) very quickly
\end{enumerate}
\end{exercice}

\begin{exercice}[Grammaire : compléter avec le bon temps]
Complete the following sentences with the correct form of the verb in brackets (simple past or present perfect).

\begin{enumerate}
\item The two families \rule{1.5cm}{0.4pt} (quarrel) for many years before the chief intervened.
\item Yesterday, the elders \rule{1.5cm}{0.4pt} (gather) under the mango tree to discuss the problem.
\item The mediator \rule{1.5cm}{0.4pt} (just / arrive) in the village.
\item After the meeting, both sides \rule{1.5cm}{0.4pt} (shake) hands and \rule{1.5cm}{0.4pt} (sign) the agreement.
\item \rule{1.5cm}{0.4pt} you \rule{1.5cm}{0.4pt} (ever / witness) a traditional palaver?
\end{enumerate}
\end{exercice}

\courssub{Je m'entraîne}

\begin{exercice}[Traduction : français vers anglais]
Translate the following sentences into English, using the vocabulary of the lesson.

\begin{enumerate}
\item Les deux voisins se disputent souvent à cause du bruit.
\item Le chef du village a écouté les deux parties avec patience.
\item Après une longue discussion, ils sont arrivés à un accord.
\item Elle a présenté ses excuses et son voisin lui a pardonné.
\item Les anciens ont réglé le conflit pacifiquement, sans aller au tribunal.
\end{enumerate}
\end{exercice}

\begin{exercice}[Transformation de phrases]
Rewrite each sentence as indicated, without changing the meaning.

\begin{enumerate}
\item The mediator listened to both sides. (Change into the negative form.)
\item The two brothers reached an agreement. (Change into a question.)
\item The chief said: “You must forgive each other.” (Report the sentence, beginning with: \emph{The chief told them...})
\item They settled the dispute peacefully. (Change into the passive voice.)
\item Although they were angry, they shook hands. (Rewrite using \emph{but}.)
\end{enumerate}
\end{exercice}

\begin{exercice}[Texte à trous : le vocabulaire du conflit et de la paix]
Complete the following text with the words from the box. Use each word only once.

\begin{center}
\begin{tabular}{|c|c|c|c|c|}
\hline
\rowcolor{popblueL} dispute & settle & agreement & mediator & forgive \\
\hline
\end{tabular}
\end{center}

\begin{textebox}[The water tap conflict]
In our quarter, there was a serious \rule{2cm}{0.4pt} between two families about the public water tap. Nobody could \rule{2cm}{0.4pt} the problem, so the quarter head called a \rule{2cm}{0.4pt} from a neighbouring village. She spoke to each family separately and then brought them together. Finally, the two families signed an \rule{2cm}{0.4pt}: each family would fetch water at fixed hours. The father of the first family said he was ready to \rule{2cm}{0.4pt} his neighbour, and they became friends again.
\end{textebox}
\end{exercice}

\begin{exercice}[Appariement : expressions de la médiation]
Match each expression (1--6) with its French equivalent (a--f).

\begin{center}
\begin{tabularx}{\linewidth}{|X|X|}
\hline
\rowcolor{popblueL} English expression & French equivalent \\
\hline
1. to listen to both sides & a) présenter ses excuses \\
\hline
2. to reach an agreement & b) serrer la main \\
\hline
3. to apologise & c) écouter les deux parties \\
\hline
4. to shake hands & d) régler un conflit \\
\hline
5. to settle a dispute & e) parvenir à un accord \\
\hline
6. to forgive somebody & f) pardonner à quelqu'un \\
\hline
\end{tabularx}
\end{center}
\end{exercice}

\courssub{Je me prépare au Bac}

\begin{exercice}[Type Bac : compréhension de l'écrit]
Read the following text carefully and answer the questions.

\begin{textebox}[Peace Under the Mango Tree]
In the village of Mbankomo, not far from Yaoundé, two families, the Etoundis and the Nkomos, quarrelled for almost ten years over a small piece of farmland. Every planting season, the quarrel started again: insults, threats, and even broken tools. The children of the two families were forbidden to play together, and the whole village suffered from this climate of tension.

One day, the village chief decided that enough was enough. He invited the elders to gather under the old mango tree in the centre of the village, the traditional place for palavers. A respected mediator, Mama Bébé, was asked to lead the discussion. She first listened to both sides separately, taking notes in a small notebook. Then she brought the two families together and asked each of them to explain their grievances without interruption.

After a long afternoon of discussion, the truth appeared: the boundary between the two farms had never been clearly marked. The elders decided to plant a row of palm trees to separate the lands. To seal the agreement, the head of each family paid a small symbolic fine of one goat, which was shared by the whole village during a reconciliation feast. The two men apologised, forgave each other and shook hands in front of everybody.

Today, the children of the two families play football together under the mango tree, and the palm trees remind everyone that dialogue is stronger than anger.
\end{textebox}

\textbf{Questions}

\begin{enumerate}
\item \textbf{True or false?} Justify each answer with a quotation from the text. (2 marks)
\begin{enumerate}
\item The two families quarrelled for more than five years.
\item The boundary between the farms was clearly marked.
\end{enumerate}
\item \textbf{Multiple choice.} Choose the correct answer. (1 mark)\\
Mama Bébé led the discussion because she was:
\begin{enumerate}
\item the richest woman in the village
\item a respected mediator
\item the mother of the two families
\end{enumerate}
\item \textbf{Short answers.} (2 marks)
\begin{enumerate}
\item Where did the elders gather for the palaver?
\item What did the elders plant to separate the two farms?
\end{enumerate}
\item \textbf{Vocabulary.} Find in the text words or expressions meaning: (2 marks)
\begin{enumerate}
\item argued angrily (paragraph 1)
\item complaints (paragraph 2)
\item to stop being angry with somebody (paragraph 3)
\end{enumerate}
\item \textbf{Language.} Rewrite the following sentence in the passive voice: “The elders decided to plant a row of palm trees.” (1 mark)
\end{enumerate}
\end{exercice}

\begin{exercice}[Type Bac : traduction]
Translate the following passage into French. (5 marks)

\begin{textebox}[For translation]
One day, the village chief decided that enough was enough. He invited the elders to gather under the old mango tree in the centre of the village, the traditional place for palavers. A respected mediator, Mama Bébé, was asked to lead the discussion. She first listened to both sides separately, taking notes in a small notebook.
\end{textebox}
\end{exercice}

\begin{exercice}[Type Bac : expression écrite]
Choose \textbf{one} of the following topics and write your composition.

\begin{enumerate}
\item \textbf{Topic 1 (guided paragraph, 50--60 words).} Your two neighbours are always quarrelling. Write a short paragraph explaining the problem and proposing a peaceful solution. Use at least three expressions from the lesson: \eng{settle a dispute}, \eng{listen to both sides}, \eng{reach an agreement}, \eng{apologise}, \eng{forgive}.
\item \textbf{Topic 2 (composition, 150--180 words).} Describe a conflict that took place in your community (in your family, your quarter or your school) and explain how it was settled. Say who acted as mediator, what solution was found, and give your opinion on traditional ways of settling disputes compared to going to court. Use link words: \eng{first}, \eng{then}, \eng{after that}, \eng{finally}.
\end{enumerate}
\end{exercice}

\newpage

\coursec{Corrigés des exercices}

\courssub{Je vérifie mes connaissances}

\begin{corrige}[Exercice 1 — QCM : comprendre le texte]
\begin{enumerate}
\item \textbf{b) a piece of land.} Le texte dit : \eng{“two traders at Mokolo market quarrelled over a piece of land”}.
\item \textbf{c) a respected old woman.} Le texte dit : \eng{“a mediator, an old woman respected by everybody”}.
\item \textbf{a) two hours.} Le texte dit : \eng{“After two hours of discussion”}.
\item \textbf{b) reached an agreement.} Le texte dit : \eng{“the two traders reached an agreement”}.
\end{enumerate}
\textbf{Méthode :} il s'agit ici d'un exercice de \cle{scanning} : on repère dans le texte les mots-clés de la question (\eng{quarrelled over}, \eng{mediator}, \eng{discussion}, \eng{at the end}) et on lit uniquement la phrase qui les contient.
\end{corrige}

\begin{corrige}[Exercice 2 — Vrai ou faux, avec justification]
\begin{enumerate}
\item \textbf{True.} \eng{“The quarrel became so loud that customers ran away.”} (Les clients se sont enfuis : la querelle les a effrayés.)
\item \textbf{False.} \eng{“She listened to both sides patiently.”} (Elle a écouté \emph{les deux} parties, pas une seule.)
\item \textbf{True.} \eng{“each of them would use the land on alternate days.”} (Chacun utilisera le terrain un jour sur deux : ils le partagent.)
\item \textbf{False.} \eng{“They apologised to each other and shook hands.”} (Le conflit s'est terminé par des excuses et une poignée de main, pas par une bagarre.)
\end{enumerate}
\textbf{Méthode :} une justification n'est valable que si elle cite le texte entre guillemets. Une réponse \eng{true}/\eng{false} sans citation ne rapporte pas de point au Bac.
\end{corrige}

\begin{corrige}[Exercice 3 — Vocabulaire : déduire le sens des mots]
\begin{enumerate}
\item \textbf{a) argued angrily.} \eng{quarrel} = se quereller, se disputer. Prétérit régulier : \eng{quarrelled} (doublement du \eng{l} à l'orthographe britannique).
\item \textbf{b) to come together.} \eng{gather} = se rassembler, se réunir.
\item \textbf{a) a sum of money paid as a punishment.} \eng{a fine} = une amende. \textbf{Attention au faux ami :} \eng{fine} ne signifie pas « fin » (la fin = \eng{the end}) ni « fin » au sens de « mince » (\eng{thin}).
\item \textbf{c) to stop being angry with somebody.} \eng{forgive} = pardonner (verbe irrégulier : \eng{forgive – forgave – forgiven}).
\item \textbf{b) without violence.} \eng{peacefully} = pacifiquement (adverbe formé sur \eng{peaceful} + \eng{-ly}).
\end{enumerate}
\textbf{Méthode :} pour deviner le sens d'un mot, on observe le contexte (les mots qui l'entourent) et la famille du mot : \eng{peace} (paix) $\rightarrow$ \eng{peaceful} $\rightarrow$ \eng{peacefully}.
\end{corrige}

\begin{corrige}[Exercice 4 — Grammaire : compléter avec le bon temps]
\begin{enumerate}
\item \eng{The two families \textbf{quarrelled} for many years before the chief intervened.} — Simple past : action révolue, située avant une autre action passée (\eng{intervened}).
\item \eng{Yesterday, the elders \textbf{gathered} under the mango tree to discuss the problem.} — Simple past obligatoire avec \eng{yesterday}, marqueur de temps précis et terminé.
\item \eng{The mediator \textbf{has just arrived} in the village.} — Present perfect avec \eng{just} : action passée récente, sans date précise, avec un lien avec le présent.
\item \eng{After the meeting, both sides \textbf{shook} hands and \textbf{signed} the agreement.} — Simple past : succession d'actions terminées. \eng{shake} est irrégulier : \eng{shake – shook – shaken}.
\item \eng{\textbf{Have} you \textbf{ever witnessed} a traditional palaver?} — Present perfect avec \eng{ever} : expérience de vie, sans date.
\end{enumerate}
\textbf{Règle à retenir :} un marqueur de temps précis et terminé (\eng{yesterday}, \eng{last week}, \eng{in 2010}) impose le simple past ; \eng{just}, \eng{ever}, \eng{never}, \eng{already}, \eng{yet} appellent le present perfect.
\end{corrige}

\courssub{Je m'entraîne}

\begin{corrige}[Exercice 5 — Traduction : français vers anglais]
\begin{enumerate}
\item \eng{The two neighbours often quarrel about the noise.} — « Souvent » = \eng{often}, placé avant le verbe ; présent simple pour une habitude.
\item \eng{The village chief listened to both sides patiently.} — Prétérit ; \eng{listen to} (préposition \eng{to} obligatoire) ; « les deux parties » = \eng{both sides}.
\item \eng{After a long discussion, they reached an agreement.} — « Parvenir à un accord » = \eng{reach an agreement} ; ne jamais dire \eng{*arrive to an agreement}.
\item \eng{She apologised and her neighbour forgave her.} — \eng{forgive} est irrégulier : prétérit \eng{forgave} ; « présenter ses excuses » = \eng{apologise}.
\item \eng{The elders settled the conflict peacefully, without going to court.} — « Régler un conflit » = \eng{settle a conflict/a dispute} ; « aller au tribunal » = \eng{go to court} (sans article).
\end{enumerate}
\textbf{Points de vigilance :} \eng{apologise} se construit sans complément direct (\eng{she apologised}, ou \eng{she apologised to him}) ; \eng{forgive} prend un complément de personne (\eng{forgive somebody}).
\end{corrige}

\begin{corrige}[Exercice 6 — Transformation de phrases]
\begin{enumerate}
\item \eng{The mediator did not (didn't) listen to both sides.} — Négation au simple past : \eng{did not} + base verbale (le verbe perd sa marque de prétérit).
\item \eng{Did the two brothers reach an agreement?} — Question au simple past : \eng{Did} + sujet + base verbale (\eng{reached} $\rightarrow$ \eng{reach}).
\item \eng{The chief told them (that) they had to forgive each other.} — Discours rapporté : \eng{must} devient \eng{had to} ; le présent recule au passé ; on supprime les guillemets.
\item \eng{The dispute was settled peacefully (by them).} — Passif au simple past : \eng{was} + participe passé \eng{settled} ; l'ancien COD (\eng{the dispute}) devient sujet.
\item \eng{They were angry, but they shook hands.} — \eng{although} (subordonnée) est remplacé par \eng{but} (coordination) ; on ne peut jamais cumuler les deux dans la même phrase.
\end{enumerate}
\textbf{Attention :} à la question 2, l'erreur classique est d'écrire \eng{*Did they reached...} ; après \eng{did}, le verbe revient toujours à la base verbale.
\end{corrige}

\begin{corrige}[Exercice 7 — Texte à trous : le vocabulaire du conflit et de la paix]
\begin{enumerate}
\item \eng{dispute} : \eng{“there was a serious \textbf{dispute} between two families”} (nom, précédé de l'article \eng{a} et de l'adjectif \eng{serious}).
\item \eng{settle} : \eng{“Nobody could \textbf{settle} the problem”} (base verbale après le modal \eng{could}).
\item \eng{mediator} : \eng{“the quarter head called a \textbf{mediator} from a neighbouring village”} (personne qui intervient ; la phrase suivante commence par \eng{She}, ce qui confirme qu'il s'agit d'une personne).
\item \eng{agreement} : \eng{“the two families signed an \textbf{agreement}”} (nom, précédé de \eng{an} — voyelle initiale — et du verbe \eng{signed}).
\item \eng{forgive} : \eng{“he was ready to \textbf{forgive} his neighbour”} (base verbale après \eng{to}).
\end{enumerate}
\textbf{Méthode :} avant de remplir, identifier la nature du mot attendu (nom après un article, verbe après un modal ou après \eng{to}). Cela élimine immédiatement la moitié des choix.
\end{corrige}

\begin{corrige}[Exercice 8 — Appariement : expressions de la médiation]
\begin{center}
\begin{tabularx}{\linewidth}{|X|X|}
\hline
\rowcolor{popblueL} English expression & French equivalent \\
\hline
1. to listen to both sides & c) écouter les deux parties \\
\hline
2. to reach an agreement & e) parvenir à un accord \\
\hline
3. to apologise & a) présenter ses excuses \\
\hline
4. to shake hands & b) serrer la main \\
\hline
5. to settle a dispute & d) régler un conflit \\
\hline
6. to forgive somebody & f) pardonner à quelqu'un \\
\hline
\end{tabularx}
\end{center}
\textbf{Remarques :} \eng{shake hands} est irrégulier (\eng{shake – shook – shaken}) ; \eng{apologise} s'écrit aussi \eng{apologize} en anglais américain ; attention au faux ami : \eng{an apology} = des excuses, pas « une apologie ».
\end{corrige}

\courssub{Je me prépare au Bac}

\begin{corrige}[Exercice 9 — Type Bac : compréhension de l'écrit]
\textbf{Barème indicatif : 8 points.}

\begin{enumerate}
\item \textbf{True or false (2 marks : 0,5 pour la réponse + 0,5 pour la citation).}
\begin{enumerate}
\item \textbf{True.} \eng{“two families, the Etoundis and the Nkomos, quarrelled for almost ten years”} (presque dix ans, donc plus de cinq).
\item \textbf{False.} \eng{“the boundary between the two farms had never been clearly marked”} (la limite n'avait \emph{jamais} été clairement marquée).
\end{enumerate}
\item \textbf{Multiple choice (1 mark).} \textbf{b) a respected mediator.} Citation : \eng{“A respected mediator, Mama Bébé, was asked to lead the discussion.”}
\item \textbf{Short answers (2 marks : 1 par réponse).}
\begin{enumerate}
\item \eng{They gathered under the old mango tree in the centre of the village.}
\item \eng{They planted a row of palm trees (to separate the lands).}
\end{enumerate}
\item \textbf{Vocabulary (2 marks).}
\begin{enumerate}
\item \eng{quarrelled} (paragraphe 1) = argued angrily.
\item \eng{grievances} (paragraphe 2) = complaints.
\item \eng{forgave} (paragraphe 3, de \eng{forgive}) = to stop being angry with somebody.
\end{enumerate}
\item \textbf{Language (1 mark).} \eng{“A row of palm trees was decided to be planted (by the elders).”} — Formulation plus naturelle acceptée : \eng{“It was decided by the elders to plant a row of palm trees.”} Le passif exige \eng{was} + participe passé.
\end{enumerate}
\textbf{Points de vigilance :} répondre par des phrases complètes ; recopier les citations sans les modifier ; à la question 3, reprendre le sujet (\eng{They}) et accorder le temps (simple past).
\end{corrige}

\begin{corrige}[Exercice 10 — Type Bac : traduction]
\textbf{Barème indicatif : 5 points} (1 point par segment correct, moins les fautes de grammaire et de vocabulaire).

\textbf{Traduction proposée :}

« Un jour, le chef du village décida qu'il en avait assez (que c'en était trop). Il invita les anciens à se rassembler sous le vieux manguier au centre du village, le lieu traditionnel des palabres. Une médiatrice respectée, Mama Bébé, fut chargée de diriger (conduire) la discussion. Elle écouta d'abord les deux parties séparément, en prenant des notes dans un petit carnet. »

\textbf{Points de vigilance :}
\begin{itemize}
\item \eng{enough was enough} : expression idiomatique ; traduire « c'en était trop » ou « il en avait assez », jamais mot à mot.
\item \eng{the elders} : « les anciens » (du village), pas « les personnes âgées ».
\item \eng{was asked to lead} : passif anglais rendu par « fut chargée de » ou « on demanda à... de » — le passif lourd « fut demandée de » est à éviter.
\item \eng{taking notes} : gérondif rendu par « en prenant des notes ».
\item \eng{palavers} : « palabres » (mot admis en français dans ce contexte culturel).
\end{itemize}
\end{corrige}

\begin{corrige}[Exercice 11 — Type Bac : expression écrite]
\textbf{Barème indicatif :} respect de la consigne et de la longueur (2 pts), exactitude grammaticale (4 pts), vocabulaire du thème (2 pts), cohérence et connecteurs (2 pts).

\textbf{Topic 1 — modèle de paragraphe (environ 55 mots) :}

\begin{textebox}[Model answer — Topic 1]
My two neighbours always quarrel about noise and their animals. Last week, I decided to help them settle their dispute. I listened to both sides calmly, and then I proposed a solution: the animals must stay in a fence, and music must stop at nine o'clock. Finally, they reached an agreement, apologised to each other and forgave each other. Now our street is peaceful again.
\end{textebox}

\textbf{Topic 2 — modèle de composition (environ 165 mots) :}

\begin{textebox}[Model answer — Topic 2]
Last year, a serious conflict took place in my quarter in Yaoundé. Two families disagreed about a mango tree that stood between their compounds: both families wanted its fruits.

First, the quarrel started with insults every morning. Then, the children of the two families stopped talking to each other, and the situation became worse. After that, the quarter head decided to act as mediator. He invited the elders and the two families to a palaver on a Sunday afternoon. He listened to both sides patiently and asked everyone to speak without interruption. Finally, the two families reached an agreement: they would share the fruits equally every season. The two fathers apologised, forgave each other and shook hands in front of the whole quarter.

In my opinion, traditional ways of settling disputes are often better than going to court. A palaver is free, quick and friendly, and it restores peace in the community, whereas a court case is expensive and creates more hatred. Dialogue is always stronger than anger.
\end{textebox}

\textbf{Points de vigilance :}
\begin{itemize}
\item Employer le simple past pour le récit (\eng{took}, \eng{started}, \eng{decided}, \eng{reached}) et le présent pour l'opinion générale (\eng{is}, \eng{restores}).
\item Insérer les connecteurs imposés : \eng{first}, \eng{then}, \eng{after that}, \eng{finally}.
\item Ne pas écrire \eng{*more better} ni \eng{*more quickly than to go} : comparatif correct \eng{better than going to court}.
\item \eng{whereas} = « tandis que » : connecteur de contraste très apprécié au Bac.
\end{itemize}
\end{corrige}

\newpage

\coursec{Fiche bilan}

\begin{popfigure}
\begin{tikzpicture}[mindmap, grow cyclic, every node/.style={concept, font=\small}, concept color=popblue, text=white, level 1/.append style={level distance=3.6cm, sibling angle=60}, level 2/.append style={level distance=2.2cm}]
\node[concept, font=\small\bfseries, align=center]{Reading:\\ Settling\\ community\\ disputes}
  child[concept color=poporange] { node[concept, font=\scriptsize, align=center]{Key vocabulary\\ dispute, mediator,\\ reconcile, verdict} }
  child[concept color=popgreen] { node[concept, font=\scriptsize, align=center]{Skimming\\ read fast for\\ the main idea} }
  child[concept color=poppurple] { node[concept, font=\scriptsize, align=center]{Scanning\\ find dates, names,\\ specific facts} }
  child[concept color=poppink] { node[concept, font=\scriptsize, align=center]{Guessing\\ meaning from\\ context clues} }
  child[concept color=popgold] { node[concept, font=\scriptsize, align=center]{Comprehension\\ questions\\ true/false, wh-} }
  child[concept color=popblue!70] { node[concept, font=\scriptsize, align=center]{Production\\ summary +\\ opinion} };
\end{tikzpicture}
\legende{Carte mentale de la leçon : lire un texte sur le règlement des conflits communautaires.}
\end{popfigure}

\begin{bilan}
\textbf{1. Lexique clé du thème (peace and conflict)}
\begin{itemize}
  \item \eng{a dispute} (n.) : un différend, un conflit ; \eng{to settle a dispute} : régler un différend.
  \item \eng{a mediator} (n.) : un médiateur ; \eng{to mediate} : servir de médiateur.
  \item \eng{to reconcile} (v.) : réconcilier ; \eng{reconciliation} (n.) : la réconciliation.
  \item \eng{a verdict / a ruling} (n.) : un verdict / une décision ; \eng{to rule} : statuer.
  \item \eng{a plaintiff} (n.) : un plaignant ; \eng{a defendant} (n.) : un défendeur ; \eng{a witness} (n.) : un témoin.
  \item \eng{compensation} (n.) : une compensation, des dédommagements ; \eng{to compensate} : dédommager.
  \item \eng{a council of elders} : un conseil d'anciens ; \eng{a traditional ruler / chief} : un chef traditionnel.
  \item \eng{to apologise} : présenter des excuses ; \eng{forgiveness} (n.) : le pardon ; \eng{to forgive -- forgave -- forgiven} : pardonner.
  \item \eng{peaceful} (adj.) : pacifique ; \eng{hostile} (adj.) : hostile ; \eng{fair / unfair} : juste / injuste.
\end{itemize}

\textbf{2. Points de langue à connaître par cœur}
\begin{center}
\begin{tabularx}{\linewidth}{|X|X|X|}
\hline
\rowcolor{popblueL} Règle & Exemple anglais & Traduction \\
\hline
Prétérit simple pour raconter des faits passés & \eng{The elders settled the dispute last year.} & Les anciens ont réglé le conflit l'année dernière. \\
\hline
\eng{to} + base verbale pour exprimer le but & \eng{They met to discuss the land problem.} & Ils se sont réunis pour discuter du problème de terre. \\
\hline
Discours rapporté : recul des temps & \eng{He said (that) they had agreed.} & Il a dit qu'ils s'étaient mis d'accord. \\
\hline
Passif pour insister sur l'action & \eng{The verdict was announced by the chief.} & Le verdict a été annoncé par le chef. \\
\hline
\eng{should} pour donner un avis & \eng{Both parties should forgive each other.} & Les deux parties devraient se pardonner. \\
\hline
\end{tabularx}
\end{center}

\textbf{3. Méthodes express de lecture}
\begin{enumerate}
  \item \cle{Skimming} : lire le titre, la première et la dernière phrase de chaque paragraphe pour trouver l'idée générale (qui ? où ? quel problème ? quelle solution ?).
  \item \cle{Scanning} : repérer vite un nom propre, une date, un chiffre ou un mot-clé demandé par la question, sans tout relire.
  \item \cle{Guessing} : deviner le sens d'un mot inconnu grâce au contexte (mots voisins, famille de mots, logique de la phrase) avant de chercher dans le glossaire.
  \item Répondre aux questions en citant le texte : \eng{According to the text, ...} / \eng{Line 3 shows that ...}.
  \item Résumer en 3 phrases : problème + médiation + résultat ; puis donner son opinion avec \eng{In my opinion...}, \eng{I think that...}, \eng{I agree / disagree because...}.
\end{enumerate}
\end{bilan}

\begin{aretenir}[Checklist avant le Bac]
\begin{itemize}
  \item $\square$ Je connais le lexique du conflit et de la paix : \eng{dispute, mediator, reconcile, verdict, compensation, witness}.
  \item $\square$ Je sais faire du \emph{skimming} : trouver l'idée générale en moins de deux minutes.
  \item $\square$ Je sais faire du \emph{scanning} : repérer rapidement un nom, une date ou un chiffre.
  \item $\square$ Je sais deviner le sens d'un mot inconnu grâce au contexte.
  \item $\square$ Je réponds aux questions en appuyant ma réponse sur le texte.
  \item $\square$ Je distingue \eng{plaintiff} et \eng{defendant} sans les confondre.
  \item $\square$ Je conjugue correctement les verbes irréguliers du thème : \eng{forgive -- forgave -- forgiven}, \eng{meet -- met -- met}, \eng{swear -- swore -- sworn}.
  \item $\square$ J'utilise le prétérit simple pour raconter les faits et le passif pour insister sur l'action.
  \item $\square$ Je sais exprimer le but avec \eng{to} + base verbale.
  \item $\square$ Je sais résumer un texte en trois phrases : problème, médiation, résultat.
  \item $\square$ Je sais donner mon opinion en anglais : \eng{In my opinion...}, \eng{I believe that...}.
  \item $\square$ Je fais attention aux faux amis : \eng{to settle} (régler) ne veut pas dire « s'installer » dans ce contexte.
\end{itemize}
\end{aretenir}

\findecours

\end{document}
